% L’économie moderne au Cameroun — Géographie Tle Tech — Cours signé OnBuch+
% Programme officiel MINESEC. Compiler avec : tectonic N03-economie-moderne-cameroun.tex
\newcommand{\DOCMATIERE}{Géographie}
\newcommand{\DOCNIVEAU}{Tle Tech}
\newcommand{\DOCMODULE}{Géographie – classe de Terminale technique}
\newcommand{\DOCLECON}{N03}
\newcommand{\DOCTITRE}{L’économie moderne au Cameroun}
% OnBuch+ — préambule « cours signés » ENRICHI (compilé avec tectonic / XeLaTeX)
% Système visuel « Pop » d'OnBuch+ : fond crème (jamais blanc), bordures encre
% épaisses, ombres « dures » décalées, titres Archivo Black en capitales.
% Version MATHÉMATIQUES : graphes (pgfplots/tikz), boîtes pédagogiques
% (définition, propriété, méthode, attention, le savais-tu, exercice résolu,
% exercice, TP, bilan), page de garde et signature OnBuch+ ancrée en bas.
%
% Macros attendues dans main.tex AVANT \input{preamble} :
%   \DOCMATIERE  \DOCNIVEAU  \DOCMODULE  \DOCLECON (numéro)  \DOCTITRE
\documentclass[11pt]{article}

\usepackage{fontspec}
\usepackage[french]{babel}
\usepackage[a4paper,margin=2cm,top=3.4cm,bottom=2.8cm,headheight=1.4cm,footskip=1.3cm]{geometry}
\usepackage{amsmath,amssymb,amsfonts}
\usepackage{mathtools} % \xmapsto, \coloneqq... souvent utilisés par les modèles
\usepackage{cancel} % simplifications barrées (bilans de forces)
% \abs{z} (module, valeur absolue) et \norm{u} (norme), utilisés par les modèles
\providecommand{\abs}{}\let\abs\relax\DeclarePairedDelimiter{\abs}{\lvert}{\rvert}
\providecommand{\norm}{}\let\norm\relax\DeclarePairedDelimiter{\norm}{\lVert}{\rVert}
\usepackage[table]{xcolor} % [table] : fournit aussi \rowcolors (alternance de lignes)
\usepackage{pagecolor}
\usepackage{tikz}
\usetikzlibrary{arrows.meta,positioning,calc,shapes.geometric,shapes.misc,shapes.arrows,decorations.pathmorphing,decorations.pathreplacing,decorations.markings,patterns,shadows,mindmap,trees,fit,backgrounds,matrix,intersections,angles,quotes}
\usepackage{pgfplots}
\usepackage[version=4]{mhchem} % équations nucléaires \ce{^{235}_{92}U}
\usepackage[european,siunitx]{circuitikz} % schémas électriques
\pgfplotsset{compat=1.18}
\usepgfplotslibrary{fillbetween,groupplots} % aires entre courbes (calcul intégral)
\usepackage{enumitem}
\usepackage{fancyhdr}
% [most] charge toutes les bibliothèques tcolorbox (listings, minted, raster,
% documentation...) et gonfle inutilement la mémoire TeX sur un document long
% (~300 boîtes) : on ne charge que ce qui est réellement utilisé.
\usepackage[skins,breakable]{tcolorbox}
\usepackage{graphicx}
\usepackage{colortbl}
\usepackage{booktabs}
\usepackage{tabularx}
\usepackage{multirow}
\usepackage{diagbox} % case d'en-tête diagonale des tableaux à double entrée
% Colonnes extensibles centrée / gauche / droite (C, L, R), souvent utilisées par les modèles
\newcolumntype{C}{>{\centering\arraybackslash}X}
\newcolumntype{L}{>{\raggedright\arraybackslash}X}
\newcolumntype{R}{>{\raggedleft\arraybackslash}X}
\usepackage{array}
\usepackage{multicol}
\usepackage{siunitx}
% parse-numbers=false : accepte aussi \SI{2,5 \times 10^{-3}}{...}
\sisetup{locale=FR,output-decimal-marker={,},group-minimum-digits=4,parse-numbers=false}
\DeclareSIUnit\ohm{\ensuremath{\symup{\Omega}}} % U+2126 absent de la police
\DeclareSIUnit\division{div} % réglages d'oscilloscope (ms/div, V/div)
\DeclareSIUnit\tr{tr} % tours (vitesse de rotation, ex. \SI{1500}{\tr\per\minute})
\DeclareSIUnit\molar{M} % concentration molaire (ex. \SI{3}{\molar}, \SI{1}{\micro\molar})
\DeclareSIUnit\an{an} % année (le modèle confond parfois avec \year, un compteur TeX) — ex. \SI{5}{\kilo\meter\per\an}
\DeclareSIUnit\FCFA{FCFA} % franc CFA (ex. \SI{500}{\FCFA\per\kilo\gram})
\DeclareSIUnit\degreeBrix{\textdegree Brix} % teneur en sucre d'un sirop/jus (réfractométrie)
\DeclareSIUnit\month{mois} % le modèle utilise parfois \month, un compteur TeX, comme unité de durée
\DeclareSIUnit\microliter{\micro\liter} % le modèle écrit parfois \microliter en un seul token
\DeclareSIUnit\million{millions} % dénombrement (ex. \SI{15}{\million\per\mL} spermatozoïdes)
\usepackage{qrcode}
% \degree : commande fréquemment utilisée par les modèles pour un angle ou une
% température en dehors de \SI{}{} (non fournie par les paquets chargés).
\providecommand{\degree}{^{\circ}}
% \qed (fin de démonstration), utilisé par les modèles sans amsthm
\providecommand{\qed}{\hfill$\square$}
% \barre : double barre des valeurs interdites dans un tableau de variations (tkz-tab)
\providecommand{\barre}{\ensuremath{\|}}
% environnement proof (amsthm non chargé) utilisé par les modèles
\ifdefined\proof\else\newenvironment{proof}[1][Démonstration]{\par\noindent\textit{#1.}\ }{\qed\par}\fi
\usepackage{lastpage}
\usepackage{hyperref}
\hypersetup{hidelinks}

% ============================================================
% Palette « Pop » OnBuch+ — mêmes hex que la classe P (plus_ui.dart)
% ============================================================
\definecolor{poporange}{HTML}{F59321}
\definecolor{popdark}{HTML}{E07A0C}
\definecolor{popcream}{HTML}{FFF6E8}
\definecolor{popink}{HTML}{1C1714}
\definecolor{poppurple}{HTML}{7A5AE0}
\definecolor{popgreen}{HTML}{2E9E6A}
\definecolor{poppink}{HTML}{E0457B}
\definecolor{popblue}{HTML}{2F7DE1}
\definecolor{popgold}{HTML}{C98A12}
\definecolor{popmuted}{HTML}{8A7F76}
\definecolor{popline}{HTML}{EADFD2}
\definecolor{popgray}{HTML}{8A7F76}
\colorlet{popgrey}{popgray}
% Alias tolérants (le modèle nomme parfois les couleurs différemment)
\colorlet{popwhite}{white}
\colorlet{popblack}{popink}
\colorlet{popred}{poppink}
\colorlet{popyellow}{popgold}
% Teintes claires (fonds de tableaux/encadrés) — le modèle nomme parfois
% les couleurs avec un suffixe L (ex. popblueL) sans les avoir définies.
\colorlet{poporangeL}{poporange!20}
\colorlet{popdarkL}{popdark!20}
\colorlet{poppurpleL}{poppurple!20}
\colorlet{popgreenL}{popgreen!20}
\colorlet{poppinkL}{poppink!20}
\colorlet{popblueL}{popblue!20}
\colorlet{popgoldL}{popgold!20}
\colorlet{popredL}{popred!20}
\colorlet{popgrayL}{popgray!20}
% Teintes claires pour fonds de tableaux / zones
\colorlet{popblueL}{popblue!10!white}
\colorlet{poporangeL}{poporange!14!white}
\colorlet{popgreenL}{popgreen!12!white}
\colorlet{poppurpleL}{poppurple!12!white}
\colorlet{poppinkL}{poppink!12!white}
\colorlet{popgoldL}{popgold!14!white}

\pagecolor{popcream}

% ============================================================
% Typographie
% ============================================================
\defaultfontfeatures{Ligatures=TeX}
\setmainfont{PlusJakartaSans-Medium.ttf}[Path=../../../fonts/,BoldFont=PlusJakartaSans-Bold.ttf]
% Maths en sans-serif (Fira Math) avec chiffres et lettres droites de Plus
% Jakarta, pour que les formules se fondent dans le texte.
\usepackage[math-style=ISO,bold-style=ISO]{unicode-math}
\setmathfont{FiraMath-Regular.otf}[Path=../../../fonts/]
\setmathfont{PlusJakartaSans-Medium.ttf}[Path=../../../fonts/,range={up/num,up/{Latin,latin},\mathup}]
\usepackage{icomma} % virgule décimale sans espace en mode math
% Caractères mathématiques Unicode absents des polices de texte (Plus Jakarta,
% Archivo Black) : rendus par leur équivalent mathématique, en texte comme en
% formule — sinon ils s'affichent en carré vide (ex. « Arithmétique dans ℤ »).
\usepackage{newunicodechar}
\newunicodechar{ℝ}{\ensuremath{\mathbb{R}}}
\newunicodechar{ℤ}{\ensuremath{\mathbb{Z}}}
\newunicodechar{ℂ}{\ensuremath{\mathbb{C}}}
\newunicodechar{ℕ}{\ensuremath{\mathbb{N}}}
\newunicodechar{ℚ}{\ensuremath{\mathbb{Q}}}
\newunicodechar{𝔻}{\ensuremath{\mathbb{D}}}
\newunicodechar{α}{\ensuremath{\alpha}}
\newunicodechar{β}{\ensuremath{\beta}}
\newunicodechar{γ}{\ensuremath{\gamma}}
\newunicodechar{δ}{\ensuremath{\delta}}
\newunicodechar{ε}{\ensuremath{\varepsilon}}
\newunicodechar{θ}{\ensuremath{\theta}}
\newunicodechar{λ}{\ensuremath{\lambda}}
\newunicodechar{μ}{\ensuremath{\mu}}
\newunicodechar{σ}{\ensuremath{\sigma}}
\newunicodechar{τ}{\ensuremath{\tau}}
\newunicodechar{φ}{\ensuremath{\varphi}}
\newunicodechar{ψ}{\ensuremath{\psi}}
\newunicodechar{ω}{\ensuremath{\omega}}
\newunicodechar{Σ}{\ensuremath{\Sigma}}
% majuscules produites par \MakeUppercase dans les titres (α-aminés -> Α)
\newunicodechar{Α}{\ensuremath{\alpha}}
\newunicodechar{Β}{\ensuremath{\beta}}
\newunicodechar{Γ}{\ensuremath{\Gamma}}
\newunicodechar{Δ}{\ensuremath{\Delta}}
\newunicodechar{Ε}{\ensuremath{\varepsilon}}
\newunicodechar{Θ}{\ensuremath{\Theta}}
\newunicodechar{Λ}{\ensuremath{\Lambda}}
\newunicodechar{Μ}{\ensuremath{\mu}}
\newunicodechar{Τ}{\ensuremath{\tau}}
\newunicodechar{Φ}{\ensuremath{\Phi}}
\newunicodechar{Ψ}{\ensuremath{\Psi}}
\newunicodechar{Ω}{\ensuremath{\Omega}}
\newunicodechar{↦}{\ensuremath{\mapsto}}
\newunicodechar{∈}{\ensuremath{\in}}
\newunicodechar{∉}{\ensuremath{\notin}}
\newunicodechar{∧}{\ensuremath{\wedge}}
\newunicodechar{⇔}{\ensuremath{\Leftrightarrow}}
\newunicodechar{⟺}{\ensuremath{\Longleftrightarrow}}
\newunicodechar{⇒}{\ensuremath{\Rightarrow}}
\newunicodechar{⟹}{\ensuremath{\Longrightarrow}}
\newunicodechar{∘}{\ensuremath{\circ}}
\newunicodechar{∪}{\ensuremath{\cup}}
\newunicodechar{∩}{\ensuremath{\cap}}
\newunicodechar{≡}{\ensuremath{\equiv}}
\newunicodechar{⊂}{\ensuremath{\subset}}
\newunicodechar{⊥}{\ensuremath{\perp}}
\newunicodechar{∃}{\ensuremath{\exists}}
\newunicodechar{∀}{\ensuremath{\forall}}
\newunicodechar{∛}{\ensuremath{\sqrt[3]{\,}}}
\newunicodechar{∜}{\ensuremath{\sqrt[4]{\,}}}
\newunicodechar{⃗}{\ensuremath{{}^{\to}}}
\newunicodechar{ⁿ}{\ifmmode^{n}\else\textsuperscript{\ensuremath{n}}\fi}
\newunicodechar{ˣ}{\ifmmode^{x}\else\textsuperscript{\ensuremath{x}}\fi}
\newunicodechar{ᵇ}{\ifmmode^{b}\else\textsuperscript{\ensuremath{b}}\fi}
\newunicodechar{ʳ}{\ifmmode^{r}\else\textsuperscript{\ensuremath{r}}\fi}
\newunicodechar{ᵘ}{\ifmmode^{u}\else\textsuperscript{\ensuremath{u}}\fi}
\newunicodechar{ʸ}{\ifmmode^{y}\else\textsuperscript{\ensuremath{y}}\fi}
\newunicodechar{ᵃ}{\ifmmode^{a}\else\textsuperscript{\ensuremath{a}}\fi}
\newunicodechar{ᵉ}{\ifmmode^{e}\else\textsuperscript{\ensuremath{e}}\fi}
\newunicodechar{ᵏ}{\ifmmode^{k}\else\textsuperscript{\ensuremath{k}}\fi}
\newunicodechar{ᵖ}{\ifmmode^{p}\else\textsuperscript{\ensuremath{p}}\fi}
\newunicodechar{ᴺ}{\ifmmode^{N}\else\textsuperscript{\ensuremath{N}}\fi}
\newunicodechar{ᶜ}{\ifmmode^{c}\else\textsuperscript{\ensuremath{c}}\fi}
\newunicodechar{ʰ}{\ifmmode^{h}\else\textsuperscript{\ensuremath{h}}\fi}
\newunicodechar{ᵠ}{\ifmmode^{q}\else\textsuperscript{\ensuremath{q}}\fi}
\newunicodechar{ₙ}{\ifmmode_{n}\else\textsubscript{\ensuremath{n}}\fi}
\newunicodechar{ₐ}{\ifmmode_{a}\else\textsubscript{\ensuremath{a}}\fi}
\newunicodechar{ᵢ}{\ifmmode_{i}\else\textsubscript{\ensuremath{i}}\fi}
\newunicodechar{ᵣ}{\ifmmode_{r}\else\textsubscript{\ensuremath{r}}\fi}
\newunicodechar{ₖ}{\ifmmode_{k}\else\textsubscript{\ensuremath{k}}\fi}
\newunicodechar{ᵦ}{\ifmmode_{\beta}\else\textsubscript{\ensuremath{\beta}}\fi}
\newunicodechar{ₓ}{\ifmmode_{x}\else\textsubscript{\ensuremath{x}}\fi}
\newfontfamily\popdisplay{ArchivoBlack-Regular.ttf}[Path=../../../fonts/]
\newfontfamily\poplogo{SpaceGrotesk-Bold.ttf}[Path=../../../fonts/]
\newfontfamily\popsemi{PlusJakartaSans-SemiBold.ttf}[Path=../../../fonts/]
\newfontfamily\popxbold{PlusJakartaSans-ExtraBold.ttf}[Path=../../../fonts/]
\newfontfamily\popxboldls{PlusJakartaSans-ExtraBold.ttf}[Path=../../../fonts/,LetterSpace=3.0]

\setlist[enumerate]{leftmargin=1.7em,itemsep=5pt,topsep=4pt}
\setlist[itemize]{leftmargin=1.7em,itemsep=4pt,topsep=4pt,label={\color{poporange}\textbullet}}
\setlength{\parindent}{0pt}
\setlength{\parskip}{5pt}
\renewcommand{\arraystretch}{1.25}

% pgfplots : style Pop par défaut
\pgfplotsset{
  every axis/.append style={axis line style={popink,line width=1pt},tick style={popink},
    grid style={popline,line width=0.5pt},label style={font=\small},tick label style={font=\footnotesize}},
  popcurve/.style={poporange,line width=1.8pt,no markers,smooth},
}
% Même style disponible dans un \draw TikZ simple (hors pgfplots)
\tikzset{popcurve/.style={poporange,line width=1.8pt}}
\tikzset{popgrid/.style={popline,very thin}}
\colorlet{popcurve}{poporange} % le nom du style est parfois utilisé comme couleur

% ============================================================
% Pill / squiggle
% ============================================================
\newcommand{\poppill}[3]{%
  \tikz[baseline=-0.55ex]\node[draw=popink,line width=1.2pt,rounded corners=2.6pt,%
    fill=#2,inner xsep=6.5pt,inner ysep=3pt]{%
    \popxboldls\fontsize{7}{7}\selectfont\color{#3}\MakeUppercase{#1}};%
}
\newcommand{\popsquiggle}[1]{%
  \tikz[baseline=-0.4ex]{
    \draw[#1,line width=2.6pt,line cap=round]
      plot [smooth] coordinates {(0,0.09) (0.34,0.01) (0.68,0.21) (1.02,0.01) (1.36,0.10)};
  }%
}
% Mot-clé surligné (marqueur orange sous le texte)
\newcommand{\cle}[1]{{\popxbold\color{popdark}#1}}

% ============================================================
% En-tête / pied
% ============================================================
\pagestyle{fancy}
\fancyhf{}
\renewcommand{\headrulewidth}{0pt}
\renewcommand{\footrulewidth}{0pt}
\lhead{\raisebox{-0.15ex}{\poplogo\fontsize{13}{13}\selectfont\color{popink}OnBuch\color{poporange}+}}
\rhead{\poppill{\DOCMATIERE\ \textbullet\ \DOCNIVEAU\ \textbullet\ Leçon \DOCLECON}{white}{popink}}
\lfoot{\popxbold\fontsize{7.6}{7.6}\selectfont\color{popmuted}\MakeUppercase{Cours signé OnBuch+}\ \textbullet\ {\popsemi\DOCTITRE}}
\rfoot{\tikz[baseline=-0.6ex]\node[rounded corners=8.5pt,fill=popink,minimum height=17pt,minimum width=17pt,inner xsep=5pt,inner sep=0]{\popxbold\fontsize{8}{8}\selectfont\color{popcream}\thepage\,/\,\pageref*{LastPage}};}

% ============================================================
% Titres
% ============================================================
\newcommand{\chaptitre}[2]{%
  \begin{tcolorbox}[enhanced,colback=poporange,colframe=popink,boxrule=2.6pt,arc=11pt,
      drop shadow=popink,left=15pt,right=15pt,top=12pt,bottom=11pt,before skip=3pt,after skip=18pt]
    \poppill{#2}{popink}{popcream}\par\vspace{9pt}
    {\raggedright\hyphenpenalty=10000\exhyphenpenalty=10000\popdisplay\fontsize{22}{24}\selectfont\color{popcream}\MakeUppercase{#1}\par}
  \end{tcolorbox}
}
% Section numérotée : pastille orange avec le numéro + titre Archivo
\newcounter{popsec}
\newcommand{\coursec}[1]{%
  \stepcounter{popsec}\setcounter{popsub}{0}%
  \par\vspace{16pt}\noindent
  \begin{minipage}{\linewidth}
  \tikz[baseline=-0.7ex]\node[draw=popink,line width=1.6pt,fill=poporange,rounded corners=4pt,minimum size=22pt,inner sep=0]{\popdisplay\fontsize{12}{12}\selectfont\color{popcream}\thepopsec};%
  \hspace{8pt}{\popdisplay\fontsize{15}{17}\selectfont\color{popink}\MakeUppercase{#1}}\par
  \vspace{4pt}\noindent{\color{poporange}\rule{36pt}{3.6pt}}
  \end{minipage}\par\nopagebreak\vspace{8pt}
}
\newcounter{popsub}
\newcommand{\courssub}[1]{%
  \stepcounter{popsub}%
  \par\vspace{9pt}\noindent{\popxbold\fontsize{11.5}{13}\selectfont\color{popink}{\color{poporange}\thepopsec.\thepopsub}\hspace{6pt}#1}\par\nopagebreak\vspace{4pt}
}

% ============================================================
% Boîtes pédagogiques — même recette : fond blanc, bordure épaisse colorée,
% ombre dure encre, badge plein. Titre optionnel [#1].
% ============================================================
\tcbset{popbase/.style={breakable,enhanced,colback=white,boxrule=2.3pt,arc=9pt,
  shadow={3pt}{-3pt}{0pt}{popink},left=14pt,right=14pt,top=11pt,bottom=11pt,before skip=11pt,after skip=15pt,
  fontupper=\popsemi\fontsize{10.6}{14.5}\selectfont\color{popink},
  fontlower=\popsemi\fontsize{10.6}{14.5}\selectfont\color{popink}}}
% Coche dessinée (le glyphe ✓ n'existe pas dans les polices du document)
\newcommand{\popcheck}[1]{\tikz[baseline=-0.5ex]\draw[#1,line width=1.6pt,line cap=round,line join=round](-0.13,0.0)--(-0.04,-0.09)--(0.13,0.1);}
\providecommand{\coche}{\popcheck{popgreen}}
\providecommand{\resultat}[1]{\par\smallskip\noindent\fcolorbox{popgreen}{popgreenL}{\parbox{\dimexpr\linewidth-2\fboxsep-2\fboxrule\relax}{\textbf{Résultat.} #1}}\par\smallskip}
\newcommand{\popboxhead}[3]{\poppill{#1}{#2}{white}\ifx\relax#3\relax\else\hspace{6pt}{\popxbold\fontsize{10.6}{12}\selectfont\color{popink}#3}\fi\par\vspace{7pt}}

\newtcolorbox{aretenir}[1][]{popbase,colframe=popgreen,before upper={\popboxhead{À retenir}{popgreen}{#1}}}
\newtcolorbox{exemplebox}[1][]{popbase,colframe=poporange,before upper={\popboxhead{Exemple}{poporange}{#1}}}
\newtcolorbox{definition}[1][]{popbase,colframe=popblue,before upper={\popboxhead{Définition}{popblue}{#1}}}
\newtcolorbox{propriete}[1][]{popbase,colframe=poppurple,before upper={\popboxhead{Propriété}{poppurple}{#1}}}
\newtcolorbox{methode}[1][]{popbase,colframe=popgold,before upper={\popboxhead{Méthode}{popgold}{#1}}}
\newtcolorbox{attention}[1][]{popbase,colframe=poppink,before upper={\popboxhead{Attention}{poppink}{#1}}}
\newtcolorbox{experience}[1][]{popbase,colframe=popblue,colback=popblueL,before upper={\popboxhead{Expérience}{popblue}{#1}}}
% « Le savais-tu ? » : carte violette pleine (contexte, culture, Cameroun)
\newtcolorbox{savaistu}[1][]{popbase,colback=poppurple,colframe=popink,
  fontupper=\popsemi\fontsize{10.4}{14.2}\selectfont\color{white},
  before upper={\poppill{Le savais-tu\,?}{white}{poppurple}\ifx\relax#1\relax\else\hspace{6pt}{\popxbold\color{white}#1}\fi\par\vspace{7pt}}}
% Exercice résolu : énoncé en haut, \tcblower puis la solution sur fond crème
\newcounter{popexo}\newcounter{popexr}
\newtcolorbox{exoresolu}[1][]{popbase,colframe=poporange,colbacklower=popcream,
  segmentation style={popink,line width=1.2pt,dashed},
  before upper={\stepcounter{popexr}\popboxhead{Exercice résolu \thepopexr}{poporange}{#1}},
  before lower={\poppill{Solution}{popink}{popcream}\par\vspace{6pt}}}
% Exercice (non résolu en place) — la correction est dans la partie Corrigés
\newtcolorbox{exercice}[1][]{popbase,colframe=popink,
  before upper={\stepcounter{popexo}\popboxhead{Exercice \thepopexo}{popink}{#1}}}
% Corrigé
\newtcolorbox{corrige}[1][]{popbase,colframe=popgreen,colback=popgreenL,
  before upper={\popboxhead{Corrigé}{popgreen}{#1}}}
% Objectifs / prérequis / bilan
\newtcolorbox{objectifs}{popbase,colframe=popink,colback=poporangeL,
  before upper={\popboxhead{Objectifs de la leçon}{popink}{}}}
\newtcolorbox{prerequis}{popbase,colframe=popink,colback=popblueL,
  before upper={\popboxhead{Prérequis}{popblue}{}}}
\newtcolorbox{bilan}{popbase,colframe=popink,colback=white,boxrule=2.8pt,
  before upper={\popboxhead{Fiche bilan}{poporange}{L'essentiel à connaître pour l'examen}}}
% Figure encadrée avec légende
\newtcolorbox{popfigure}[1][]{enhanced,breakable=false,colback=white,colframe=popline,boxrule=1.4pt,arc=8pt,
  left=10pt,right=10pt,top=10pt,bottom=8pt,before skip=10pt,after skip=12pt,halign=center,
  lower separated=false,halign lower=center,
  fontlower=\popsemi\fontsize{9}{11.5}\selectfont\color{popmuted},
  #1}
\newcommand{\legende}[1]{\par\vspace{6pt}{\popsemi\fontsize{9}{11.5}\selectfont\color{popmuted}#1}}

% ============================================================
% Signature OnBuch+
% ============================================================
\newcommand{\poplogoinline}[1]{{\poplogo\fontsize{#1}{#1}\selectfont\color{popink}OnBuch\color{poporange}+}}
\newcommand{\signatureonbuch}{%
  \begin{tcolorbox}[enhanced,colback=popink,colframe=popink,boxrule=2.2pt,arc=10pt,
      drop shadow=poporange,left=14pt,right=14pt,top=10pt,bottom=10pt,before skip=0pt,after skip=0pt]
    \tikz[baseline=-0.7ex]\node[circle,fill=popgreen,draw=popcream,line width=1.3pt,minimum size=22pt,inner sep=0]{\popcheck{white}};%
    \hspace{9pt}\begin{minipage}[c]{0.62\linewidth}
      {\popxbold\fontsize{12}{13}\selectfont\color{popcream}Signé\ \textbullet\ {\poplogo OnBuch\color{poporange}+}}\par\vspace{2pt}
      {\raggedright\popsemi\fontsize{8}{10.5}\selectfont\color{popline}\MakeUppercase{Équipe pédagogique OnBuch+}\\\MakeUppercase{Conforme au programme officiel MINESEC}\par}
    \end{minipage}\hfill
    {\poplogo\fontsize{22}{22}\selectfont\color{popcream}OnBuch\color{poporange}+}
  \end{tcolorbox}%
}

% ============================================================
% Page de garde
% #1 titre · #2 matière · #3 niveau · #4 module · #5 n° leçon · #6 durée ·
% #7 description / objectifs (texte ou itemize)
% ============================================================
\newcommand{\pagegarde}[7]{%
  \thispagestyle{empty}
  \newgeometry{a4paper,margin=1.7cm,top=1.6cm,bottom=1.4cm}%
  \begin{tikzpicture}[remember picture,overlay]
    \fill[poporange!18] ([xshift=-3.2cm,yshift=-3.6cm]current page.north east) circle (4.6cm);
    \fill[poppurple!14] ([xshift=2.4cm,yshift=7.6cm]current page.south west) circle (3.8cm);
  \end{tikzpicture}%
  {\poplogo\fontsize{30}{30}\selectfont\color{popink}OnBuch\color{poporange}+}\hfill
  \raisebox{0.6ex}{\poppill{Cours signé \textbullet\ Premium}{popink}{popcream}}\par
  \vspace{1pt}\popsquiggle{poppurple}\par
  \vspace{8pt}
  \begin{tcolorbox}[enhanced,colback=poporange,colframe=popink,boxrule=2.8pt,arc=13pt,
      drop shadow=popink,left=18pt,right=18pt,top=12pt,bottom=13pt,after skip=10pt]
    \poppill{#2 \textbullet\ #3}{popink}{popcream}\hspace{5pt}\poppill{#4}{white}{popink}\par\vspace{8pt}
    {\popdisplay\fontsize{14}{14}\selectfont\color{popink}LEÇON #5}\par\vspace{4pt}
    {\raggedright\hyphenpenalty=10000\exhyphenpenalty=10000\popdisplay\fontsize{25}{27}\selectfont\color{popcream}\MakeUppercase{#1}\par}
    \vspace{8pt}
    {\popxbold\fontsize{9.5}{9.5}\selectfont\color{popink}\MakeUppercase{Durée conseillée : #6}}
  \end{tcolorbox}
  \begin{tcolorbox}[enhanced,colback=white,colframe=popink,boxrule=2.2pt,arc=10pt,
      drop shadow=popink,left=16pt,right=16pt,top=10pt,bottom=10pt,after skip=8pt]
    \poppill{Dans cette leçon}{poppurple}{white}\par\vspace{5pt}
    \popsemi\fontsize{9.3}{12.6}\selectfont\color{popink}#7
  \end{tcolorbox}
  \noindent\begin{minipage}[t]{0.487\linewidth}
    \begin{tcolorbox}[enhanced,colback=popgreen,colframe=popink,boxrule=2.2pt,arc=10pt,
        drop shadow=popink,left=11pt,right=11pt,top=10pt,bottom=10pt,height=3.1cm,valign=center]
      \tcbox[colback=white,colframe=popink,boxrule=1.3pt,arc=5pt,boxsep=0pt,nobeforeafter,
        left=3pt,right=3pt,top=3pt,bottom=3pt,valign=center]{\qrcode[height=1.9cm]{https://play.google.com/store/apps/details?id=com.luvvix.onbuch&hl=fr}}%
      \hspace{8pt}\begin{minipage}[c]{0.5\linewidth}
        {\popdisplay\fontsize{11}{12.5}\selectfont\color{white}\MakeUppercase{Télécharge OnBuch}}\par\vspace{4pt}
        {\raggedright\popsemi\fontsize{8.4}{11}\selectfont\color{white}Gratuit \textbullet\ Google Play\\Tuteur IA, annales, concours, résultats\par}
      \end{minipage}
    \end{tcolorbox}
  \end{minipage}\hfill
  \begin{minipage}[t]{0.487\linewidth}
    \begin{tcolorbox}[enhanced,colback=poppurple,colframe=popink,boxrule=2.2pt,arc=10pt,
        drop shadow=popink,left=11pt,right=11pt,top=10pt,bottom=10pt,height=3.1cm,valign=center]
      \tikz[baseline=-0.7ex]\node[circle,fill=white,draw=popink,line width=1.3pt,minimum size=1.6cm,inner sep=0]{\popdisplay\fontsize{24}{24}\selectfont\color{poppurple}+};%
      \hspace{8pt}\begin{minipage}[c]{0.6\linewidth}
        {\popdisplay\fontsize{11}{12.5}\selectfont\color{white}\MakeUppercase{Passe à OnBuch+}}\par\vspace{4pt}
        {\raggedright\popsemi\fontsize{8.4}{11}\selectfont\color{white}Tous les cours signés, Tuteur IA illimité, fiches PDF — onglet OnBuch+ de l'app\par}
      \end{minipage}
    \end{tcolorbox}
  \end{minipage}
  \vspace{8pt}
  \popcontenu
  \vfill
  \signatureonbuch
  \restoregeometry
  \newpage
}

% Rangée « Ce cours contient » (page de garde)
\newcommand{\poptile}[3]{% #1 couleur #2 gros texte #3 légende
  \begin{tcolorbox}[enhanced,colback=#1,colframe=popink,boxrule=2pt,arc=9pt,shadow={2.5pt}{-2.5pt}{0pt}{popink},
      left=6pt,right=6pt,top=9pt,bottom=9pt,halign=center,valign=center,height=1.75cm,width=\linewidth,nobeforeafter]
    {\popdisplay\fontsize{12.5}{13}\selectfont\color{white}\MakeUppercase{#2}}\par\vspace{3pt}
    {\centering\popsemi\fontsize{7.8}{9.5}\selectfont\color{white}#3\par}
  \end{tcolorbox}}
\newcommand{\popcontenu}{%
  \par\noindent{\popxboldls\fontsize{7.5}{7.5}\selectfont\color{popmuted}CE COURS SIGNÉ CONTIENT}\par\vspace{7pt}
  \noindent\begin{minipage}[t]{0.235\linewidth}\poptile{popblue}{Cours}{complet, illustré, pas à pas}\end{minipage}\hfill
  \begin{minipage}[t]{0.235\linewidth}\poptile{poporange}{Méthodes}{et exercices résolus}\end{minipage}\hfill
  \begin{minipage}[t]{0.235\linewidth}\poptile{poppink}{Exercices}{type examen + corrigés}\end{minipage}\hfill
  \begin{minipage}[t]{0.235\linewidth}\poptile{popgreen}{Fiche bilan}{l'essentiel à retenir}\end{minipage}\par}

% Sommaire
\newtcolorbox{sommaire}{enhanced,colback=white,colframe=popink,boxrule=2.2pt,arc=10pt,
  drop shadow=popink,left=18pt,right=18pt,top=13pt,bottom=13pt,before skip=6pt,after skip=20pt,
  before upper={\poppill{Sommaire}{poppurple}{white}\par\vspace{9pt}\popsemi\fontsize{10.6}{14.5}\selectfont\color{popink}}}

% Signature de fin de cours — poussée tout en bas de la dernière page
\newcommand{\findecours}{%
  \par\vfill\nopagebreak
  \begin{center}\popsquiggle{poporange}\end{center}\vspace{4pt}
  \signatureonbuch
}
\AtBeginDocument{\def\llbracket{\mathopen{[\mkern-3.5mu[}}\def\rrbracket{\mathclose{]\mkern-3.5mu]}}} % intervalles d'entiers (glyphe ⟦ absent de la police)
% Glyphes absents de Fira Math : redéfinis à partir de glyphes disponibles
\AtBeginDocument{%
  \def\setminus{\mathbin{\backslash}}\let\smallsetminus\setminus
  \def\vdots{\mathord{\vbox{\baselineskip3.2pt\lineskiplimit0pt\kern4pt\hbox{.}\hbox{.}\hbox{.}}}}%
  \def\ddots{\mathinner{\mkern1mu\raise6.5pt\hbox{.}\mkern2mu\raise3.6pt\hbox{.}\mkern2mu\raise0.7pt\hbox{.}\mkern1mu}}%
  \def\triangle{\mathord{\scalebox{0.9}{$\symup{\Delta}$}}}%
  \def\nabla{\mathord{\raisebox{\depth}{\scalebox{1}[-1]{$\symup{\Delta}$}}}}%
  \def\checkmark{\popcheck{popdark}}%
  \def\star{\mathbin{\ast}}\def\bigstar{\mathord{\ast}}%
  \def\diamond{\mathbin{\scalebox{0.6}{$\Diamond$}}}%
  \def\bigcup{\mathop{\mathchoice{\raisebox{-0.3em}{\scalebox{1.7}{$\cup$}}}{\scalebox{1.25}{$\cup$}}{\cup}{\cup}}\displaylimits}%
  \def\bigcap{\mathop{\mathchoice{\raisebox{-0.3em}{\scalebox{1.7}{$\cap$}}}{\scalebox{1.25}{$\cap$}}{\cap}{\cap}}\displaylimits}%
  \def\lesseqqgtr{\mathrel{\lessgtr}}\def\bigoplus{\mathop{\oplus}\displaylimits}%
}
\providecommand{\ssi}{si et seulement si} % abréviation fréquente
\AtBeginDocument{\providecommand{\wideparen}{\overparen}} % arc de cercle

%% FIGFIX-BEGIN v5 — correctifs globaux des figures (docs/onbuch-generation/tools/figfix)
\usepackage{adjustbox}\usepackage{etoolbox}\tcbuselibrary{raster}
% toute figure TikZ/pgfplots plus large que la ligne est réduite à la largeur disponible
\BeforeBeginEnvironment{tikzpicture}{\begin{adjustbox}{max width=\linewidth}}
\AfterEndEnvironment{tikzpicture}{\end{adjustbox}}
% légendes pgfplots lisibles par-dessus les courbes
\pgfplotsset{every axis/.append style={legend style={fill=white,fill opacity=0.92,text opacity=1,draw=black!15,inner sep=3pt}}}
% étiquettes d'arêtes (syntaxe edge["..."]) avec fond blanc
\tikzset{every edge quotes/.append style={fill=white,inner sep=1.2pt}}
%% FIGFIX-END
\begin{document}

\pagegarde{L’économie moderne au Cameroun}{Géographie}{Tle Tech}{Géographie – classe de Terminale technique}{N03}{2 séances (fiche de progression harmonisée)}{Cette leçon analyse les deux piliers de l'économie moderne camerounaise : l'agriculture de plantation à vocation d'exportation et le secteur industriel. Elle met en lumière leur organisation spatiale, leurs acteurs historiques et actuels, leurs enjeux socio-économiques et environnementaux, à travers des études de cas, des croquis et l'exploitation de données chiffrées officielles.
\vspace{4pt}
\begin{multicols}{2}\small
\begin{itemize}
\item À la fin de cette leçon, je sais caractériser l'agriculture de grandes plantations (origines, statuts fonciers, principales spéculations, acteurs).
\item Je sais localiser avec précision les grands bassins de plantation (CDC, HEVECAM, SOCAPALM, PHP) et les pôles industriels majeurs (Douala-Bonabéri, Edéa, Yaoundé, Kribi, Nord) sur une carte muette.
\item Je sais analyser l'organisation spatiale et productive d'une firme agro-industrielle type (noyau central, planteurs villageois, chaîne de valeur, infrastructures sociales).
\item Je sais expliquer les facteurs de localisation de l'industrie camerounaise (énergie, port, matières premières, marché, transport) et en hiérarchiser l'importance.
\item Je sais différencier les types d'industries (agro-alimentaires, chimiques, métallurgiques, matériaux de construction) et identifier leurs implantations régionales.
\item Je sais construire et commenter un croquis de synthèse "Les pôles de l'économie moderne au Cameroun" respectant la sémiologie graphique et une légende hiérarchisée.
\end{itemize}\end{multicols}}

\begin{sommaire}
\begin{multicols}{2}
\begin{enumerate}
\item Cadre général : Dualisme et poids de l'économie moderne
\item L'agriculture de grandes plantations : Genèse, acteurs et statuts fonciers
\item Étude de cas guidée : Organisation spatiale et enjeux d'une firme agro-industrielle (CDC / SOCAPALM / HEVECAM)
\item Le secteur industriel : Typologie, facteurs de localisation et pôles majeurs
\item TP Intégré : Croquis de synthèse « Les pôles de l'économie moderne au Cameroun »
\item Enjeux transversaux : Agro-industrie, emploi, environnement et trajectoire d'émergence
\item Activité d'intégration
\item Méthodes et exercices résolus
\item Exercices
\item Corrigés des exercices
\item Fiche bilan
\end{enumerate}
\end{multicols}
\end{sommaire}

\begin{objectifs}
\begin{itemize}
\item À la fin de cette leçon, je sais caractériser l'agriculture de grandes plantations (origines, statuts fonciers, principales spéculations, acteurs).
\item Je sais localiser avec précision les grands bassins de plantation (CDC, HEVECAM, SOCAPALM, PHP) et les pôles industriels majeurs (Douala-Bonabéri, Edéa, Yaoundé, Kribi, Nord) sur une carte muette.
\item Je sais analyser l'organisation spatiale et productive d'une firme agro-industrielle type (noyau central, planteurs villageois, chaîne de valeur, infrastructures sociales).
\item Je sais expliquer les facteurs de localisation de l'industrie camerounaise (énergie, port, matières premières, marché, transport) et en hiérarchiser l'importance.
\item Je sais différencier les types d'industries (agro-alimentaires, chimiques, métallurgiques, matériaux de construction) et identifier leurs implantations régionales.
\item Je sais construire et commenter un croquis de synthèse "Les pôles de l'économie moderne au Cameroun" respectant la sémiologie graphique et une légende hiérarchisée.
\item Je sais évaluer les impacts socio-économiques (emploi, recettes d'exportation, infrastructures) et environnementaux (déforestation, pollution, conflits fonciers) de l'économie moderne.
\item Je sais rédiger une argumentation structurée sur la place de l'agro-industrie et de la transformation locale dans la trajectoire d'émergence du Cameroun (PND30, Vision 2035).
\end{itemize}
\end{objectifs}

\begin{prerequis}
\begin{itemize}
\item Notions d'espace rural, espace urbain, activités primaires/secondaires (Programme 4ème/3ème).
\item Facteurs de localisation des activités humaines et notion de pôle de développement (Début année Terminale / Géographie 1ère).
\item Repères historiques : colonisation allemande, mandats français/britannique, indépendances (Histoire/Géographie 1ère).
\item Indicateurs économiques de base : PIB, Balance commerciale, Investissements Directs Étrangers (IDE), Taux de transformation locale (Économie/Geographie 1ère).
\item Maîtrise de base de la sémiologie graphique (points, surfaces, traits, couleurs) pour la cartographie.
\end{itemize}
\end{prerequis}

\newpage

\coursec{Cadre général : Dualisme et poids de l'économie moderne}

Dans le rayon frais d'un supermarché de Douala, des bananes « Plantations du Haut-Penja » côtoient des ananas « Afrit », tandis qu'à l'usine SOSUCAM de Nkoteng, la canne à sucre devient sucre raffiné pour le marché national. Parallèlement, les fumées de l'aluminerie ALUCAM à Edéa et les cuves de la raffinerie SONARA à Limbé dessinent le profil industriel du pays. Comment ces grandes exploitations agricoles et ces unités industrielles structurent-elles l'espace national, génèrent-elles de la valeur ajoutée et relèvent-elles les défis de l'emploi et de la transformation locale pour l'émergence à l'horizon 2035 ?

\courssub{Définition et contours de l'économie moderne}

L'économie moderne au Cameroun désigne l'ensemble des activités de production de biens et de services organisées selon des normes \textbf{formelles}, \textbf{capitalistiques} et \textbf{salariées}. Elle se distingue par l'usage intensif de capital (machines, infrastructures, technologies), une productivité du travail élevée, une insertion dans les circuits commerciaux structurés (exportation, grande distribution nationale) et une soumission à la réglementation de l'État (droit du travail, fiscalité, normes environnementales).

\begin{definition}[Économie moderne (secteur formel)]
Ensemble des unités de production (firmes agro-industrielles, industries manufacturières, services modernes) caractérisées par : l'intensité en capital, la forte productivité du travail, le salariat déclaré, la comptabilité normalisée, l'orientation vers le marché solvable (exportation ou grande consommation urbaine) et l'encadrement juridique par l'État.
\end{definition}

Deux piliers la composent :
\begin{itemize}
    \item \textbf{L'agriculture de grandes plantations} (banane, caoutchouc, palmier à huile, thé, ananas, canne à sucre) : exploitations de plusieurs centaines à dizaines de milliers d'hectares, mono-culture, main-d'œuvre permanente et saisonnière nombreuse, fort taux d'exportation.
    \item \textbf{L'industrie} : transformation des ressources locales (agro-alimentaire, bois, aluminium, raffinage, ciment, textile) et industries d'assemblage ou de fabrication (métallurgie, chimie, BTP).
\end{itemize}

\courssub{Poids dans le PIB : des contributions contrastées}

Les données de l'Institut National de la Statistique (INS) et de la Banque des États de l'Afrique Centrale (BEAC) pour les exercices 2022-2023 révèlent une structure tertiarisée de l'économie, où le secteur moderne pèse significativement dans la création de richesse mais reste étroit en volume d'emplois.

\begin{popfigure}
\begin{tikzpicture}
    % Camembert manuel (somme = 100)
    \def\angle{90}
    \def\radius{3}
    \foreach \percent/\color/\label in {
        49/popblue!80/Tertiaire (49\%),
        28/popgreen!70/Secondaire (28\%),
        10/poporange!70/Primaire traditionnel (10\%),
        6/poppurple!70/Primaire moderne (6\%),
        7/poppink!70/Autres (7\%)
    } {
        \pgfmathsetmacro{\newangle}{\angle - \percent/100*360}
        \draw[fill=\color, thick] (0,0) -- (\angle:\radius) arc (\angle:\newangle:\radius) -- cycle;
        \pgfmathsetmacro{\midangle}{(\angle+\newangle)/2}
        % Légende externe sera dans \legende
        \global\let\angle\newangle
    }
    \node[above] at (0,3.8) {\textbf{Structure du PIB par grande catégorie (estimation 2022-2023)}};
\end{tikzpicture}
\legende{Source : INS, Comptes nationaux 2022-2023 ; BEAC, Rapports annuels. Le secteur primaire moderne (plantations, pétrole, bois industriel) représente environ 6\% du PIB mais tire l'essentiel des recettes d'exportation hors pétrole. L'industrie manufacturière stagne autour de 15-17\% du PIB.}
\end{popfigure}

\begin{aretenir}[Poids sectoriels clés (estimations 2022-2023)]
\begin{itemize}
    \item \textbf{Secteur primaire total} : $\sim$ 16\% du PIB (dont \textbf{primaire moderne} $\sim$ 6\% : plantations, pétrole, bois industriel).
    \item \textbf{Secteur secondaire} : $\sim$ 28\% du PIB (Industrie manufacturière $\sim$ 16\%, BTP $\sim$ 9\%, Énergie/Eau $\sim$ 3\%).
    \item \textbf{Secteur tertiaire} : $\sim$ 49\% du PIB (Services marchands, administration, commerce, télécoms, finance).
    \item \textbf{Autres} (impôts nets, etc.) : $\sim$ 7\%.
\end{itemize}
\end{aretenir}

\begin{methode}[Lire et interpréter un graphique de structure du PIB ou des exportations]
\begin{enumerate}
    \item \textbf{Identifier le type de graphique} : diagramme circulaire (parts relatives), histogramme (évolution temporelle), courbes (parts sectorielles dans le temps).
    \item \textbf{Repérer les légendes et unités} : \% du PIB, milliards de FCFA, tonnes, valeur FOB.
    \item \textbf{Isoler les composantes « modernes »} : dans le primaire, distinguer \emph{moderne} (plantations, pétrole, bois scié) et \emph{traditionnel} (vivrier, cacao/café paysan). Dans le secondaire, isoler l'\emph{industrie manufacturière} du BTP.
    \item \textbf{Quantifier la dépendance} : calculer la part des hydrocarbures (pétrole brut + gaz), du bois (grumes + sciages) et de l'agro-industrie (banane, caoutchouc, palme, coton, aluminium) dans les exportations totales.
    \item \textbf{Conclure sur la vulnérabilité} : une part > 60-70\% pour 3-4 produits de base signale une économie de rente peu diversifiée.
\end{enumerate}
\end{methode}

\begin{exemplebox}[Application : Lecture des exportations 2022]
En 2022, les exportations camerounaises s'élèvent à environ \num{3 500} milliards FCFA.
\begin{itemize}
    \item Pétrole brut + Gaz : $\sim$ 45-50\%.
    \item Bois (grumes, sciages, contreplaqué) : $\sim$ 12-15\%.
    \item Agro-industrie (Banane, Caoutchouc, Huile de palme, Coton, Ananas, Café/cacao transformés) : $\sim$ 15-18\%.
    \item Aluminium (ALUCAM) : $\sim$ 5-7\%.
\end{itemize}
\textbf{Constat} : Hors pétrole, l'agro-industrie, le bois et l'aluminium représentent \textbf{environ 35-40\% des recettes d'exportation totales} (soit 60-75\% des exportations non-pétrolières). La diversification reste faible.
\end{exemplebox}

\courssub{Le dualisme structurel : une coexistence inégale}

L'économie camerounaise présente un \textbf{dualisme} marqué, hérité de l'histoire coloniale et entretenu par les choix de développement post-indépendance.

\begin{definition}[Dualisme économique]
Coexistence, au sein d'un même espace national, de deux secteurs aux logiques antagonistes :
\begin{itemize}
    \item \textbf{Secteur moderne (formel)} : capitalistique, productif, exportateur, salarié, fiscalisé, concentré spatialement (littoral, grands axes).
    \item \textbf{Secteur traditionnel / informel} : travail familial, faible capitalisation, productivité faible, autoconsommation ou marché local précaire, emploi massif (≈ 90\% de la population active), peu fiscalisé, diffus spatialement.
\end{itemize}
\end{definition}

\begin{center}
\begin{tabularx}{\linewidth}{|l|X|X|}
\hline
\rowcolor{popblueL} \textbf{Critère} & \textbf{Secteur Moderne (Formel)} & \textbf{Secteur Traditionnel / Informel} \\ \hline
\textbf{Acteurs} & Multinationales (Bolloré, Socfin, Olam), Grandes EP (CDC, SONARA, ALUCAM), PME formelles & Paysans, artisans, commerçants ambulants, micro-entreprises non déclarées \\ \hline
\textbf{Capital / Travail} & Capital intensif (K/L fort) & Travail intensif (L/K fort) \\ \hline
\textbf{Productivité} & Élevée (VA / tête > 5-10 millions FCFA/an) & Faible (VA / tête < 500 000 FCFA/an) \\ \hline
\textbf{Marché} & Exportation + Grande distribution urbaine & Autoconsommation + Marchés locaux informels \\ \hline
\textbf{Emploi} & $\sim$ 10-15\% de la pop. active (salarié déclaré) & $\sim$ 85-90\% de la pop. active (sous-emploi massif) \\ \hline
\textbf{Fiscalité} & Contribuable majeur (IS, TVA, droits douaniers) & Quasi inexistante (niche fiscale) \\ \hline
\textbf{Localisation} & Pôles côtiers (Douala, Limbé, Kribi), Edéa, Yaoundé, axes routiers bitumés & Partout, dominant en zone rurale et quartiers périphériques \\ \hline
\end{tabularx}
\end{center}

\begin{attention}[Erreur fréquente : Confondre « agriculture moderne » et « agriculture de rente paysanne »]
\begin{itemize}
    \item \textbf{Agriculture moderne (plantations)} : Grandes surfaces (\textgreater 100 ha, souvent 1 000-20 000 ha), mono-culture, salariat, capital étranger ou public, forte technicité (irrigation, drainage, clonal), \textbf{exportation quasi totale}. Ex : CDC (caoutchouc/palme), PHP (banane), HEVECAM (hévéa), SOSUCAM (canne).
    \item \textbf{Agriculture de rente paysanne} : Petites exploitations (1-5 ha), polyculture-assolement, travail familial, faibles intrants, commercialisation via filières coopératives/collecteurs, \textbf{marché national + exportation}. Ex : Cacao (Sud, Centre, Est), Café (Ouest, Nord-Ouest), Coton (Nord, Adamaoua).
\end{itemize}
\textbf{Piège} : Dire « le cacao est une plantation » est faux au Cameroun (c'est paysan). Dire « la banane d'exportation est paysanne » est faux (c'est plantation).
\end{attention}

\courssub{Le rôle de l'État : actionnaire historique et régulateur}

L'État camerounais a été le premier entrepreneur industriel et agricole via les \textbf{Entreprises Publiques (EP)} créées dès les années 1960-70 pour maîtriser les « hauteurs de l'économie ».

\begin{itemize}
    \item \textbf{Agro-industrie} : \textbf{CDC} (Cameroun Development Corporation, 1947, reprise État 1960), \textbf{HEVECAM} (1975), \textbf{SOCAPALM} (1968), \textbf{SOSUCAM} (1965), \textbf{SPFS} (Société des Plantations de Mbanga), \textbf{UNEX PALM}.
    \item \textbf{Industrie lourde} : \textbf{ALUCAM} (Aluminium, Edéa, 1957), \textbf{SONARA} (Raffinage, Limbé, 1973), \textbf{CIMENCAM} (Ciments), \textbf{REGIFER/CAMRAIL} (Transport ferroviaire).
\end{itemize}

Depuis les années 1990 (Plans d'Ajustement Structurel), l'État se retire de la gestion directe (\textbf{privatisations} : HEVECAM $\to$ GMG/Socfin, SOSUCAM $\to$ Somdiaa, CIMENCAM $\to$ Lafarge/Holcim, partie CDC $\to$ Pamol) pour endosser le rôle de \textbf{régulateur} et de \textbf{facilitateur} :
\begin{itemize}
    \item \textbf{MINMIDT} (Ministère de l'Industrie, des Mines et du Développement Technologique) : politique industrielle, zones industrielles, promotion du « Made in Cameroon ».
    \item \textbf{MINADER} (Ministère de l'Agriculture et du Développement Rural) : encadrement des filières, recherche (IRAD), vulgarisation.
    \item \textbf{Agence de Promotion des Investissements (API)} : guichet unique, accompagnement investisseurs.
\end{itemize}

\courssub{Ouverture aux IDE et cadres incitatifs}

Pour combler le déficit d'investissement productif, le Cameroun mise sur les \textbf{Investissements Directs Étrangers (IDE)}.

\begin{propriete}[Cadre juridique d'attraction des IDE (Loi n°2013/004 du 18 avril 2013)]
\begin{itemize}
    \item \textbf{Régime commun} : Liberté d'investissement, rapatriement des capitaux et revenus, égalité de traitement investisseurs nationaux/étrangers.
    \item \textbf{Régimes incitatifs} : \emph{Régime d'incitation à l'investissement} (exonérations IS 5-10 ans, droits de douane sur équipements, TVA) conditionné à : montant investi (\textgreater 50 millions FCFA), création d'emplois, formation locale, transfert de technologie, localisation (zones prioritaires).
    \item \textbf{Zones Franches Industrielles (ZFI)} : Kribi (Zone à Statut Particulier), Douala-Bassa, Yaoundé-Nsimalen. Exonérations totales (IS, TVA, droits douaniers) pour entreprises exportant \textgreater 80\% de la production.
    \item \textbf{Partenariats Public-Privé (PPP)} : Loi 2006/012 modifiée. Port de Kribi (Kribi Conteneur Terminal - KCT), Barrage de Nachtigal (NHPC), Autoroute Yaoundé-Douala (projet), Centrales thermiques.
\end{itemize}
\end{propriete}

\begin{savaistu}[Le Cameroun, premier exportateur africain de banane vers l'UE]
Grâce aux accords de \textbf{Cotonou} (2000) puis aux \textbf{Accords de Partenariat Économique (APE)} intérimaires (2009/2014) avec l'Union Européenne, le Cameroun bénéficie d'un accès \textbf{sans droits de douane ni contingents} pour ses bananes. En 2022-2023, il exporte $\sim$ \num{300 000} à \num{350 000} tonnes/an, devançant la Côte d'Ivoire et le Ghana. Cette rente de position (préférence tarifaire) est menacée par la libéralisation progressive (tarif « Most Favoured Nation » à 114 €/tonne) et la concurrence de l'Équateur (premier exportateur mondial, coûts inférieurs).
\end{savaistu}

\courssub{Vulnérabilités structurelles de l'économie moderne}

Malgré son poids dans le PIB et les exportations, l'économie moderne camerounaise souffre de faiblesses structurelles qui freinent l'émergence.

\begin{enumerate}
    \item \textbf{Dépendance aux cours mondiaux (Choc exogène)} : Pétrole, bois, aluminium, banane, caoutchouc, coton sont des \emph{commodities}. La volatilité des prix (ex: cours du caoutchouc divisé par 3 entre 2011 et 2020) impacte directement la balance commerciale, le budget de l'État (recettes pétrolières) et la trésorerie des firmes.
    \item \textbf{Faible taux de transformation locale} :
    \begin{itemize}
        \item Bois : Taux de transformation locale $\sim$ 30-35\% (objectif 100\% grumes interdites à l'export depuis 2017, mais contournements).
        \item Cacao : $\sim$ 30-35\% transformé localement (poudre, beurre, liqueur) vs 65-70\% fèves brutes.
        \item Coton : $\sim$ 30-40\% égrené/filé localement (CICAM, Sodecoton) vs export fibre.
        \item Pétrole : Raffinage local (SONARA) couvre $\sim$ 60-70\% des besoins, mais vétusté + sous-capacité.
    \end{itemize}
    \item \textbf{Balance commerciale structurellement déficitaire (hors pétrole/bois)} : Le pays importe massivement biens d'équipement, produits manufacturés, hydrocarbures raffinés (si SONARA arrêtée), médicaments, céréales (blé, riz). Le solde commercial non-pétrolier est souvent négatif.
    \item \textbf{Enclavement logistique et coûts de transaction} : Douala (port unique saturé), routes dégradées, coûts énergie (coupures), télécoms chers, lourdeurs administratives (classement \emph{Doing Business} faible).
\end{enumerate}

\begin{exoresolu}[Analyse de la vulnérabilité : Le cas du caoutchouc naturel (HEVECAM / CDC)]
\textbf{Énoncé} : En 2011, le cours mondial du caoutchouc naturel (TSR20) atteint 4,80 \$/kg. En 2020, il chute à 1,30 \$/kg. HEVECAM (42 000 ha) et CDC (35 000 ha hévéa) emploient ensemble $\sim$ 25 000 permanents. Analyser les conséquences sur l'économie moderne locale.
\tcblower
\textbf{Solution détaillée} :
\begin{enumerate}
    \item \textbf{Chute du Chiffre d'Affaires (CA)} : Pour un rendement moyen de 1,5 t/ha/an et 77 000 ha totaux $\to$ $\sim$ 115 500 t/an. CA 2011 $\sim$ 115 500 $\times$ 4,80 $\times$ 500 (taux de change) $\approx$ 277 milliards FCFA. CA 2020 $\sim$ 115 500 $\times$ 1,30 $\times$ 550 $\approx$ 83 milliards FCFA. \textbf{Perte $\sim$ 194 milliards FCFA/an}.
    \item \textbf{Trésorerie \& Investissement} : Gel des replantations (cycle 25-30 ans), report maintenance usines, difficultés paiement salaires/fournisseurs.
    \item \textbf{Effet d'entraînement (Multiplicateur local)} : Baisse commande PME sous-traitantes (transport, maintenance, BTP), baisse consommation ménages employés (commerces locaux, marchés).
    \item \textbf{Réponse stratégique} : Diversification (Palmier à huile chez HEVECAM/CDC), certification durable (RSPO) pour primes de marché, recherche variétés à haut rendement/précocité (IRAD/PRCC), transformation locale (usines de gants, chambres à air) pour capter valeur ajoutée.
\end{enumerate}
\textbf{Conclusion} : La mono-spécialisation expose les bassins d'emploi (Sud-Ouest, Sud, Littoral) à une instabilité sociale et économique majeure.
\end{exoresolu}

\courssub{Synthèse cartographique : Les 5 grands pôles de l'économie moderne}

L'espace national est structuré par quelques pôles concentriques où se concentrent capital, infrastructures, main-d'œuvre qualifiée et pouvoir de décision.

\begin{popfigure}
\begin{tikzpicture}[scale=0.75]
    % Contour très simplifié du Cameroun (polygone approximatif)
    \draw[popdark, thick, fill=popcream!30] 
        (0,0) -- (10,0.5) -- (11,4) -- (10.5,7) -- (8,9) -- (5,10) -- (2,9) -- (0.5,6) -- (-0.5,3) -- cycle;
    
    % Villes repères
    \node[circle, fill=popblue, inner sep=2pt, label=below:\small Douala] (D) at (3.5, 2.5) {};
    \node[circle, fill=popblue, inner sep=2pt, label=above:\small Yaoundé] (Y) at (5.5, 5.5) {};
    \node[circle, fill=poporange, inner sep=2pt, label=right:\small Edéa] (E) at (4.5, 3.5) {};
    \node[circle, fill=popgreen, inner sep=2pt, label=left:\small Limbé] (L) at (2.0, 3.0) {};
    \node[circle, fill=poppurple, inner sep=2pt, label=right:\small Kribi] (K) at (6.0, 1.5) {};
    \node[circle, fill=poppink, inner sep=2pt, label=below:\small Buea/Tiko] (B) at (2.5, 4.0) {};
    \node[circle, fill=poppink, inner sep=2pt, label=above:\small Nkoteng] (N) at (6.5, 4.5) {};
    \node[circle, fill=poppink, inner sep=2pt, label=above:\small Mbanga] (M) at (4.0, 3.0) {};

    % Axes routiers principaux
    \draw[popline, very thick] (D) -- (Y) node[midway, above, sloped, font=\tiny] {N3};
    \draw[popline, thick] (D) -- (K) node[midway, below, sloped, font=\tiny] {N7 / Autop.};
    \draw[popline, thick] (D) -- (B);
    \draw[popline, thick] (Y) -- (N);
    \draw[popline, thick, dashed] (Y) -- (8,7) node[above, font=\tiny] {Vers Nord (N1)};
    
    % Chemin de fer
    \draw[popdark, decorate, decoration={snake, amplitude=1pt, segment length=4pt}, thick] (D) -- (Y) -- (8,7);
    \draw[popdark, decorate, decoration={snake, amplitude=1pt, segment length=4pt}, thick] (D) -- (N);

    % PÔLE 1 : DOUALA - Capitale économique, Port, Industrie diverse
    \node[draw=popblue, thick, rounded corners, fill=popblue!10, inner sep=4pt, anchor=south west, align=center] at (10.5, 2){
        \textbf{Pôle 1 : Douala - Bonabéri}\\
        \textcolor{popblue}{\rule{1.5ex}{1.5ex}} \textbf{Industrie} : Brasseries, Textile, Chimie, BTP, Mécanique\\
        \textcolor{popblue}{\rule{1.5ex}{1.5ex}} \textbf{Port} : 90\% trafic maritime\\
        \textcolor{popblue}{\rule{1.5ex}{1.5ex}} \textbf{Logistique} : Hub CEMAC
    };

    % PÔLE 2 : EDÉA - Énergie / Métallurgie
    \node[draw=poporange, thick, rounded corners, fill=poporange!10, inner sep=4pt, anchor=north west, align=center] at (10.5, 0){
        \textbf{Pôle 2 : Edéa}\\
        \textcolor{poporange}{\rule{1.5ex}{1.5ex}} \textbf{Énergie} : Barrages (Song Loulou, Edéa)\\
        \textcolor{poporange}{\rule{1.5ex}{1.5ex}} \textbf{Métallurgie} : ALUCAM (Aluminium)\\
        \textit{Électro-intensif}
    };

    % PÔLE 3 : YAOUNDÉ - Administration / Services / Agro-industrie périphérie
    \node[draw=poppurple, thick, rounded corners, fill=poppurple!10, inner sep=4pt, anchor=south east, align=center] at (-1.5, 7){
        \textbf{Pôle 3 : Yaoundé}\\
        \textcolor{poppurple}{\rule{1.5ex}{1.5ex}} \textbf{Administration / Décision}\\
        \textcolor{poppurple}{\rule{1.5ex}{1.5ex}} \textbf{Services / Finance / ICT}\\
        \textcolor{poppurple}{\rule{1.5ex}{1.5ex}} \textbf{Périphérie} : Nkoteng (SOSUCAM), Mbandjock
    };

    % PÔLE 4 : SUD-OUEST / CDC - Agro-industrie plantation
    \node[draw=popgreen, thick, rounded corners, fill=popgreen!10, inner sep=4pt, anchor=north east, align=center] at (-1.5, 0){
        \textbf{Pôle 4 : Sud-Ouest (CDC / PHP / PAMOL)}\\
        \textcolor{popgreen}{\rule{1.5ex}{1.5ex}} \textbf{Plantations} : Banane (PHP, CDC), Palmier (CDC, Pamol), Hévéa (CDC)\\
        \textcolor{popgreen}{\rule{1.5ex}{1.5ex}} \textbf{Port Limbé} : Export banane/bois\\
        \textcolor{popgreen}{\rule{1.5ex}{1.5ex}} \textbf{Emploi} : $\sim$ 25 000 permanents
    };

    % PÔLE 5 : KRIBI / SUD - Nouveau pôle énergétique / portuaire / minier
    \node[draw=poppink, thick, rounded corners, fill=poppink!10, inner sep=4pt, anchor=north west, align=center] at (10.5, -2.5){
        \textbf{Pôle 5 : Kribi / Sud (Nouveau)}\\
        \textcolor{poppink}{\rule{1.5ex}{1.5ex}} \textbf{Port en eau profonde} (KCT)\\
        \textcolor{poppink}{\rule{1.5ex}{1.5ex}} \textbf{Gaz / Énergie} : GNL, Centrale Kribi Power\\
        \textcolor{poppink}{\rule{1.5ex}{1.5ex}} \textbf{Mine} : Fer Mbalam-Nabéba (projet)\\
        \textcolor{poppink}{\rule{1.5ex}{1.5ex}} \textbf{ZFI} : Zone Franche Industrielle
    };

    % Flèche Nord
    \node[draw=popdark, fill=popblue!20, circle, inner sep=3pt] at (9.5, 9.5) {\textbf{N}};
    
    % Légende symboles
    \node[anchor=north west, font=\small, fill=white, fill opacity=0.8, draw=popline, rounded corners, align=center] at (-1.5, -1.5){
        \textbf{Légende} \\
        \textcolor{popblue}{\rule{1.5ex}{1.5ex}} Industrie lourde / Manufacturière \\
        \textcolor{poporange}{\rule{1.5ex}{1.5ex}} Énergie / Métallurgie \\
        \textcolor{poppurple}{\rule{1.5ex}{1.5ex}} Administration / Services supérieurs \\
        \textcolor{popgreen}{\rule{1.5ex}{1.5ex}} Agro-industrie (Plantations) \\
        \textcolor{poppink}{\rule{1.5ex}{1.5ex}} Mines / Port profond / ZFI \\
        \textcolor{popdark}{\rule{1.5ex}{1.5ex}} Chemin de fer (Camrail) \\
        \textcolor{popline}{\rule{1.5ex}{1.5ex}} Routes nationales bitumées
    };
\end{tikzpicture}
\legende{Localisation schématique des 5 grands pôles de l'économie moderne camerounaise. La concentration est forte sur l'axe Douala-Edéa-Yaoundé (triangle productif) et le bassin CDC (Sud-Ouest). Kribi émerge comme 5ème pôle stratégique (Port, Gaz, Mine, ZFI).}
\end{popfigure}

\begin{methode}[Réaliser un croquis de synthèse « Les pôles de l'économie moderne au Cameroun » (TP Intégré)]
\begin{enumerate}
    \item \textbf{Fond de carte} : Tracer le contour simplifié du Cameroun, les grands axes routiers (N1, N3, N7, autoroute), le chemin de fer (Douala-Yaoundé-Ngaoundéré, Douala-Nkongsamba/Kumba) et les deux ports (Douala, Kribi).
    \item \textbf{Repérage des pôles} : Positionner les 5 pôles (Douala, Edéa, Yaoundé, Sud-Ouest/CDC, Kribi) par des symboles de forme/couleur distincts (carré industrie, cercle énergie, triangle services, losange plantation, pentagone mine/port).
    \item \textbf{Activités dominantes} : Pour chaque pôle, inscrire 2-3 activités clés (ex: Douala → Port, Brasseries, Raffinage ; Edéa → ALUCAM, Barrages ; Kribi → Port profond, GNL, Fer Mbalam).
    \item \textbf{Flux et liaisons} : Représenter les flux d'exportation (flèches épaisses vers ports), les flux intrants (pétrole brut vers SONARA, bauxite/alumine vers ALUCAM), les flux de main-d'œuvre (flèches fines zones rurales → pôles).
    \item \textbf{Légende organisée} : Classer les symboles par thème (Industrie, Énergie, Agro-industrie, Mines/Port, Infrastructures). Indiquer l'échelle approximative et l'orientation (Nord).
    \item \textbf{Titre et conclusion} : Titre explicite. Une phrase de synthèse sur la concentration spatiale et l'émergence du pôle Kribi.
\end{enumerate}
\end{methode}

\begin{aretenir}[Enjeux transversaux : Agro-industrie, emploi, environnement, émergence]
\begin{itemize}
    \item \textbf{Emploi} : Le secteur moderne formel crée $\sim$ 150 000-200 000 emplois directs (dont $\sim$ 100 000 dans les plantations). Défi : passer de l'emploi précaire (saisonnier) à l'emploi stable et qualifié (transformation, maintenance, R\&D).
    \item \textbf{Environnement} : Déforestation (palmier, hévéa), pollution industrielle (SONARA, ALUCAM, CIMENCAM), érosion des sols (plantations en pente), conflits d'usage (aires protégées vs concessions). Nécessité de la certification (RSPO, FSC) et des Évaluations d'Impact Environnemental (EIE) rigoureuses.
    \item \textbf{Trajectoire d'émergence (SND30)} : Objectif : taux de transformation locale > 50\% pour le bois, le cacao, le coton ; développement de chaînes de valeur intégrées (ex: hévéa $\to$ gants/chambres à air) ; attractivité des ZFI pour industries manufacturières d'exportation ; maîtrise de l'énergie (barrages, gaz) pour compétitivité.
\end{itemize}
\end{aretenir}

\begin{exercice}[Application : Caractériser un pôle]
À partir de la carte et de vos connaissances, compléter le tableau ci-dessous pour le pôle de \textbf{Kribi} :
\begin{center}
\begin{tabularx}{\linewidth}{|l|X|}
\hline
\textbf{Domaine} & \textbf{Élément caractéristique (à compléter)} \\ \hline
Infrastructure portuaire & \\ \hline
Ressource énergétique & \\ \hline
Projet minier majeur & \\ \hline
Cadre incitatif spécifique & \\ \hline
Enjeu environnemental & \\ \hline
\end{tabularx}
\end{center}
\end{exercice}

\begin{corrige}[Exercice n : Caractériser le pôle de Kribi]
\begin{center}
\begin{tabularx}{\linewidth}{|l|X|}
\hline
\textbf{Domaine} & \textbf{Élément caractéristique} \\ \hline
Infrastructure portuaire & Port en eau profonde de Kribi (KCT), capacité 1,4 M EVP, quai 16-18m, hub transbordement Golfe de Guinée \\ \hline
Ressource énergétique & Gisements gaziers offshore (Sanaga Sud, SNH/Perenco/Glencore), Usine GNL (FLNG), Centrale thermique Kribi Power (216 MW), Projet Centrale Mpolongwé \\ \hline
Projet minier majeur & Fer de Mbalam-Nabéba (Sundance Resources / AustSino / Bestway Finance) : 700-800 Mt, chemin de fer dédié 510 km $\to$ Kribi, terminal minéralier prévu \\ \hline
Cadre incitatif spécifique & Zone Franche Industrielle (ZFI) Kribi : exonérations fiscales/douanières totales pour exportateurs > 80\%, guichet unique API \\ \hline
Enjeu environnemental & Parc national de Campo Ma'an (biodiversité, éléphants, gorilles), mangroves, risques pollution portuaire/gazière, conflits usage sol (agro-industrie HEVECAM/SOCAPALM vs mine/port) \\ \hline
\end{tabularx}
\end{center}
\end{corrige}

\coursec{L'agriculture de grandes plantations : Genèse, acteurs et statuts fonciers}

Dans la section précédente, nous avons montré que l'économie camerounaise est marquée par un \cle{dualisme} profond : d'un côté, un secteur traditionnel (agriculture vivrière, artisanat, petit commerce) qui occupe l'immense majorité de la population ; de l'autre, un secteur moderne, capitaliste, intégré aux marchés mondiaux, qui ne concerne qu'une minorité mais fournit l'essentiel des exportations agricoles et une part importante des recettes fiscales. C'est au sein de ce secteur moderne que se situe l'\cle{agriculture de grandes plantations}, forme d'exploitation agricole héritée de la colonisation et encore aujourd'hui au cœur des débats économiques, fonciers et environnementaux du pays. Cette section en retrace la genèse, identifie les acteurs et analyse le statut foncier qui la rend possible.

\courssub{Définition et caractéristiques de la grande plantation}

\begin{definition}[Grande plantation]
Une \cle{grande plantation} (ou agro-industrie) est une exploitation agricole de très grande superficie (de plusieurs centaines à plusieurs dizaines de milliers d'hectares), spécialisée dans une ou quelques \cle{spéculations} d'exportation (banane, palmier à huile, hévéa, canne à sucre, thé), qui combine en un même ensemble : la culture sur de vastes étendues, une unité de \cle{transformation industrielle} sur place (huilerie, usine de traitement du latex, sucrerie, station de conditionnement) et une main-d'œuvre salariée nombreuse. Elle est généralement détenue par une société de capitaux et orientée vers l'exportation.
\end{definition}

L'intuition est simple : contrairement au paysan qui cultive 1 à 3 ha pour nourrir sa famille et vendre le surplus, la plantation fonctionne comme une \emph{usine à ciel ouvert}. Elle mobilise des capitaux importants, une organisation hiérarchisée du travail, des infrastructures propres (routes, cités ouvrières, hôpitaux, écoles) et des débouchés extérieurs sécurisés par des contrats internationaux. Au Cameroun, elle repose sur un pilier juridique particulier : la terre n'appartient pas à l'exploitant, elle lui est \emph{concédée} par l'État.

\begin{definition}[Concession foncière et bail emphytéotique]
La \cle{concession foncière} est un acte par lequel l'État, propriétaire d'une terre domaniale, en confie l'usage à une personne morale (société) pour une durée déterminée, en échange d'une redevance et d'obligations (mise en valeur, emplois, infrastructures). Le \cle{bail emphytéotique} est la forme typique de cette concession : c'est un droit réel immobilier de longue durée (jusqu'à 99 ans au Cameroun), moyennant une \cle{redevance foncière} annuelle. Le preneur jouit de la terre comme un propriétaire, mais ne peut la vendre ; à l'expiration du bail, la terre et les améliorations reviennent à l'État (clause de réversion).
\end{definition}

\courssub{Genèse : des origines coloniales aux privatisations}

\textbf{1. Les origines allemandes (1884-1916).}\par
Dès les années 1890, l'administration coloniale allemande crée les premières plantations du Cameroun, attirée par les sols volcaniques fertiles et le climat humide du \cle{Sud-Ouest} (région de Victoria, actuelle Limbé, et de Buea) et du \cle{Littoral}. Des firmes comme la \emph{Westafrikanische Pflanzungsgesellschaft Victoria} (WAPV, 1892) mettent en culture la banane, le cacao, le palmier à huile et le caoutchouc. Le modèle est déjà celui de la concession : de vastes domaines concédés par le pouvoir colonial, une main-d'œuvre recrutée parfois de force, et une production intégralement exportée vers l'Europe.

\textbf{2. L'héritage mandataire (1916-1960).}\par
Après la défaite allemande, le Cameroun est partagé entre la France et le Royaume-Uni sous mandat de la SDN puis tutelle de l'ONU. Chaque puissance prolonge le modèle :
\begin{itemize}
\item côté britannique, les anciennes plantations allemandes confisquées sont regroupées en 1946-1947 au sein de la \cle{CDC} (\emph{Cameroon Development Corporation}), entreprise publique qui devient le deuxième employeur du pays après l'État ;
\item côté français, l'État encourage la création de sociétés d'économie mixte : \cle{SOCAPALM} (palmier à huile, 1963), \cle{SOSUCAM} (sucre, 1965), puis après l'indépendance \cle{HEVECAM} (hévéa, 1975) et l'\cle{UNVDA} (riziculture dans la plaine de Ndop, au Nord-Ouest).
\end{itemize}
Le modèle repose sur deux piliers : la \emph{concession foncière de longue durée} (jusqu'à 99 ans) et une \emph{main-d'œuvre migrante}, recrutée dans les régions du Nord et dans les pays voisins (Nigeria, Tchad), logée dans des camps de plantation.

\textbf{3. L'évolution post-indépendance.}\par
Après 1960, l'État camerounais nationalise ou crée des \cle{sociétés d'État} (SODECAO pour le cacao, HEVECAM, CDC, SOSUCAM) dans le cadre du « libéralisme planifié ». La crise économique des années 1980-1990 (effondrement des cours du cacao, du café et du pétrole) conduit aux \cle{Plans d'ajustement structurel} (PAS) imposés par le FMI et la Banque mondiale, qui exigent le désengagement de l'État. Dans les années 2000, la plupart des agro-industries sont \cle{privatisées} partiellement ou totalement : SOCAPALM et SAFACAM passent au groupe SOCFIN (lié à Bolloré), HEVECAM au groupe GMG puis Halcyon Agri, SOSUCAM au groupe SOMDIAA (Castel).

\begin{popfigure}
\begin{tikzpicture}[x=1cm,y=1cm,>=Stealth]
\draw[line width=1.2pt,popline] (0,0) -- (13.2,0);
\foreach \x/\y/\titre/\txt/\col in {
  0.6/1.1/{1890}/{1\textsuperscript{res} plantations allemandes (Victoria, Buea)}/popblue,
  3.2/-1.1/{1946-1947}/{Création de la CDC (zone britannique)}/poporange,
  5.8/1.1/{1963-1975}/{SOCAPALM, SOSUCAM, HEVECAM}/popgreen,
  8.6/-1.1/{1990-2005}/{PAS et privatisations (SOCFIN, Castel)}/poppink,
  11.4/1.1/{2010-2020}/{Certifications RSPO, Fairtrade}/poppurple}
{
  \fill[\col] (\x,0) circle (0.12);
  \draw[\col,line width=0.8pt] (\x,0) -- (\x,\y*0.55);
  \node[anchor=south west,align=left,text width=2.3cm,font=\scriptsize] at (\x-0.5,\y*0.55+0.05) {\textbf{\titre}\\ \txt};
}
\end{tikzpicture}
\legende{Frise chronologique : les grandes étapes de l'agriculture de plantation au Cameroun (1890-2020).}
\end{popfigure}

\courssub{Les acteurs actuels}

Le paysage des plantations camerounaises associe aujourd'hui trois catégories d'acteurs :

\begin{itemize}
\item les \cle{multinationales}, détentrices des capitaux et des débouchés : le groupe \textbf{SOCFIN} (lié à Bolloré) contrôle SOCAPALM, SAFACAM (en partenariat avec Wilmar) et une partie de l'hévéaculture ; les \textbf{Plantations du Haut-Penja} (PHP, groupe Compagnie Fruitière) dominent la banane d'exportation ; SOMDIAA contrôle SOSUCAM ;
\item l'\cle{État}, qui reste propriétaire du foncier, actionnaire (majoritaire à la CDC et à SOSUCAM, minoritaire ailleurs) et garant du respect des cahiers des charges ;
\item les \cle{planteurs villageois} (\emph{outgrowers}) : paysans installés autour des plantations, qui cultivent sous contrat pour la firme (plants, intrants et encadrement fournis, production rachetée à prix convenu).
\end{itemize}

\begin{exemplebox}[Quelques chiffres repères]
\begin{center}\begin{tabularx}{\linewidth}{|l|X|l|}\hline
\rowcolor{popblueL} \textbf{Firme} & \textbf{Spéculation} & \textbf{Surface} \\ \hline
SOCAPALM & Palmier à huile (Littoral, Sud-Ouest, Centre) & $\approx$ \SI{58000}{ha} \\ \hline
CDC & Banane, caoutchouc, palme (Sud-Ouest) & $\approx$ \SI{40000}{ha} \\ \hline
HEVECAM & Hévéa (Sud, région de Kribi-Niete) & $\approx$ \SI{40000}{ha} \\ \hline
SOSUCAM & Canne à sucre (Moungo, Bénoué) & $\approx$ \SI{20000}{ha} \\ \hline
PHP & Banane, ananas (Littoral, Penja) & $\approx$ \SI{3000}{ha} \\ \hline
\end{tabularx}\end{center}
\end{exemplebox}

\begin{savaistu}[Le modèle « noyau central + planteurs villageois »]
Chez SOCAPALM et HEVECAM, la plantation industrielle (le « noyau ») n'est pas seule : elle est entourée de \cle{planteurs villageois} qui fournissent une part importante (jusqu'à 30 à 40\,\% selon les sites) des régimes de palme ou du latex traités par les usines. Ce modèle, dit \emph{nucleus estate}, permet à la firme d'augmenter sa production sans agrandir ses concessions, tout en encadrant les paysans (plants sélectionnés, formation, prix garanti).
\end{savaistu}

\begin{attention}[Ne pas confondre planteurs villageois et petits producteurs indépendants]
Erreur fréquente à l'examen : assimiler les \cle{planteurs villageois} (outgrowers) aux petits producteurs indépendants. Le planteur villageois est \emph{lié par contrat} à une firme : il reçoit plants, intrants et encadrement technique, mais il est \emph{obligé de vendre} sa production à cette firme, au prix fixé par elle. Le petit producteur indépendant, lui, est libre de vendre où il veut (collecteurs, marchés locaux, coopératives). Cette distinction est essentielle pour analyser les rapports de dépendance dans l'agro-industrie.
\end{attention}

\courssub{Le statut foncier : un problème juridique et social}

\begin{propriete}[Évolution du statut juridique des terres de plantation]
Sous la colonisation, les terres dites « vacantes et sans maître » sont déclarées \cle{domaine public}. Après l'indépendance, les ordonnances de 1974 (ordonnance n°74/1) créent le \cle{domaine national} et le domaine privé de l'État, seul susceptible d'être concédé. La réglementation postérieure (décrets des années 2000) encadre les \emph{concessions à des personnes morales de droit privé} : bail emphytéotique de longue durée (jusqu'à 99 ans), redevance foncière, cahier des charges obligatoire. À l'expiration du bail, la terre revient à l'État (réversion).
\end{propriete}

Ce cadre juridique génère trois problèmes majeurs :
\begin{enumerate}
\item l'\textbf{expiration des baux} : plusieurs concessions coloniales ou mandataires arrivent à échéance, posant la question de leur renouvellement, de leur rétrocession aux communautés ou de leur remise en adjudication ;
\item le \textbf{chevauchement avec les terres coutumières} : les concessions ont souvent été tracées sur des terres que les communautés locales considèrent comme leurs terres ancestrales, sans consultation ni indemnisation réelle ;
\item les \textbf{revendications autochtones} : les populations Bakossi (contre SOCAPALM), Bassa et Douala réclament la restitution de terres, des compensations ou une renégociation des cahiers des charges.
\end{enumerate}

\begin{methode}[Analyser un extrait de cahier des charges de concession]
\begin{enumerate}
\item \textbf{Identifier les parties} : qui concède (l'État, quel ministère) ? Qui reçoit (quelle société, quels actionnaires) ?
\item \textbf{Relever les données foncières} : superficie concédée (en ha), localisation, durée du bail, montant de la redevance annuelle (en FCFA/ha).
\item \textbf{Repérer les obligations sociales} : emplois créés, construction d'écoles, de dispensaires, de routes, logement des travailleurs.
\item \textbf{Repérer les obligations environnementales} : zones tampons, protection des cours d'eau, certification (RSPO pour le palmier), interdiction de déforestation.
\item \textbf{Chercher la clause de réversion} : que deviennent terres et infrastructures à l'expiration du bail ?
\item \textbf{Conclure de façon critique} : le contrat est-il équilibré ? Les communautés locales sont-elles mentionnées ? Le contrat est-il publié (transparence) ?
\end{enumerate}
\end{methode}

\courssub{Répartition spatiale des spéculations}

La localisation des plantations n'est pas aléatoire : elle suit les conditions écologiques (pluies abondantes, sols profonds) et l'héritage colonial (proximité des ports de Douala et Limbé, des axes routiers et ferroviaires).

\begin{popfigure}
\begin{tikzpicture}[scale=1.05,>=Stealth]
% Contour simplifié du Cameroun
\draw[fill=popcream,line width=1pt,popdark] (2.2,0.2) -- (1.2,1.6) -- (0.9,3.2) -- (1.6,4.6) -- (2.4,5.6) -- (2.9,7.2) -- (3.6,8.6) -- (4.4,9.6) -- (5.2,9.4) -- (5.0,8.2) -- (5.8,7.0) -- (6.4,5.8) -- (7.6,5.0) -- (8.0,3.6) -- (7.2,2.2) -- (6.0,1.4) -- (4.6,0.6) -- (3.2,0.1) -- cycle;
% Nord
\draw[->,line width=1pt,popdark] (8.6,9.2) -- (8.6,10.0) node[above,font=\scriptsize]{N};
% Villes
\fill[popdark] (1.35,2.6) circle (0.07) node[right,font=\scriptsize]{Douala};
\fill[popdark] (0.95,1.9) circle (0.06) node[right,font=\scriptsize]{Limbé};
\fill[popdark] (1.7,3.1) circle (0.06) node[right,font=\scriptsize]{Buea};
\fill[popdark] (3.6,2.6) circle (0.07) node[right,font=\scriptsize]{Yaoundé};
\fill[popdark] (2.6,1.2) circle (0.06) node[right,font=\scriptsize]{Kribi};
\fill[popdark] (2.1,3.4) circle (0.05) node[right,font=\scriptsize]{Penja};
% Plantations (points proportionnels)
\fill[popgreen,opacity=0.85] (1.5,2.2) circle (0.42);
\fill[poporange,opacity=0.85] (1.1,2.0) circle (0.35);
\fill[poppurple,opacity=0.85] (2.9,1.5) circle (0.35);
\fill[popgold,opacity=0.85] (2.0,3.2) circle (0.30);
\fill[poppink,opacity=0.85] (3.0,3.0) circle (0.22);
\fill[popblue,opacity=0.85] (1.9,4.4) circle (0.22);
% Axes
\draw[popmuted,line width=0.9pt] (1.35,2.6) -- (3.6,2.6);
\draw[popmuted,line width=0.9pt,dashed] (1.35,2.6) -- (2.6,1.2);
\draw[popmuted,line width=0.9pt] (1.35,2.6) -- (2.1,3.4);
% Légende
\node[anchor=north west,font=\scriptsize,align=left] at (5.6,4.4) {%
\textcolor{popgreen}{$\bullet$} Palmier à huile (SOCAPALM)\\
\textcolor{poporange}{$\bullet$} Banane / caoutchouc / palme (CDC)\\
\textcolor{poppurple}{$\bullet$} Hévéa (HEVECAM)\\
\textcolor{popgold}{$\bullet$} Banane / ananas (PHP)\\
\textcolor{poppink}{$\bullet$} Canne à sucre (SOSUCAM)\\
\textcolor{popblue}{$\bullet$} Thé (Nord-Ouest, Tole, Ndu)\\
\textcolor{popmuted}{---} route \quad \textcolor{popmuted}{- -} rail};
\end{tikzpicture}
\legende{Croquis schématique : localisation des grandes plantations par spéculation (surfaces approximatives, points proportionnels). Les plantations se concentrent dans le Sud forestier humide, près des ports et des axes de transport ; seul le thé occupe les hauts plateaux du Nord-Ouest.}
\end{popfigure}

\begin{center}\begin{tabularx}{\linewidth}{|l|X|}\hline
\rowcolor{popblueL} \textbf{Spéculation} & \textbf{Localisation principale} \\ \hline
Banane & Sud-Ouest (CDC), Moungo (PHP, Penja) \\ \hline
Caoutchouc naturel & Sud-Ouest, Centre, Sud (HEVECAM, CDC, SAFACAM) \\ \hline
Palmier à huile & Littoral, Sud-Ouest, Centre (SOCAPALM) \\ \hline
Sucre & Moungo et Bénoué (SOSUCAM) \\ \hline
Thé & Nord-Ouest (hauts plateaux de Tole, Ndu) \\ \hline
Ananas & Littoral (PHP) \\ \hline
\end{tabularx}\end{center}

\begin{exoresolu}[Calculer et comparer les emprises foncières]
Énoncé : SOCAPALM occupe environ \SI{58000}{ha}, la CDC \SI{40000}{ha} et PHP \SI{3000}{ha}. (1) Quelle part la surface de PHP représente-t-elle par rapport à celle de SOCAPALM ? (2) Quelle surface totale ces trois firmes occupent-elles ?
\tcblower
Solution : (1) Part de PHP par rapport à SOCAPALM :
\[ \frac{3000}{58000} \times 100 \approx 5,2\,\% \]
PHP n'occupe donc qu'environ 5\,\% de la surface de SOCAPALM, malgré son poids dans la banane d'exportation. (2) Surface totale :
\[ 58000 + 40000 + 3000 = \SI{101000}{ha} \]
soit environ \SI{1010}{km^2} (rappel : $1~\mathrm{km}^2 = 100$ ha), l'équivalent d'un département entier. Conclusion : l'agro-industrie concentre d'immenses emprises foncières entre quelques firmes, ce qui explique l'acuité des conflits fonciers avec les communautés riveraines.
\end{exoresolu}

\begin{aretenir}[L'essentiel]
L'agriculture de grandes plantations au Cameroun est un héritage colonial (Allemands, puis mandats français et britannique), fondé sur la \cle{concession} de terres domaniales par bail emphytéotique (jusqu'à 99 ans). Après les nationalisations des années 1970 et la crise des années 1980-1990, les PAS ont conduit aux privatisations des années 2000 au profit de multinationales (SOCFIN/Bolloré, Wilmar, Castel). Les spéculations dominantes (banane, palme, hévéa, sucre, thé) se concentrent dans le Sud humide, près des ports. Le modèle « noyau central + planteurs villageois » et les conflits fonciers avec les communautés (Bakossi, Bassa, Douala) en font un secteur économiquement stratégique mais socialement contesté.
\end{aretenir}

\begin{exercice}[Pour s'entraîner]
(1) Définis : concession foncière, bail emphytéotique, planteur villageois. (2) Reconstitue la chronologie : création de la CDC, premières plantations allemandes, privatisations, création de SOCAPALM. (3) Explique en cinq lignes pourquoi les concessions de SOCAPALM font l'objet de revendications des populations Bakossi. (4) À l'aide du croquis, justifie la concentration des plantations dans les régions du Littoral et du Sud-Ouest (deux arguments : un naturel, un historique).
\end{exercice}

\coursec{Étude de cas guidée : Organisation spatiale et enjeux d'une firme agro-industrielle (CDC / SOCAPALM / HEVECAM)}

Nous avons vu, dans la section précédente, comment les grandes plantations sont nées au Cameroun, qui en sont les acteurs (État, firmes, planteurs villageois) et sur quels statuts fonciers elles reposent (concessions, baux emphytéotiques). Mais une plantation n'est pas seulement un champ : c'est un \cle{espace organisé}, une véritable petite ville industrielle greffée sur la forêt ou la savane. Pour comprendre concrètement cette organisation, rien ne vaut une \cle{étude de cas}. Nous choisissons ici la \cle{CDC} (\emph{Cameroon Development Corporation}), firme la plus complexe et la plus ancienne (multi-spéculations : banane, caoutchouc, palme), en la comparant ponctuellement à la \cle{SOCAPALM}, leader camerounais de l'huile de palme, et à \cle{HEVECAM} pour l'hévéaculture.

\courssub{Pourquoi étudier une firme agro-industrielle comme un « système spatial » ?}

\begin{definition}[Firme agro-industrielle]
Une \cle{firme agro-industrielle} est une entreprise qui combine, sur un même territoire concédé, la \cle{production agricole} (plantations), la \cle{transformation industrielle} (usines de première transformation) et la \cle{logistique d'exportation} (routes, rail, port). Elle contrôle donc l'ensemble de la \cle{chaîne de valeur}, du plant au produit exporté.
\end{definition}

L'intuition est simple : un régime de bananes ou une grappe de palme pourrit en quelques jours. Il faut donc que l'usine soit \emph{à côté} du champ, la route \emph{à côté} de l'usine, et le port \emph{au bout} de la route. Toute l'organisation spatiale de la firme découle de cette contrainte de \cle{périssabilité} et de \cle{masse} (tonnages énormes à évacuer).

\begin{exemplebox}[Deux géants chiffrés]
\begin{itemize}
\item \textbf{CDC} : environ \num{22000} employés permanents (premier employeur privé du Cameroun) ; production d'environ \num{180000} tonnes de bananes par an et \num{25000} tonnes de caoutchouc par an ; siège social à \textbf{Buea}, siège opérationnel et portuaire à \textbf{Bota-Limbé}, plantations dans le Fako et la Meme.
\item \textbf{SOCAPALM} : environ \num{6000} employés, près de \num{180000} tonnes d'huile brute de palme par an, réparties sur plusieurs sites (Edéa, Dibombari, Mbongo, Mbambé, Apouh, Kienké...).
\end{itemize}
\end{exemplebox}

\courssub{La structure spatiale de la firme : une organisation en réseau hiérarchisé}

Observons une carte d'une concession CDC ou SOCAPALM : on y voit toujours la même logique en quatre niveaux.

\begin{enumerate}
\item \textbf{Le siège social et la direction générale} : cerveau de la firme, il est excentré par rapport aux champs — à \textbf{Buea} pour la CDC (direction générale), \textbf{Bota-Limbé} pour la direction des opérations portuaires, à \textbf{Douala} pour la direction générale de SOCAPALM (groupe Socfin). C'est là que se prennent les décisions d'investissement et que transite l'information.
\item \textbf{Les unités de production (estates ou divisions)} : la concession est découpée en grandes unités (ex. : \emph{Tiko, Missellele, Mondoni} pour la banane CDC ; \emph{Mbongo, Edéa} pour SOCAPALM), chacune gérée par un directeur d'exploitation et subdivisée en blocs de culture.
\item \textbf{Les usines de première transformation} : \cle{huileries} de palme (grappes $\rightarrow$ huile brute), usines de \cle{caoutchouc} (latex $\rightarrow$ plaques ou granulés), \cle{pack-houses} bananiers (régimes $\rightarrow$ cartons de bananes contrôlées). Elles se placent au centre de gravité des champs pour limiter le transport de matière fraîche.
\item \textbf{Les infrastructures d'évacuation} : pistes internes et routes goudronnées, réseau ferroviaire \textbf{interne à la CDC} (évacuation banane/caoutchouc vers Limbé/Douala), ligne de chemin de fer national \cle{Camrail} (Douala–Yaoundé–Ngaoundéré, utilisé pour les flux Nord-Sud), et surtout les ports : \textbf{Douala} (principal), \textbf{Limbé} et \textbf{Kribi} (port en eau profonde, débouché croissant). Les bananes partent en navires réfrigérés, l'huile et le caoutchouc en vrac ou en conteneurs.
\end{enumerate}

\begin{attention}[Ne pas confondre siège et site de production]
Au Bac, beaucoup de copies placent le siège social « au milieu des plantations ». Erreur : la direction se situe dans une \textbf{ville} (accès aux banques, aux administrations, au port), tandis que les usines se situent \textbf{dans la concession}, au plus près des champs. C'est la distinction classique entre \cle{fonction de commandement} (urbaine) et \cle{fonction de production} (rurale).
\end{attention}

\courssub{Le schéma fonctionnel : flux de matière, d'énergie et d'information}

Pour synthétiser, on représente la plantation comme un \cle{système} : des flux entrent (intrants, capitaux, main-d'œuvre), circulent (matière récoltée, informations de gestion) et sortent (produits exportés ou vendus localement), avec des \emph{boucles de retour} (encadrement des planteurs villageois, réinvestissement).

\begin{popfigure}
\begin{center}
\begin{tikzpicture}[font=\small,
  bloc/.style={draw=popdark, thick, rounded corners=2pt, fill=popblueL, minimum width=2.6cm, minimum height=0.9cm, align=center},
  flux/.style={-{Stealth[length=2.5mm]}, very thick, popblue},
  retour/.style={-{Stealth[length=2.5mm]}, thick, dashed, poporange}]
\node[bloc, fill=popgreenL, align=center] (intrants) at (0,0){Intrants\\ (plants, engrais,\\ pesticides)};
\node[bloc, fill=popgreenL, align=center] (champ) at (4.2,0){Champs\\ (estates)\\ + main-d'œuvre};
\node[bloc, fill=popgoldL, align=center] (usine) at (8.4,0){Usine de 1\textsuperscript{re}\\ transformation\\ (huilerie, pack-house)};
\node[bloc, fill=poppurpleL, align=center] (port) at (12.6,0){Port\\ (Douala, Limbé,\\ Kribi)};
\node[bloc, fill=poppinkL, align=center] (export) at (12.6,-2.6){Exportation /\\ marché local};
\node[bloc, fill=poporangeL, align=center] (pv) at (4.2,-2.6){Planteurs\\ villageois\\ (outgrowers)};
\node[bloc, fill=popcream, align=center] (siege) at (0,-2.6){Siège / Direction\\ (Buea, Douala,\\ Yaoundé)};
\draw[flux] (intrants) -- (champ);
\draw[flux] (champ) -- node[above, font=\footnotesize]{récolte fraîche} (usine);
\draw[flux] (usine) -- node[above, font=\footnotesize]{route / rail interne} (port);
\draw[flux] (port) -- (export);
\draw[flux] (pv) -- node[right, font=\footnotesize]{livraison à l'usine} (champ);
\draw[retour] (siege) -- node[left, font=\footnotesize]{capitaux, décisions} (intrants);
\draw[retour] (siege.east) -- node[below, font=\footnotesize]{intrants + encadrement technique} (pv.west);
\draw[retour] (export.west) -- ++(-1.2,0) |- node[pos=0.25, right, font=\footnotesize]{devises} (siege.east);
\end{tikzpicture}
\end{center}
\legende{Schéma fonctionnel d'une plantation agro-industrielle type : flux de matière (flèches pleines bleues), flux d'information et de capitaux (flèches pointillées orange). La boucle « planteurs villageois » montre l'approvisionnement complémentaire de l'usine.}
\end{popfigure}

\courssub{L'organisation sociale : la plantation, une ville dans la brousse}

Une grande plantation emploie des milliers de personnes qui vivent sur place. Elle doit donc assumer des fonctions que, normalement, l'État assure : logement, santé, éducation. C'est ce qu'on appelle la \cle{cité ouvrière}.

\begin{itemize}
\item \textbf{La cité ouvrière} : camps de logements (souvent hiérarchisés : « camps staff » pour les cadres, camps ouvriers), écoles, dispensaires et hôpitaux de plantation (l'hôpital de la CDC à Tiko est réputé), marchés, terrains de sport, lieux de culte. La firme devient un \cle{État dans l'État}.
\item \textbf{L'encadrement technique} : ingénieurs agronomes, chefs de culture, contremaîtres ; ils contrôlent la qualité, planifient les récoltes et gèrent les intrants.
\item \textbf{La main-d'œuvre} : d'une part les \cle{permanents} (cueilleurs de palme, coupeurs de bananes, saigneurs d'hévéa, ouvriers d'usine), d'autre part les \cle{saisonniers} recrutés aux pics de récolte — souvent des migrants venus du Nord ou des pays voisins.
\item \textbf{Les planteurs villageois} : paysans riverains intégrés au circuit de la firme.
\end{itemize}

\begin{definition}[Planteurs villageois / Outgrowers]
Les \cle{planteurs villageois} (ou \emph{outgrowers}) sont des producteurs indépendants, installés autour de la concession, qui cultivent sous \cle{contrat d'approvisionnement} avec la firme : celle-ci leur fournit plants sélectionnés, intrants à crédit et encadrement technique ; en échange, ils s'engagent à livrer toute leur récolte à l'usine, à un \cle{prix fixé} par la firme. Ce système augmente l'approvisionnement de l'usine sans agrandir la concession.
\end{definition}

\courssub{La chaîne de valeur : du champ au navire}

Suivons une grappe de palme chez SOCAPALM : elle illustre parfaitement les étapes de la \cle{chaîne de valeur}.

\begin{enumerate}
\item \textbf{Plantation} : préparation du sol, mise en place des plants, entretien, fertilisation (intrants), main-d'œuvre permanente.
\item \textbf{Récolte} : coupe des régimes de noix de palme mûrs, toute l'année.
\item \textbf{Transport interne} : évacuation rapide par pistes et bennes vers l'huilerie (délai maximal de quelques heures, sinon l'acidité de l'huile augmente).
\item \textbf{Transformation primaire} : stérilisation, pressage $\rightarrow$ \cle{huile brute de palme} (CPO) et palmiste. À la CDC : latex $\rightarrow$ caoutchouc ; régimes de bananes $\rightarrow$ lavage, tri, mise en cartons au pack-house.
\item \textbf{Export ou marché local} : huile brute en vrac vers Douala/Kribi ; bananes en navires réfrigérés vers l'Europe (certification Fairtrade/Max Havelaar pour la filière banane CDC « Boh Plantations ») ; une partie de l'huile est raffinée localement (raffineries de la zone de Douala, ex. Olavea) pour le marché national.
\end{enumerate}

\begin{exoresolu}[Calculer un taux d'extraction]
Une huilerie SOCAPALM reçoit \num{900000} tonnes de régimes de palme frais par an et produit environ \num{180000} tonnes d'huile brute. Calculer le taux d'extraction (rendement usine) en pourcentage.
\tcblower
Le taux d'extraction est le rapport entre la masse d'huile obtenue et la masse de régimes traités :
\[ \text{Taux} = \frac{\num{180000}}{\num{900000}} \times 100 = 20\ \%. \]
L'huilerie extrait donc \textbf{20 \%} de la masse des régimes, ce qui correspond au rendement moyen d'une huilerie moderne (18 à 22 \%). Ce chiffre montre que 80 \% de la matière transportée devient des \emph{déchets} (fibres, coques, effluents) : d'où l'importance de localiser l'usine au plus près des champs et de valoriser ces résidus.
\end{exoresolu}

\courssub{Enjeux environnementaux et certifications}

La conversion de la forêt en plantation a un coût écologique que les marchés internationaux ne tolèrent plus sans garanties.

\begin{itemize}
\item \textbf{Déforestation historique} : les concessions ont été ouvertes sur forêt dense (Fako, Meme, Sanaga-Maritime). Les nouvelles extensions doivent désormais épargner les zones \cle{HCS} (\emph{High Carbon Stock}, forêts riches en carbone) et \cle{HCV} (\emph{High Conservation Value}, forêts à forte valeur de conservation : biodiversité, sites culturels).
\item \textbf{Gestion des effluents} : les huileries rejettent le \cle{POME} (\emph{Palm Oil Mill Effluent}), un liquide très chargé en matière organique qui, déversé dans les rivières, asphyxie les cours d'eau.
\item \textbf{Pesticides et engrais} : risques de pollution des sols et des nappes, d'où des programmes d'agriculture raisonnée.
\item \textbf{Certifications} : \cle{RSPO} (\emph{Roundtable on Sustainable Palm Oil}) pour l'huile de palme durable (SOCAPALM engagée dans ce processus), norme \cle{ISO 14001} (management environnemental), label \cle{Fairtrade/Max Havelaar} pour la banane (filière banane de la CDC « Boh Plantations »), qui garantit un prix plancher et une prime de développement aux travailleurs.
\end{itemize}

\begin{savaistu}[Des déchets qui éclairent des villages]
Les effluents d'huilerie (POME), longtemps rejetés dans les rivières, sont de plus en plus \cle{méthanisés} : enfermés dans des bassins couverts, ils dégagent du \cle{biogaz} brûlé pour produire de l'électricité qui alimente l'usine, voire les villages voisins. Les fibres et coques de palmiste servent aussi de combustible pour les chaudières, et les résidus compostés retournent aux champs comme engrais organique. L'huilerie devient ainsi une mini-centrale électrique.
\end{savaistu}

\courssub{Conflits fonciers actuels : quand le bail arrive à expiration}

Le foncier est le point le plus sensible. Les concessions de la CDC et de SOCAPALM reposent sur des \cle{baux emphytéotiques} de longue durée consentis par l'État ; plusieurs arrivent à échéance ou font l'objet de renégociation. Trois acteurs s'affrontent ou négocient :

\begin{itemize}
\item \textbf{L'État}, propriétaire du domaine national, qui peut renouveler, réduire ou reprendre les assiettes foncières ;
\item \textbf{les firmes}, qui veulent sécuriser leurs investissements ;
\item \textbf{les communautés riveraines} (ex. : villages de \textbf{Mbonge}, chefferies \textbf{Bakossi} autour des concessions CDC ; villages riverains de SOCAPALM Edéa ou Mbongo), qui revendiquent la restitution de terres coutumières devenues insuffisantes pour leurs champs vivriers.
\end{itemize}

Les négociations en cours portent sur la \cle{réduction de l'assiette foncière} (rétrocession de milliers d'hectares aux communautés), la création de \cle{fonds de développement local} alimentés par la firme, et l'encadrement des planteurs villageois comme compromis entre emploi et terre.

\begin{attention}[Autochtones et allogènes : ne pas tout mélanger]
Dans un conflit foncier de plantation, il faut distinguer deux groupes : les \cle{populations autochtones}, qui tiennent la terre de leurs ancêtres et invoquent les \textbf{droits coutumiers} (ex. : Bakossi, Bakweri face à la CDC), et les \cle{populations allogènes}, main-d'œuvre migrante installée dans les cités ouvrières, qui réclament surtout des \textbf{droits sociaux} (salaires, logement, retraite). Confondre les deux dans une copie, c'est fausser l'analyse du conflit.
\end{attention}

\courssub{Méthode : réaliser le croquis d'organisation spatiale d'une plantation}

\begin{methode}[Cinq étapes pour un croquis de plantation réussi]
\begin{enumerate}
\item \textbf{Délimiter la concession} : tracer le contour de l'emprise (trait plein), placer les villages riverains à l'extérieur ou en lisière, les cours d'eau qui la traversent.
\item \textbf{Localiser les noyaux d'habitat et les usines} : cités ouvrières (points noirs), huilerie ou pack-house (carrés), au centre de gravité des champs.
\item \textbf{Figurer les cultures par aplats} : un aplat de couleur par spéculation (palmeraie, bananeraie, hévéas), sans chercher la précision parcelle par parcelle.
\item \textbf{Figurer les réseaux de transport} : routes (traits doubles ou rouges), voie ferrée interne (trait noir crénelé), pistes internes (pointillés), et flécher l'évacuation vers le port.
\item \textbf{Organiser la légende hiérarchisée} : regrouper les signes par rubriques (Agriculture / Industrie / Infrastructures / Habitat / Conflits), du plus important au moins important.
\end{enumerate}
\end{methode}

\begin{popfigure}
\begin{center}
\begin{tikzpicture}[font=\footnotesize, scale=0.95]
% contour concession
\draw[very thick, popdark, fill=popcream!60] (0.3,0.5) -- (2.2,0.2) -- (5.5,0.4) -- (8.6,0.8) -- (9.3,2.6) -- (8.4,5.2) -- (5.8,6.1) -- (3.1,5.7) -- (0.8,4.6) -- (0.2,2.4) -- cycle;
% cours d'eau
\draw[thick, popblue, decorate, decoration={snake, amplitude=0.6mm}] (0.5,3.4) -- (9.0,3.9);
\node[popblue, font=\scriptsize] at (2.0,3.7) {cours d'eau};
% route principale
\draw[very thick, poporange] (0.6,1.0) -- (4.6,2.2) -- (9.2,1.6);
\node[poporange, font=\scriptsize] at (7.6,1.15) {route vers le port};
% rail interne
\draw[thick, popdark, dashed] (1.0,0.6) -- (5.0,1.4) -- (8.8,1.1);
\node[popdark, font=\scriptsize] at (3.2,0.85) {rail interne};
% villages riverains
\foreach \p/\n in {(-0.5,2.0)/Village A, (9.9,3.2)/Village B, (4.5,-0.4)/Village C}
  {\fill[poppurple] \p circle (2.2pt); \node[font=\scriptsize, poppurple] at \p [below=1pt] {\n};}
% nord
\draw[-{Stealth}, very thick, popdark] (10.6,5.2) -- (10.6,6.2) node[above]{N};
% echelle
\draw[thick] (0.5,-0.9) -- (2.5,-0.9) node[midway, below, font=\scriptsize]{5 km (indicatif)};
% zones vierges a completer
\node[font=\scriptsize\itshape, popmuted] at (4.7,4.6) {emprise de la concession (à compléter :};
\node[font=\scriptsize\itshape, popmuted] at (4.7,4.2) {cultures, usines, cités ouvrières)};
\end{tikzpicture}
\end{center}
\legende{Base de croquis vierge d'une concession type (inspirée de la CDC) : limites de la concession, axe routier et voie ferrée interne d'évacuation, cours d'eau, villages riverains. En TP, vous y placerez les aplats de cultures, l'usine, les cités ouvrières et vous construirez la légende.}
\end{popfigure}

\begin{aretenir}[Attendus au Bac pour un croquis commenté de plantation]
\begin{itemize}
\item \textbf{Titre précis} : « Organisation spatiale de la concession de \ldots{} (firme, région) ».
\item \textbf{Échelle} indicative et \textbf{orientation} (flèche du nord) obligatoires.
\item \textbf{Légende organisée} en rubriques : Agriculture / Industrie / Infrastructures / Habitat / Conflits fonciers.
\item \textbf{Commentaire structuré} en trois temps : \emph{localisation} (site, situation, accès au port), \emph{organisation} (hiérarchie siège--estates--usines, cité ouvrière, planteurs villageois), \emph{enjeux} (emploi, exportations, environnement, foncier).
\end{itemize}
\end{aretenir}

En définitive, la firme agro-industrielle camerounaise apparaît comme un \cle{système spatial complet} : un commandement urbain, des champs découpés en estates, des usines au cœur des cultures, une cité ouvrière qui fait vivre des milliers de familles, et des flux convergent vers les ports. Mais ce système performant est aussi un espace de tensions — foncières, sociales, environnementales — que l'État, les firmes et les communautés tentent aujourd'hui d'apaiser. C'est précisément cette double lecture, organisation et enjeux, que le croquis de synthèse du TP intégré (section 5) vous demandera de restituer.

\coursec{Le secteur industriel : Typologie, facteurs de localisation et pôles majeurs}

Après avoir étudié l'organisation spatiale d'une firme agro-industrielle comme la SOCAPALM ou la CDC, nous avons vu que la transformation des matières premières agricoles constitue déjà une forme d'industrie. Mais l'industrie camerounaise ne se réduit pas à l'agro-industrie : elle comprend aussi des usines lourdes (aluminium, raffinage, ciment) et des industries légères de consommation. Où se situent ces usines ? Pourquoi là et pas ailleurs ? C'est toute la question de la \cle{localisation industrielle}, qui permet de comprendre la carte des grands pôles industriels du Cameroun.

\courssub{La typologie de l'industrie camerounaise}

\begin{definition}[Industrie et pôle industriel]
L'\cle{industrie} est l'ensemble des activités économiques qui transforment des matières premières (agricoles, forestières, minérales, énergétiques) en biens de consommation ou en biens intermédiaires. Un \cle{pôle industriel} est une concentration d'unités de production sur un même espace, liées entre elles par des \cle{externalités d'agglomération} : main-d'œuvre disponible et formée, services partagés (banques, réparation, transport), infrastructures communes (port, énergie, routes).
\end{definition}

L'industrie camerounaise se classe en trois grandes familles.

\textbf{1. Les industries agro-alimentaires (IAA).} C'est la branche dominante : premier employeur industriel et premier chiffre d'affaires du secteur. Elle comprend :
\begin{itemize}
\item les \cle{brasseries} : SABC (Société Anonyme des Brasseries du Cameroun) et UCB (Union Camerounaise de Brasseries), implantées à Douala, Yaoundé, Bafoussam, Garoua et Sembe ;
\item les \cle{sucreries} : la SOSUCAM (Société Sucrière du Cameroun) à Nkoteng et Mbandjock, dans la région du Centre ;
\item les \cle{huileries} : SOCAPALM, SAFACAM et AZUR (savonnerie et raffinage d'huile à Edéa) ;
\item les conserveries, minoteries et aliments pour bétail.
\end{itemize}

\textbf{2. Les industries lourdes et chimiques.} Elles exigent beaucoup de capital et d'énergie :
\begin{itemize}
\item \cle{ALUCAM} (Compagnie Camerounaise de l'Aluminium) à Edéa : production d'aluminium à partir d'alumine importée ;
\item la \cle{SONARA} (Société Nationale de Raffinage) à Limbé : raffinage du pétrole brut ;
\item \cle{CIMENCAM} (Cimenteries du Cameroun) à Douala-Bonabéri et Nomayos (près de Yaoundé) : production de ciment ;
\item les industries chimiques : PPC (Produits Pharmaceutiques du Cameroun) à Edéa, SOCAVER (verrerie) à Douala.
\end{itemize}

\textbf{3. Les industries légères.} Elles produisent des biens de consommation courante : textile (CICAM à Douala et Garoua, entreprise en grande difficulté face à la concurrence des importations), chaussures, plastiques, et la filière bois (première transformation : sciages ; seconde transformation : contreplaqués, meubles), concentrée dans l'Est et le Sud.

\begin{exemplebox}[Quelques chiffres repères]
\begin{itemize}
\item ALUCAM consomme une part très importante de la production électrique nationale (ENEO, ex-AES-SONEL) : l'aluminium est une industrie « électrovore ».
\item La SONARA a une capacité nominale de raffinage d'environ 2,5 millions de tonnes de pétrole brut par an.
\item CIMENCAM produit environ 2,5 millions de tonnes de ciment par an.
\item Les industries agro-alimentaires représentent environ 45 \% de la valeur ajoutée industrielle (ordres de grandeur INS).
\end{itemize}
\end{exemplebox}

\courssub{Les facteurs de localisation des industries}

Pourquoi ALUCAM est-elle à Edéa et non à Maroua ? Pourquoi les brasseries sont-elles près des villes ? La réponse tient à la notion de \cle{facteur de localisation}.

\begin{propriete}[Les facteurs de localisation industrielle (A. Weber, F. Perroux)]
Selon la théorie d'Alfred Weber, une entreprise choisit son implantation en minimisant ses coûts, selon une hiérarchie :
\begin{enumerate}
\item les \cle{coûts de transport} : on se rapproche soit des matières premières (si elles sont lourdes ou périssables), soit du marché (si le produit fini est plus coûteux à transporter) ;
\item les \cle{coûts de main-d'œuvre} : recherche d'une main-d'œuvre abondante et peu coûteuse ;
\item les \cle{facteurs d'agglomération} (Perroux) : avantages liés à la proximité d'autres entreprises, des services et des infrastructures.
\end{enumerate}
Au Cameroun, pour l'industrie lourde, deux facteurs dominent : l'\cle{énergie hydroélectrique} et l'\cle{accès portuaire}.
\end{propriete}

Appliquons cette grille au Cameroun :

\begin{enumerate}
\item \textbf{L'énergie hydroélectrique.} Les barrages d'Edéa et de Song Loulou (sur la Sanaga), complétés par Memve'ele, fournissent une électricité abondante et relativement bon marché. C'est ce qui explique la présence d'ALUCAM et d'AZUR à Edéa.
\item \textbf{La position portuaire.} Les ports de Douala et de Kribi permettent d'importer les intrants (alumine pour ALUCAM, clinker pour le ciment) et d'exporter les produits. La SONARA à Limbé bénéficie de la façade maritime du Sud-Ouest, tandis que le pipeline Tchad-Cameroun, dont le terminal d'exportation est à Kribi, a renforcé l'attractivité pétrolière du littoral.
\item \textbf{Les matières premières.} Le bois attire les scieries dans l'Est et le Sud ; le coton attire les usines de la SODECOTON dans le Nord (Garoua, Maroua) ; le caoutchouc et le palmier à huile fixent les usines de transformation dans le Sud-Ouest et le Littoral.
\item \textbf{Le marché de consommation.} Douala et Yaoundé, grandes agglomérations, concentrent brasseries, cimenteries et agro-alimentaire léger : la bière et le ciment se transportent mal, on les produit près des consommateurs.
\item \textbf{Les transports.} Le chemin de fer Camrail (Douala–Yaoundé–Ngaoundéré), les routes bitumées et le pipeline Tchad-Cameroun structurent les axes industriels.
\end{enumerate}

\courssub{Les grands pôles industriels du Cameroun}

\begin{popfigure}
\begin{center}
\begin{tikzpicture}[scale=0.95, font=\small]
% Contour très simplifié du Cameroun
\draw[thick, popdark, fill=popcream] (2.2,7.6) -- (3.4,7.9) -- (4.6,7.2) -- (4.9,6.2) -- (4.4,5.2) -- (4.7,4.2) -- (4.3,3.2) -- (4.6,2.2) -- (3.9,1.0) -- (2.9,0.3) -- (1.9,0.9) -- (1.3,2.2) -- (0.9,3.4) -- (1.2,4.6) -- (1.6,5.6) -- (1.5,6.6) -- cycle;
% Nord
\draw[->, thick, popdark] (0.35,7.3) -- (0.35,7.9); \node[above] at (0.35,7.9) {\scriptsize N};
% Pôles
\node[circle, fill=popblue, minimum size=0.62cm, inner sep=0pt] (douala) at (1.35,3.35) {};
\node[circle, fill=poporange, minimum size=0.45cm, inner sep=0pt] (edea) at (2.15,3.15) {};
\node[circle, fill=popgreen, minimum size=0.5cm, inner sep=0pt] (yaounde) at (2.7,2.75) {};
\node[circle, fill=poppurple, minimum size=0.35cm, inner sep=0pt] (kribi) at (1.95,1.85) {};
\node[circle, fill=popgold, minimum size=0.35cm, inner sep=0pt] (garoua) at (3.7,5.6) {};
\node[circle, fill=popgold, minimum size=0.28cm, inner sep=0pt] (maroua) at (4.15,6.7) {};
\node[circle, fill=popblue!60, minimum size=0.3cm, inner sep=0pt] (limbe) at (0.95,2.75) {};
% Noms
\node[right, align=left] at (1.7,3.5) {\textbf{Douala} \scriptsize (1\textsuperscript{er} pôle, diversifié)};
\node[right] at (2.4,3.2) {\scriptsize \textbf{Edéa} (lourd : ALUCAM, AZUR)};
\node[right] at (2.95,2.75) {\scriptsize \textbf{Yaoundé} (consommation)};
\node[right] at (2.2,1.85) {\scriptsize \textbf{Kribi} (port en eau profonde)};
\node[right] at (3.9,5.6) {\scriptsize \textbf{Garoua} (coton)};
\node[right] at (4.35,6.7) {\scriptsize \textbf{Maroua}};
\node[left] at (0.85,2.75) {\scriptsize \textbf{Limbé} (SONARA)};
% Flux
\draw[->, very thick, poppink] (4.9,3.9) -- (2.1,2.1) node[midway, above right, sloped]{\scriptsize Pipeline Tchad--Cameroun};
\draw[->, thick, popgreen!70!black] (4.5,1.2) -- (1.7,3.1) node[midway, below right, sloped]{\scriptsize Bois (Est/Sud)};
\draw[->, thick, popgold!80!black] (3.7,5.5) -- (1.6,3.5) node[midway, above left, sloped]{\scriptsize Coton (Nord)};
% Chemin de fer
\draw[dashed, thick, popmuted] (1.35,3.35) -- (2.7,2.75) -- (3.4,4.4) -- (3.7,5.6) -- (3.9,6.4);
\node[popmuted] at (3.35,4.05) {\scriptsize Camrail};
\end{tikzpicture}
\end{center}
\legende{Croquis de localisation des pôles industriels du Cameroun (schématique). La taille des cercles est proportionnelle au poids industriel ; les couleurs indiquent la spécialisation : bleu = diversifié/portuaire, orange = industrie lourde, vert = consommation, violet = pôle émergent, doré = agro-ressources du Nord. Les flèches représentent les flux de matières premières.}
\end{popfigure}

\textbf{Le pôle Douala-Bonabéri} est le premier pôle industriel du pays : il concentre plus de la moitié de la production industrielle nationale. Diversifié (agro-alimentaire, chimie, mécanique, bois, BTP), il doit sa puissance à son port (premier port du pays), à son aéroport international et à sa position de nœud routier et ferroviaire. Mais il souffre de \cle{congestion} (embouteillages, port saturé), d'un foncier devenu très cher et de pollution.

\textbf{Le pôle d'Edéa} est spécialisé dans l'industrie lourde (ALUCAM, AZUR, PPC) grâce à l'énergie des barrages d'Edéa et de Song Loulou, à son port fluvial sur la Sanaga et à sa position routière entre Douala et Kribi. C'est une zone industrielle planifiée dès l'époque coloniale.

\textbf{Le pôle de Yaoundé}, deuxième marché du pays, accueille des industries de consommation : brasseries, ciment (Nomayos), plastiques, produits pharmaceutiques (CINPHARM). Ses zones industrielles (Mvan, Nsam, Ekounou) servent surtout le marché intérieur : la capitale est peu exportatrice.

\textbf{Les pôles émergents et spécialisés.} Kribi dispose depuis 2013 d'un port en eau profonde et d'une zone franche industrielle, avec des perspectives gazières (GNL) et minières (projet de fer de Mbalam). Dans le Nord, Garoua, Kousseri et Maroua s'appuient sur le coton (SODECOTON), les huileries et les abattoirs ; mais ces pôles souffrent de l'\cle{enclavement} et d'une énergie thermique coûteuse.

\begin{attention}[Ne pas « doualiser » toute l'industrie camerounaise]
Erreur fréquente à l'examen : limiter l'industrie camerounaise à Douala. Douala est certes le premier pôle, mais Edéa (industrie lourde), Yaoundé (industries de consommation), le Nord (agro-ressources : coton, huiles) et Kribi (pôle d'avenir) sont indispensables dans toute analyse complète. Attention aussi à ne pas situer la SOSUCAM dans le Nord : ses sucreries sont à Nkoteng et Mbandjock, dans la région du Centre. Pensez enfin à citer les limites de chaque pôle (congestion à Douala, enclavement du Nord).
\end{attention}

\courssub{Lire et construire des graphiques sur la structure industrielle}

\begin{methode}[Construire un diagramme de la structure industrielle par branche]
À partir d'un tableau statistique (INS, MINMIDT) :
\begin{enumerate}
\item Identifier le caractère étudié (valeur ajoutée, effectifs, chiffre d'affaires) et les branches retenues.
\item Vérifier que les données sont en pourcentages ; sinon, calculer la part de chaque branche : part = (valeur de la branche / total) $\times$ 100.
\item Choisir le graphique adapté : barres empilées ou camembert pour une répartition en parts ; barres simples pour comparer des régions.
\item Tracer les axes ou les secteurs avec exactitude, ajouter titre, légende, unité et source.
\item Rédiger une phrase de lecture : repérer la branche dominante, puis les contrastes.
\end{enumerate}
\end{methode}

\begin{popfigure}
\begin{center}
\begin{tikzpicture}
\begin{axis}[
    ybar, bar width=1.1cm,
    width=0.85\linewidth, height=6.2cm,
    ymin=0, ymax=55,
    ylabel={Part de la valeur ajoutée industrielle (\%)},
    symbolic x coords={Agro-alimentaire, Bois, Chimie/Pétrole, BTP/Matériaux, Autres},
    xtick=data, x tick label style={rotate=15, anchor=east, font=\scriptsize},
    nodes near coords, nodes near coords style={font=\scriptsize},
    axis lines=left, enlarge x limits=0.12
]
\addplot[fill=popblue, draw=popdark] coordinates {(Agro-alimentaire,45) (Bois,15) (Chimie/Pétrole,20) (BTP/Matériaux,10) (Autres,10)};
\end{axis}
\end{tikzpicture}
\end{center}
\legende{Répartition approximative de la valeur ajoutée industrielle par branche au Cameroun (données INS, ordres de grandeur). Les industries agro-alimentaires dominent nettement la structure industrielle nationale.}
\end{popfigure}

\begin{exoresolu}[Calculer et interpréter une part de marché industriel]
Énoncé : la valeur ajoutée industrielle totale du Cameroun est estimée à 2 000 milliards de FCFA. Les industries agro-alimentaires en représentent 900 milliards. (1) Calculez la part des IAA en pourcentage. (2) Comparez avec la branche bois (300 milliards) et concluez en deux phrases.
\tcblower
Solution : (1) Part des IAA $= \dfrac{900}{2000} \times 100 = 45\,\%$. (2) Part du bois $= \dfrac{300}{2000} \times 100 = 15\,\%$. Conclusion : les industries agro-alimentaires représentent près de la moitié de la valeur ajoutée industrielle, soit trois fois plus que la filière bois ($45 / 15 = 3$). Cela confirme que l'industrie camerounaise repose d'abord sur la transformation des produits agricoles, ce qui traduit la faible diversification du tissu industriel.
\end{exoresolu}

\begin{savaistu}[La Zone Franche de Kribi]
Créée par décret en 2013, la Zone Franche de Kribi offre aux entreprises exportatrices des avantages exceptionnels : exonération d'impôt sur les sociétés pendant 10 ans, exonération de TVA et de droits de douane, liberté de change. L'objectif est d'attirer les industries d'exportation autour du port en eau profonde et de faire de Kribi le futur grand pôle industriel du Cameroun, à l'image des zones franches asiatiques.
\end{savaistu}

\begin{aretenir}[L'essentiel]
L'industrie camerounaise est dominée par les industries agro-alimentaires (environ 45 \% de la valeur ajoutée), suivies de la chimie-pétrole, du bois et des matériaux. Sa localisation obéit à des facteurs précis : énergie hydroélectrique (Edéa), accès portuaire (Douala, Kribi, Limbé), matières premières (coton au Nord, bois à l'Est) et marchés (Douala, Yaoundé). Quatre pôles structurent l'espace industriel : Douala-Bonabéri (dominant et diversifié), Edéa (lourd), Yaoundé (consommation) et les pôles émergents (Kribi) ou spécialisés (Nord).
\end{aretenir}

\coursec{TP Intégré : Croquis de synthèse « Les pôles de l'économie moderne au Cameroun »}

Après avoir analysé la typologie industrielle et les facteurs de localisation des usines, nous disposons désormais de tous les éléments pour réaliser une synthèse graphique. Ce travail de croquis est l'exercice roi de l'épreuve du Baccalauréat technique : il exige de passer de l'analyse textuelle à la modélisation spatiale, puis à l'argumentation écrite. C'est la preuve que vous maîtrisez la géographie de l'économie moderne camerounaise.

\courssub{Consigne et démarche d'investigation}

\begin{methode}[Étapes de réalisation d'un croquis géographique de synthèse]
\begin{enumerate}
    \item \textbf{Analyse du sujet et des documents.} Lire la consigne, identifier les thèmes (agriculture de plantation, industrie, infrastructures, enjeux), l'échelle (nationale) et les documents fournis (fond de carte, tableaux statistiques, textes).
    \item \textbf{Construction de la légende organisée.} Regrouper les informations par \textbf{thèmes} (colonnes), définir une \textbf{hiérarchie} visuelle (grands pôles $>$ pôles secondaires $>$ axes) et choisir les \textbf{variables visuelles} adaptées : l'aplat de couleur pour les plantations, le point proportionnel pour l'industrie, le trait pour les flux et les axes.
    \item \textbf{Figuration soignée sur la carte de base.} Travailler d'abord au crayon à papier. Respecter l'ordre d'imbrication : fond (aplats des plantations) $\rightarrow$ réseau (axes, pipeline) $\rightarrow$ points (usines, ports, barrages) $\rightarrow$ surfaces ponctuelles (conflits, zone franche) $\rightarrow$ noms (villes, régions).
    \item \textbf{Finalisation et esthétique.} Encrer au feutre fin, colorier aux crayons de couleur (aplats légers), rédiger la légende \textbf{au feutre noir} sur le côté droit de la feuille, avec des cases parfaitement alignées.
    \item \textbf{Rédaction du commentaire structuré.} Introduction (définition, localisation, problématique, annonce du plan) $\rightarrow$ développement en deux parties équilibrées $\rightarrow$ conclusion (bilan synthétique, ouverture sur l'émergence).
\end{enumerate}
\end{methode}

\courssub{Étape 1 : La légende type — matrice sémiologique complète}

La légende n'est pas une simple liste : c'est une \textbf{grammaire visuelle}. Elle doit être construite \emph{avant} de toucher à la carte. Voici la matrice de référence pour ce croquis, respectant les variables visuelles de la sémiologie graphique (forme, taille, couleur, orientation, texture, valeur).

\begin{definition}[Légende organisée et hiérarchisée]
Une légende géographique rigoureuse repose sur trois piliers :
\begin{enumerate}
    \item \textbf{Regroupement thématique :} les signes sont classés en colonnes logiques (ici : agriculture / industrie / infrastructures et enjeux).
    \item \textbf{Ordre logique :} à l'intérieur d'une colonne, on suit l'ordre : surfaces (aplats) $\rightarrow$ traits (linéaires) $\rightarrow$ points (ponctuels) $\rightarrow$ noms (textes).
    \item \textbf{Hiérarchie visuelle :} la taille du symbole (diamètre, épaisseur) est proportionnelle à l'importance quantitative (hectares, effectifs, tonnage) ; la couleur code la nature qualitative (spéculation, secteur industriel).
\end{enumerate}
\end{definition}

\begin{exemplebox}[Légende type normalisée pour le croquis de synthèse]
\begin{center}
\begin{tabularx}{\linewidth}{|X|X|X|}\hline
\rowcolor{popblueL} \textbf{AGRICULTURE DE PLANTATION} & \textbf{INDUSTRIE} & \textbf{INFRASTRUCTURES ET ENJEUX} \\ \hline
Aplat jaune : banane (CDC, PHP) & Carré rouge : industrie lourde, chimie, métallurgie (Alucam, SONARA) & Trait plein épais : route nationale (N1, N3, N10) \\ \hline
Aplat vert foncé : palmier à huile (SOCAPALM, SAFACAM) & Cercle bleu : agro-alimentaire (brasseries, SOSUCAM, huileries) & Trait tireté : chemin de fer (Douala--Ngaoundéré) \\ \hline
Aplat vert clair : hévéa (HEVECAM, SUDCAM) & Triangle vert : bois (scieries) & Trait pointillé rouge : pipeline Tchad--Cameroun (Doba--Kribi) \\ \hline
Aplat orange : canne à sucre (SOSUCAM) & Losange gris : matériaux de construction (Cimencam) & Trait tireté bleu : ligne haute tension (Song Loulou--Douala) \\ \hline
Aplat marron : thé (CDC Tole, UNVDA) & Grand symbole : grand pôle ($>$ 500 emplois) ; petit symbole : unité secondaire & Ancre : port (Douala, Kribi, Limbé, Garoua) ; triangle bleu : barrage (Edéa, Song Loulou, Lagdo, Memve'ele) \\ \hline
Grand aplat : grand pôle ($>$ \num{20000} ha) ; petit aplat : pôle secondaire (\num{5000}--\num{20000} ha) & & Hachures rouges : conflits fonciers ; cadre pointillé violet : zone franche de Kribi ; étoile : ville primaire (Douala, Yaoundé) \\ \hline
\end{tabularx}
\end{center}
Ordre impératif dans chaque colonne : surfaces $\rightarrow$ traits $\rightarrow$ points $\rightarrow$ noms.
\end{exemplebox}

\begin{attention}[Respecter la sémiologie graphique — pièges fréquents au Bac]
\begin{itemize}
    \item \textbf{Variable taille :} réservée aux données \textbf{quantitatives} (effectifs, superficie, tonnage). Il est interdit de faire varier la taille pour une simple différence de nature (banane contre palme : c'est la couleur qui change, pas la taille).
    \item \textbf{Variable forme :} réservée aux données \textbf{qualitatives} (secteur industriel : carré pour l'industrie lourde, cercle pour l'agro-alimentaire, triangle pour le bois). Ne pas utiliser de formes complexes impossibles à reproduire proprement.
    \item \textbf{Variable couleur :} elle code la \textbf{nature} (spéculation agricole, type d'industrie). Utiliser des teintes bien discriminées.
    \item \textbf{Variable texture (hachures) :} réservée aux \textbf{enjeux et risques} (conflits fonciers) ou aux zones spécifiques (zone franche). Ne jamais hachurer les plantations, qui restent en aplats pleins.
    \item \textbf{Propreté :} un croquis sale, baveux, avec une légende au crayon ou mal alignée, est lourdement pénalisé par les points de présentation.
\end{itemize}
\end{attention}

\courssub{Étape 2 : Figuration des plantations — l'emprise agro-industrielle}

L'agriculture de plantation structure l'espace littoral et méridional. On utilise des \textbf{aplats de couleur pleins} (crayon de couleur, pression légère) sur les emprises réelles. La hiérarchie se fait par la \textbf{surface du polygone} (grand pôle contre pôle secondaire) et la \textbf{couleur} (spéculation).

\begin{exemplebox}[Emprises majeures à figurer (ordres de grandeur)]
\begin{itemize}
    \item \textbf{CDC (Cameroon Development Corporation) :} région du Sud-Ouest (Tiko, Muyuka, Ekona, Idenau). Banane (jaune), palmier (vert foncé), hévéa (vert clair), thé (marron, à Tole). Plusieurs dizaines de milliers d'hectares $\rightarrow$ \textbf{grand pôle} (gros aplat).
    \item \textbf{PHP (Plantations du Haut-Penja) :} Moungo (Penja, Njombé). Banane (jaune). Quelques milliers d'hectares $\rightarrow$ \textbf{pôle secondaire}.
    \item \textbf{SOCAPALM :} plusieurs sites (Dibombari, Mbongo, Kienké, Eséka, Makené). Palmier à huile (vert foncé). Plusieurs dizaines de milliers d'hectares au total $\rightarrow$ \textbf{grands pôles} (Dibombari, Kienké) et pôles secondaires.
    \item \textbf{HEVECAM et SUDCAM :} région du Sud (vallée du Ntem, environs de Kribi et de Meyomessala). Hévéa (vert clair). Grandes concessions $\rightarrow$ \textbf{grands pôles}.
    \item \textbf{SOSUCAM :} Nkoteng et Mbandjock (Sanaga-Maritime). Canne à sucre (orange). Environ \num{10000} ha $\rightarrow$ \textbf{grand pôle}.
    \item \textbf{UNVDA (Upper Noun Valley Development Authority) :} plaine de Bamendjin (Ouest). Riz et maraîchage. Quelques milliers d'hectares $\rightarrow$ \textbf{pôle secondaire}, seul pôle de développement agricole moderne notable hors du littoral.
\end{itemize}
\end{exemplebox}

\begin{attention}[Exactitude des localisations]
Ne pas confondre les localisations : la CDC est dans le \textbf{Sud-Ouest} (et non dans le Littoral), la PHP dans le \textbf{Moungo} (Littoral), SOSUCAM dans la \textbf{Sanaga-Maritime} (Littoral), HEVECAM et SUDCAM dans le \textbf{Sud}, l'UNVDA dans l'\textbf{Ouest}. Une plantation placée dans la mauvaise région est une erreur de figuration sanctionnée.
\end{attention}

\courssub{Étape 3 : Figuration de l'industrie — symboles ponctuels proportionnels}

L'industrie occupe peu d'espace au sol mais pèse lourd en valeur. On utilise des \textbf{points proportionnels} (cercles, carrés, triangles, losanges) dont le \textbf{diamètre} croît avec l'effectif ou le chiffre d'affaires.

\begin{exoresolu}[Calibrage des symboles industriels — méthode Bac]
\textbf{Données (ordres de grandeur) :} Alucam (Edéa) : environ 800 emplois. Brasseries (Douala, Yaoundé, Garoua, Bafoussam) : 500 à \num{1000} emplois par site. Cimencam (Figuil, Douala, Yaoundé) : 300 à 500 emplois. Scieries (Douala, Bertoua, Ngaoundéré) : 100 à 300 emplois.
\textbf{Choix d'échelle :} 1 mm de diamètre pour 50 emplois.
\textbf{Calculs :}
\begin{align*}
\text{Alucam (grand pôle)} &: \varnothing = \frac{800}{50} = 16 \text{ mm} \rightarrow \textbf{carré rouge de 16 mm} \\
\text{Brasserie de Douala (grand pôle)} &: \varnothing = \frac{900}{50} = 18 \text{ mm} \rightarrow \textbf{cercle bleu de 18 mm} \\
\text{Cimencam Figuil (pôle moyen)} &: \varnothing = \frac{400}{50} = 8 \text{ mm} \rightarrow \textbf{losange gris de 8 mm} \\
\text{Scierie de Bertoua (unité secondaire)} &: \varnothing = \frac{150}{50} = 3 \text{ mm} \rightarrow \textbf{triangle vert de 3 mm}
\end{align*}
\textbf{Astuce :} dessiner les grands symboles au compas sur papier calque, reporter au crayon, puis encrer. Ne jamais superposer deux symboles au même endroit : décaler légèrement avec une flèche de rappel si deux usines sont dans la même ville.
\tcblower
\textbf{Résultat attendu sur la carte :} concentration massive de symboles (rouges, bleus) sur l'axe \textbf{Douala--Edéa--Kribi} (corridor industriel). Symboles bleus (agro-alimentaire) à Nkoteng, Mbandjock, Dibombari, Kienké. Symboles verts (bois) à Douala, Bertoua, Ngaoundéré. Symboles gris (matériaux) à Figuil, Douala, Yaoundé. La lecture de la carte montre immédiatement la polarisation littorale de l'industrie camerounaise.
\end{exoresolu}

\courssub{Étape 4 : Infrastructures structurantes — le squelette des flux}

Ces éléments sont des \textbf{traits linéaires}. Leur épaisseur et leur style (plein, pointillé, tirets) codent la nature et l'importance de l'infrastructure.

\begin{exemplebox}[Repères pour dimensionner les traits (ordres de grandeur)]
\begin{center}
\begin{tabularx}{\linewidth}{|l|X|X|}\hline
\rowcolor{popgreenL} \textbf{Infrastructure} & \textbf{Caractéristiques} & \textbf{Figuration (épaisseur, style)} \\ \hline
Port de Douala & Premier port national, l'essentiel du trafic maritime & Trait très épais (2 mm) + ancre noire \\ \hline
Port en eau profonde de Kribi & Port moderne en expansion, terminal à conteneurs & Trait très épais (2 mm) + ancre + cadre pointillé violet (zone franche) \\ \hline
Port de Limbé & Port pétrolier (SONARA) & Trait moyen (1 mm) + ancre \\ \hline
Port de Garoua & Port fluvial sur la Bénoué, saisonnier & Trait fin (0,5 mm) pointillé bleu \\ \hline
Barrage d'Edéa (Sanaga) & Environ 263 MW & Triangle bleu + ligne HT (tirets bleus) vers Douala \\ \hline
Barrage de Song Loulou & Environ 384 MW & Triangle bleu + ligne HT \\ \hline
Barrage de Lagdo (Bénoué) & Environ 72 MW & Triangle bleu \\ \hline
Barrage de Memve'ele (Ntem) & Environ 211 MW & Triangle bleu + ligne HT vers Yaoundé et le Sud \\ \hline
Pipeline Tchad--Cameroun & Environ \num{1070} km (Doba--Kribi) & Trait pointillé épais (1,5 mm) rouge, flèches orientées Tchad $\rightarrow$ Kribi \\ \hline
Chemin de fer (Camrail) & Environ \num{1000} km exploités (Douala--Ngaoundéré et embranchements) & Trait tireté épais (1 mm) noir \\ \hline
Route N1 (Douala--Yaoundé--Ngaoundéré--Garoua) & Axe structurant sud--nord & Trait plein épais (1,5 mm) noir \\ \hline
Route N3 (Edéa--Yaoundé--Bertoua--Garoua-Boulaï) & Axe vers l'Est & Trait plein moyen (1 mm) noir \\ \hline
Route N10 (Bafoussam--Bamenda--Ngaoundéré) & Axe ouest--nord & Trait plein moyen (1 mm) noir \\ \hline
\end{tabularx}
\end{center}
\end{exemplebox}

\courssub{Étape 5 : Figurer les enjeux — conflits, zone franche, métropoles}

Ces éléments se superposent aux précédents. Ils se figurent en \textbf{dernier}, au-dessus des aplats et des traits.

\begin{itemize}
    \item \textbf{Conflits fonciers (hachures rouges) :} zones de tension entre plantations, populations riveraines et aires protégées.
        \begin{itemize}
            \item \textbf{Sud-Ouest (CDC) :} conflits fonciers historiques avec les communautés riveraines, aggravés par la crise anglophone (abandons de plantations, destructions).
            \item \textbf{Sud (HEVECAM, SUDCAM) :} revendications des populations locales et autochtones (Bagyeli), déforestation, accusations d'accaparement des terres.
            \item \textbf{Moungo (PHP, SOCAPALM) :} forte pression foncière et spéculation autour des plantations.
        \end{itemize}
    \item \textbf{Zone franche de Kribi (cadre pointillé violet) :} périmètre du port autonome et de la zone industrielle bénéficiant d'un régime douanier et fiscal privilégié. Figurer un grand rectangle pointillé autour de Kribi et Lolabé.
    \item \textbf{Villes primaires (étoiles rouges) :} \textbf{Douala} (capitale économique, premier port, premier pôle industriel) et \textbf{Yaoundé} (capitale politique, deuxième pôle industriel et de services). Les nommer en \textbf{majuscules grasses}.
\end{itemize}

\begin{savaistu}[Le corridor Douala--Kribi--Edéa : cœur battant de l'économie moderne]
Ce triangle littoral concentre l'essentiel de la valeur ajoutée industrielle du pays (Alucam, SONARA à Limbé toute proche, Cimencam, brasseries, nouvelles usines de Kribi) et la quasi-totalité du trafic portuaire national (Douala et Kribi). C'est la seule région du Cameroun où se combinent : énergie hydroélectrique abondante (Edéa, Song Loulou), ports (dont un port en eau profonde), terminal du pipeline, raffinerie à proximité, réseau ferré, autoroute en construction vers Yaoundé et main-d'œuvre qualifiée. C'est le moteur de la trajectoire d'émergence à l'horizon 2035. Ses vulnérabilités (inondations à Douala, érosion côtière à Kribi, saturation routière) constituent le risque majeur du système productif national.
\end{savaistu}

\courssub{Étape 6 : Commentaire type Baccalauréat — modèle de rédaction}

Le commentaire ne décrit pas la carte (la légende s'en charge) : il \textbf{l'explique} et la \textbf{problématise}.

\begin{exoresolu}[Commentaire type — sujet : « Les pôles de l'économie moderne au Cameroun »]
\textbf{Introduction}
\begin{itemize}
    \item \textbf{Définition :} l'économie moderne au Cameroun désigne l'ensemble des activités productives (agro-industrie, industrie manufacturière) organisées en entreprises capitalistes employant des salariés, tournées vers le marché national et international, et structurées par l'État depuis l'indépendance.
    \item \textbf{Localisation :} elle s'inscrit dans un espace national d'environ \num{475000} km$^2$, contrasté entre une façade atlantique ouverte et un arrière-pays (hinterland) enclavé.
    \item \textbf{Problématique :} si l'économie moderne repose sur deux piliers — l'agriculture de plantation (héritage colonial) et l'industrie naissante —, son organisation spatiale révèle un \textbf{dualisme territorial} marqué : une façade littorale connectée et productive face à un intérieur marginalisé. Comment ce dualisme est-il organisé et quels enjeux pose-t-il pour l'émergence ?
    \item \textbf{Annonce du plan :} nous analyserons d'abord la prédominance côtière de l'agro-industrie exportatrice (I), puis la polarisation industrielle sur le corridor Douala--Kribi--Edéa (II).
\end{itemize}

\textbf{Développement}

\textbf{I. L'agriculture de plantation : une rente foncière littorale et méridionale}
\begin{itemize}
    \item \textbf{A. Une spécialisation exportatrice pérenne.} Le croquis montre une concentration d'aplats (jaune banane, vert palme et hévéa, orange sucre, marron thé) sur la façade atlantique (Sud-Ouest, Littoral, Sud). Les grandes firmes (CDC, PHP, SOCAPALM, HEVECAM, SOSUCAM) occupent des sols favorables et bénéficient d'un accès portuaire direct (Douala, Kribi, Limbé). Le Cameroun est le premier exportateur africain de banane dessert et un producteur majeur d'huile de palme et de caoutchouc.
    \item \textbf{B. Un dualisme foncier et social.} La légende signale des hachures rouges (conflits). L'emprise des plantations s'est faite au détriment des terres coutumières des communautés. La crise anglophone (Nord-Ouest et Sud-Ouest) et les revendications des populations du Sud menacent la continuité de la production. L'UNVDA (Ouest) reste l'unique pôle intérieur significatif, ce qui illustre la difficulté d'étendre ce modèle vers l'intérieur.
\end{itemize}

\textbf{II. L'industrie : une polarisation extrême sur le corridor Douala--Kribi--Edéa}
\begin{itemize}
    \item \textbf{A. Le poids de l'agro-alimentaire et de l'industrie lourde.} Les symboles proportionnels (cercles bleus, carrés rouges) s'empilent sur l'axe Douala--Edéa--Kribi. Alucam (aluminium, alimentée par le barrage d'Edéa) et la SONARA (raffinage, Limbé) sont les deux locomotives industrielles. L'agro-alimentaire (brasseries, huileries, sucreries) transforme localement les productions des plantations : c'est l'intégration verticale (SOSUCAM, SOCAPALM).
    \item \textbf{B. Les infrastructures comme facteur de verrouillage.} Le trait pointillé rouge (pipeline d'environ \num{1070} km) et les lignes haute tension convergent vers Kribi et Douala. Le chemin de fer (Douala--Ngaoundéré) et la N1 sont les principales artères d'évacuation. L'enclavement de l'Est (bois, triangles verts à Bertoua et Ngaoundéré) et du Nord (coton, losanges gris à Figuil) limite leur industrialisation : ils exportent brut.
\end{itemize}

\textbf{Conclusion}
\begin{itemize}
    \item \textbf{Bilan :} l'économie moderne camerounaise est une économie polarisée par le littoral, spécialisée dans les matières premières et l'agro-alimentaire, dépendante d'infrastructures vieillissantes (sauf Kribi) et génératrice de conflits fonciers.
    \item \textbf{Ouverture :} l'émergence à l'horizon 2035 exige de \textbf{mailler le territoire} : achever l'autoroute Yaoundé--Douala, moderniser le rail, développer l'énergie hydroélectrique (Nachtigal, Memve'ele) pour attirer l'industrie hors du corridor littoral. Le croquis de demain devra montrer des symboles industriels essaimant le long de la N10, de la N3 et de l'axe trans-camerounais.
\end{itemize}
\tcblower
\textbf{Analyse de la copie (grille Bac) :}
\begin{itemize}
    \item \textit{Introduction :} définition précise, problématique claire (dualisme), plan annoncé.
    \item \textit{Développement :} deux parties équilibrées, arguments chiffrés (pipeline de \num{1070} km, puissances des barrages), référence constante aux symboles de la carte (« le croquis montre... », « les hachures rouges signalent... »).
    \item \textit{Conclusion :} bilan synthétique et ouverture prospective (maillage du territoire, énergie).
    \item \textit{Langage :} vocabulaire géographique approprié (hinterland, intégration verticale, polarisation, maillage).
\end{itemize}
\end{exoresolu}

\courssub{Support d'examen : fond de carte vierge (TikZ)}

Ce fond de carte est fourni « nu » (contour simplifié, hydrographie, villes, axes, limites régionales indicatives). Il sert de base à l'exercice : l'élève doit le compléter intégralement.

\begin{popfigure}
\begin{tikzpicture}[scale=0.75, every node/.style={font=\tiny}]
% -- CONTOUR CAMEROUN (polygone très simplifié) --
\draw[popline, thick, fill=popcream!40]
(0,0) -- (1,0.5) -- (2,1) -- (3.5,1.2) -- (5,0.8) -- (6.5,0.5) -- (7.5,0)
-- (8.5,-0.5) -- (9.5,-1) -- (10,-1.5) -- (10.5,-2) -- (10,-2.5) -- (9,-3) -- (8,-3.2) -- (7,-3)
-- (6,-2.8) -- (5,-2.5) -- (4,-2.2) -- (3,-2) -- (2,-1.8) -- (1,-1.5) -- (0.5,-1) -- (0,-0.5) -- cycle;

% -- LIMITES RÉGIONALES (indicatives, traits fins pointillés) --
\draw[popmuted, densely dashed, thin] (2,1) -- (4,0.5) -- (5,-0.5) -- (4,-1.5) -- (3,-1.8);
\draw[popmuted, densely dashed, thin] (5,0.8) -- (6.5,0) -- (7,-1) -- (6,-1.8);
\draw[popmuted, densely dashed, thin] (6.5,0.5) -- (8,0) -- (8.5,-1);
\draw[popmuted, densely dashed, thin] (4,-2.2) -- (6,-2.5) -- (7,-2.8);
\draw[popmuted, densely dashed, thin] (2,-1.8) -- (3.5,-2.5) -- (5,-2.5);

% -- HYDROGRAPHIE --
\draw[popblue, thick] (2.5,-0.5) .. controls (3.5,-1) and (4.5,-1.2) .. (5.5,-1.5) node[right, pos=0.8, popblue] {Sanaga};
\draw[popblue, thick] (1.5,0.5) .. controls (2.5,0) and (3.5,-0.5) .. (4.5,-1) node[right, pos=0.7, popblue] {Nyong};
\draw[popblue, thick] (4.5,0.5) -- (5,0) node[above, pos=0.5, popblue] {Wouri};
\draw[popblue, thick] (8,-1.5) -- (9,-2.5) node[right, pos=0.8, popblue] {Bénoué};
\draw[popblue, thick] (9.5,-2) -- (10,-2.8) node[right, pos=0.8, popblue] {Logone};
\draw[popblue, fill=popblue!20] (8.5,-0.5) ellipse (0.3 and 0.2);
\node[popblue] at (8.5,-0.1) {Lac Tchad};

% -- VILLES PRIMAIRES --
\fill[popdark] (5.2,0.1) circle (2.5pt);
\node[anchor=south west, font=\tiny\bfseries, popdark] at (5.3,0.2) {DOUALA};
\fill[popdark] (4.0,-1.2) circle (2.5pt);
\node[anchor=south west, font=\tiny\bfseries, popdark] at (4.1,-1.1) {YAOUNDÉ};

% -- VILLES SECONDAIRES --
\foreach \x/\y/\nom in {
    2.5/0.8/Bamenda,
    3.5/-0.5/Bafoussam,
    4.8/0.5/Edéa,
    5.5/-0.5/Kribi,
    3.0/1.2/Buea,
    7.5/-1.0/Bertoua,
    6.5/-2.0/Ngaoundéré,
    9.0/-2.5/Garoua
} {
    \fill[popdark] (\x,\y) circle (1.5pt);
    \node[anchor=south west] at (\x+0.1,\y+0.1) {\nom};
}

% -- AXES DE TRANSPORT (base) --
\draw[popmuted, densely dashed, thick] (5.2,0.1) -- (4.0,-1.2) -- (6.5,-2.0);
\draw[popmuted, densely dashed, thick] (5.2,0.1) -- (3.0,1.2);
\draw[popline, thin] (5.2,0.1) -- (4.0,-1.2) -- (6.5,-2.0) -- (9.0,-2.5);
\draw[popline, thin] (4.8,0.5) -- (4.0,-1.2) -- (7.5,-1.0);
\draw[popline, thin] (3.5,-0.5) -- (2.5,0.8) -- (6.5,-2.0);
\draw[poppink, densely dotted, thick] (9.5,-1.0) -- (8.0,-0.5) -- (6.0,-0.2) -- (5.5,-0.5);

% -- FLÈCHE NORD --
\node[anchor=south, popdark] at (10.5,1.5) {\textbf{N}};
\draw[->, thick, popdark] (10.5,1.2) -- (10.5,0.5);

% -- ÉCHELLE GRAPHIQUE INDICATIVE --
\draw[popdark, thick] (0.5,-3.5) -- (2.5,-3.5) node[midway, below] {200 km (indicatif)};
\draw[popdark, thick] (0.5,-3.55) -- (0.5,-3.45);
\draw[popdark, thick] (2.5,-3.55) -- (2.5,-3.45);
\end{tikzpicture}
\legende{Fond de carte vierge du Cameroun (contour simplifié, limites régionales indicatives, hydrographie, villes principales, axes ferroviaires, routes N1/N3/N10, pipeline). Croquis schématique sans précision cartographique : à compléter par l'élève selon la légende type.}
\end{popfigure}

\courssub{Critères d'évaluation — grille de notation Bac (20 points)}

\begin{aretenir}[Critères d'évaluation du croquis au Baccalauréat technique]
\begin{center}
\begin{tabularx}{\linewidth}{|l|X|X|}\hline
\rowcolor{popblueL} \textbf{Critère} & \textbf{Barème} & \textbf{Indicateurs de réussite} \\ \hline
\textbf{1. Présentation, propreté} & \textbf{4 pts} & Carte propre, sans ratures ; coloriage régulier (aplats légers) ; écriture lisible (noms horizontaux) ; légende au feutre noir, encadrée, cases alignées ; titre de la carte centré en haut. \\ \hline
\textbf{2. Légende} & \textbf{6 pts} & Organisée (trois colonnes thématiques), hiérarchisée (surfaces $\rightarrow$ traits $\rightarrow$ points), sémiologie respectée (forme, taille, couleur, texture adaptées aux données), exhaustive, titre de légende présent. \\ \hline
\textbf{3. Figuration (carte)} & \textbf{6 pts} & Exactitude des localisations (villes, plantations, usines, barrages, pipeline, rail) ; proportionnalité des symboles industriels ; superposition logique (fond $\rightarrow$ réseau $\rightarrow$ points $\rightarrow$ enjeux) ; noms reportés. \\ \hline
\textbf{4. Commentaire} & \textbf{4 pts} & Structure (introduction, développement, conclusion) ; problématique claire ; arguments appuyés sur la carte ; vocabulaire géographique précis ; chiffres-clés intégrés ; ouverture en conclusion. \\ \hline
\end{tabularx}
\end{center}
\textbf{Conseil stratégique :} ne sacrifiez jamais la légende ni le commentaire pour « finir le coloriage ». Une légende parfaite, un commentaire solide et une carte figée à 80\% rapportent plus qu'une carte entièrement coloriée avec une légende bâclée et sans commentaire.
\end{aretenir}

\courssub{Exercice d'entraînement — sujet zéro}

\begin{exercice}[Croquis noté — durée 1 h 30]
\textbf{Sujet :} à l'aide du fond de carte ci-joint et de vos connaissances, réaliser un croquis synthétique intitulé : \textbf{« L'organisation spatiale de l'économie moderne au Cameroun »}.
\textbf{Consignes :}
\begin{enumerate}
    \item Construire une légende organisée en trois colonnes (agriculture de plantation / industrie / infrastructures et enjeux) sur le côté droit de la feuille.
    \item Figurer sur la carte : les principales plantations (CDC, PHP, SOCAPALM, HEVECAM, SOSUCAM, UNVDA) par aplats de couleur ; les principaux pôles industriels (Alucam, SONARA, brasseries, Cimencam, scieries, huileries et sucreries) par symboles proportionnels codés par secteur ; les infrastructures structurantes (ports, barrages, pipeline, rail, routes N1/N3/N10) ; les enjeux (conflits fonciers du Sud-Ouest et du Sud, zone franche de Kribi, villes primaires).
    \item Rédiger un commentaire structuré (introduction, développement en deux parties, conclusion) d'environ 300 mots.
\end{enumerate}
\end{exercice}

\begin{corrige}[Exercice n°1 — éléments attendus pour la correction]
\textbf{Légende :} vérifier la présence des trois colonnes, l'ordre surfaces/traits/points, le codage couleur (banane = jaune, palmier = vert foncé, hévéa = vert clair, sucre = orange, thé = marron ; industrie lourde = carré rouge, agro-alimentaire = cercle bleu, bois = triangle vert, matériaux = losange gris), l'épaisseur des traits (pipeline, rail, routes), les hachures des conflits, le cadre de la zone franche.
\textbf{Figuration :}
\begin{itemize}
    \item \textit{Plantations :} gros aplats jaune, vert foncé, vert clair et marron dans le Sud-Ouest (CDC) ; jaune dans le Moungo (PHP) ; vert foncé dans le Littoral et le Sud (SOCAPALM) ; vert clair dans le Sud (HEVECAM, SUDCAM) ; orange dans la Sanaga-Maritime (SOSUCAM) ; petit aplat dans l'Ouest (UNVDA).
    \item \textit{Industrie :} gros carrés rouges (Alucam à Edéa, SONARA à Limbé) ; gros cercles bleus (Douala, Yaoundé, Nkoteng, Mbandjock, Dibombari) ; triangles verts (Douala, Bertoua, Ngaoundéré) ; losanges gris (Figuil, Douala, Yaoundé). Tailles proportionnelles aux effectifs.
    \item \textit{Infrastructures :} ports (Douala, Kribi, Limbé, Garoua), barrages (Edéa, Song Loulou, Lagdo, Memve'ele), pipeline (Doba--Kribi), rail (Douala--Ngaoundéré et embranchements), routes N1, N3, N10.
    \item \textit{Enjeux :} hachures rouges (Sud-Ouest, Sud, Moungo), cadre pointillé violet (Kribi), étoiles sur Douala et Yaoundé.
\end{itemize}
\textbf{Commentaire :} vérifier la problématique (dualisme littoral / arrière-pays), les deux parties (agro-industrie littorale ; industrie du corridor Douala--Kribi--Edéa), l'usage des chiffres (pipeline d'environ \num{1070} km, puissances des barrages, rang de premier exportateur africain de banane), et la conclusion prospective (émergence 2035, maillage du territoire, énergie).
\end{corrige}

\coursec{Enjeux transversaux : Agro-industrie, emploi, environnement et trajectoire d'émergence}

Après avoir cartographié les pôles de l'économie moderne (section 5), nous devons maintenant comprendre les \textbf{dynamiques structurelles} qui conditionnent leur avenir. Le Cameroun ne peut se contenter d'être un « grenier » et un « atelier » tournés vers l'extérieur sans maîtriser la valeur ajoutée. Cette section analyse les leviers (agro-industrie, capital humain), les freins (financement, énergie, environnement) et la stratégie de rupture (PND30, Vision 2035) pour passer d'une économie de rente à une économie de transformation.

\courssub{L'agro-industrie : cœur de la création de valeur ajoutée}

L'agriculture de plantation et l'industrie ne sont pas deux mondes étanches : leur articulation dans l'\textbf{agro-industrie} est le principal levier de richesse. Il s'agit de transformer sur place la matière première brute (1\textsuperscript{re} transformation) puis de fabriquer des biens finis ou intermédiaires (2\textsuperscript{ème}/3\textsuperscript{ème} transformation) avant l'exportation ou la consommation locale.

\begin{definition}[Taux de transformation locale]
Indicateur clé de l'industrialisation d'une filière. Il se calcule :
\[
\text{Taux de transformation (\%)} = \frac{\text{Valeur de la production transformée localement}}{\text{Valeur de la production brute totale}} \times 100
\]
Il distingue la transformation \textbf{sur site} (usine intégrée à la plantation : huile brute, CRM, coton égrené) de la transformation \textbf{en aval} (usines distinctes : savonneries, filatures, menuiseries, chocolateries).
\end{definition}

\begin{attention}[1\textsuperscript{re} vs 2\textsuperscript{ème}/3\textsuperscript{ème} transformation]
Ne pas confondre :
\begin{itemize}
    \item \textbf{1\textsuperscript{re} transformation} : souvent réalisée \emph{sur le site de production} (huileries de la CDC/SOCAPALM, usines de crêpage HEVECAM, égrenages SODECOTON, scieries). Elle stabilise le produit pour le transport mais la valeur ajoutée reste modérée.
    \item \textbf{2\textsuperscript{ème}/3\textsuperscript{ème} transformation} : réalisée dans des \emph{usines spécialisées} (souvent à Douala/Bonabéri, Yaoundé, Bafoussam) : raffinage huile $\to$ oléine/stéarine/savon (AZUR, SODECOTON) ; CRM $\to$ gants/pneus/chambres à air ; bois scié $\to$ placage/contre-plaqué/meubles ; coton fibre $\to$ fil/tissu/vêtement (CICAM) ; fèves cacao $\to$ liqueur/beurre/poudre/chocolat (NEO INDUSTRY, SIC CACAOS, CHOCOCAM). C'est là que se joue l'emploi qualifié et la capture de la rente.
\end{itemize}
\end{attention}

\begin{propriete}[Chaîne de valeur (Value Chain) et intégration]
Une chaîne de valeur englobe l'ensemble des activités (conception, production, transformation, logistique, distribution, SAV) qui ajoutent de la valeur au produit.
\begin{itemize}
    \item \textbf{Intégration verticale} : la firme contrôle l'amont (plantation) et l'aval (usine, marque). Ex : \textbf{SOCAPALM} (palmeraies + huileries + raffinerie AZUR + distribution) ou \textbf{HEVECAM} (hévéas + usines CRM + projet gants). Cela sécurise l'approvisionnement et capte la marge totale.
    \item \textbf{Intégration horizontale} : diversification dans des activités connexes (ex : plantation qui produit aussi de l'huile de palme, du palmiste, du biogaz à partir des effluents).
\end{itemize}
L'objectif du PND30 est de pousser les firmes vers l'intégration verticale complète (jusqu'au produit fini grand public).
\end{propriete}

\begin{popfigure}
\begin{tikzpicture}
\begin{axis}[
    width=\linewidth,
    height=0.6\linewidth,
    ybar=5pt,
    bar width=12pt,
    ymin=0, ymax=100,
    ylabel={Taux de transformation (\%)},
    symbolic x coords={Bois, Cacao, Coton, Palme, Caoutchouc, Aluminium},
    xtick=data,
    xticklabel style={rotate=30, anchor=east, font=\small},
    legend style={at={(0.5,-0.3)}, anchor=north, legend columns=-1, font=\small},
    nodes near coords={\pgfmathprintnumber\pgfplotspointmeta\%},
    nodes near coords align={vertical},
    grid=major,
]
\addplot[popblue!80, fill=popblue!60] coordinates {
    (Bois, 8) (Cacao, 30) (Coton, 5) (Palme, 65) (Caoutchouc, 40) (Aluminium, 100)
};
\addplot[only marks, mark=*, mark size=3pt, poporange, thick] coordinates {
    (Bois, 50) (Cacao, 50) (Coton, 20) (Palme, 85) (Caoutchouc, 70) (Aluminium, 100)
};
\legend{Taux actuel (estim. 2020), Cible PND30 / Vision 2035}
\end{axis}
\node[anchor=north west, font=\footnotesize\itshape, text=popmuted] at (current bounding box.north west) {Source : MINEPAT, MINCOMMERCE, Rapports PND30 -- Aluminium (ALUCAM) = transformation locale quasi-totale de la bauxite/alumine importée.};
\end{tikzpicture}
\legende{Graphique combiné : Taux de transformation locale par filière stratégique. Barres bleues = situation actuelle (faible pour bois, coton, cacao). Points orange = cibles PND30 (ambition de rupture). La palme et l'hévéa sont les plus avancées grâce à l'intégration verticale historique.}
\end{popfigure}

\begin{exemplebox}[Filières phares : état des lieux et trajectoire]
\begin{itemize}
    \item \textbf{Palme (SOCAPALM, CDC, SPFS)} : 1\textsuperscript{re} transfo sur site (huile brute, tourteaux). 2\textsuperscript{ème} transfo : Raffinerie AZUR (Edéa) $\to$ oléine (huile de table), stéarine (margarine, savon), PFAD (biodiesel, aliments). \textbf{Atout} : Intégration quasi-complète. \textbf{Cible} : 85\% (biodiesel, oléochimie).
    \item \textbf{Hévéa (HEVECAM, CAMVERT)} : 1\textsuperscript{re} transfo : CRM (Caoutchouc Regenerated/Mill) et TSR. 2\textsuperscript{ème} transfo : Usine de gants (projet HEVECAM), rechapage pneus. \textbf{Cible} : 70\% (pneus neufs, articles médicaux).
    \item \textbf{Bois (Grumes $\to$ Sciage $\to$ Placage/Contre-plaqué/Meubles)} : Interdiction export grumes (2017, dérogations dégressives). Taux actuel $\approx$ 8\% (sciage brut). \textbf{Cible PND30 : 50\%}. Nécessite usines de 2\textsuperscript{ème}/3\textsuperscript{ème} transfo (meubles, panneaux) et énergie stable.
    \item \textbf{Cacao (NEO INDUSTRY, SIC CACAOS, CHOCOCAM)} : 1\textsuperscript{re} transfo : liqueur, beurre, poudre (export vrac). 2\textsuperscript{ème} : Chocolat (CHOCOCAM), pâte à tartiner. \textbf{Cible : 50\%}. Enjeu : qualité fèves, énergie, logistique export produit fini (chaîne du froid).
    \item \textbf{Coton (SODECOTON, CICAM)} : 1\textsuperscript{re} transfo : égrenage (fibre + graines $\to$ huile/cotol). 2\textsuperscript{ème} : Filature/Tissage (CICAM Garoua) $\to$ tissu, vêtement. \textbf{Taux actuel $\approx$ 5\%}. \textbf{Cible : 20\%}. Enjeu : compétitivité vs import Asie (tissu wax, vêtements d'occasion), coût énergie/main-d'œuvre.
\end{itemize}
\end{exemplebox}

\courssub{Emploi : le secteur moderne, pourvoyeur d'emplois formels mais insuffisant}

Le secteur moderne (plantations + industrie formelle) structure le marché du travail camerounais par sa capacité à offrir des \textbf{contrats salariés}, la protection sociale (CNPS) et des salaires supérieurs au SMIG, mais son volume reste limité face à l'informel.

\begin{aretenir}[Chiffres clés Emploi Moderne (Estimations INS/Ministères, ~2020-2023)]
\begin{itemize}
    \item Part dans l'emploi total : \textbf{10 à 15\%} seulement (vs 85-90\% secteur informel / agriculture familiale).
    \item Part dans l'emploi \textbf{formel salarié} : \textbf{$\sim$ 80\%}.
    \item \textbf{Plantations (CDC, SOCAPALM, HEVECAM, SAFACAM, PHP, SODECOTON)} : $\sim$ 50 000 à 70 000 emplois permanents + \textbf{forte saisonnalité} (cueillette, entretien) $\to$ emplois précaires, journaliers, souvent migrants internes (Nord $\to$ Sud-Ouest/Littoral).
    \item \textbf{Industrie manufacturière (Douala, Bonabéri, Yaoundé, Bafoussam)} : Emplois plus qualifiés (techniciens maintenance, chimistes, ingénieurs process, logisticiens, commerciaux). Besoin croissant en \textbf{compétences techniques} (soudure, automatisme, qualité HSE).
\end{itemize}
\end{aretenir}

\begin{attention}[Adéquation Formation-Emploi \& Précarité]
Deux écueils majeurs :
\begin{enumerate}
    \item \textbf{Gap compétences} : Les ENSET, IUT, Mines, Écoles d'ingénieurs forment, mais les programmes ne collent pas toujours aux besoins réels des usines (maintenance prédictive, automatisme Siemens/Schneider, traitement eaux, certifications ISO). Résultat : postes vacants \emph{et} diplômés au chômage.
    \item \textbf{Sous-traitance \& Précarité} : Les grandes firmes externalisent des activités à risque ou saisonnières (nettoyage, manutention, sécurité, même certaine maintenance) vers des PME sous-traitantes. Les salariés de ces PME ont des contrats plus fragiles, salaires inférieurs, moins de formation. Cela « maquille » la précarité dans les stats du donneur d'ordre.
\end{enumerate}
\end{attention}

\begin{methode}[Analyser l'adéquation formation-emploi dans une filière (Démarche d'enquête)]
\begin{enumerate}
    \item \textbf{Identifier les métiers en tension} : Consulter les rapports de l'ONEC (Observatoire National de l'Emploi et de la Formation) ou enquêter auprès DRH de 3-4 entreprises leaders (ex: SOCAPALM, ALUCAM, BRASERIES, CICAM).
    \item \textbf{Recenser l'offre de formation locale} : Programmes des IUT/ENSET/Universités proches (spécialités : Génie chimique, Agro-alimentaire, Maintenance industrielle, Qualité/HSE, Logistique portuaire).
    \item \textbf{Comparer référentiels} : Croiser « Compétences requises (fiches de poste) » vs « Compétences acquises (programmes, stages) ». Repérer les manques (ex: conduite de ligne automatisée, gestion GMAO, anglais technique, normes ISO 22000/14001).
    \item \textbf{Proposer des leviers} : Apprentissage/alternance obligatoire, chaires d'entreprises à l'université, formation continue certifiante pour salariés, équipement plateaux techniques (CFA, Lycées techniques).
\end{enumerate}
\end{methode}

\courssub{Transformation locale : les objectifs chiffrés du PND30 et de la Vision 2035}

La \textbf{Stratégie Nationale de Développement 2020-2030 (SND30/PND30)} et la \textbf{Vision 2035} fixent le cap : « Sortir de l'exportation de matières premières brutes ». Le levier principal est l'interdiction progressive de l'export brut et les incitations à l'investissement industriel.

\begin{savaistu}[L'interdiction des grumes : un précédent juridique fort]
Depuis \textbf{janvier 2017}, le Cameroun a interdit l'exportation des grumes (bois brut) via la Loi de Finances (art. 14 puis décret d'application). Des dérogations dégressives ont été accordées aux opérateurs ayant investi dans des unités de transformation locale (scieries, placage, contre-plaqué). C'est un \textbf{cas d'école de politique industrielle offensive} : forcer la transformation par la contrainte réglementaire, couplée à l'incitation (exonérations fiscales/douanières sur équipements). Résultat : taux de transformation bois passé de $<5\%$ (2010) à $\sim 8-10\%$ (2020), mais l'objectif 50\% (2030) exige une accélération massive (usines meubles, panneaux, biomasse-énergie).
\end{savaistu}

\begin{exemplebox}[Objectifs PND30 par filière (Extraits) \& Investissements]
\begin{center}
\begin{tabularx}{\linewidth}{|l|X|X|X|}\hline
\textbf{Filière} & \textbf{Taux 2020 (est.)} & \textbf{Cible 2030} & \textbf{Leviers majeurs PND30} \\ \hline
Bois & $\sim 8\%$ & \textbf{50\%} & Interdiction grumes stricte, zones industrielles dédiées (Bertoua, Douala), énergie, formation ébénisterie. \\ \hline
Cacao & $\sim 30\%$ & \textbf{50\%} & Incitations usines 2\textsuperscript{ème} transfo (beurre, poudre, chocolat), qualité fèves (fermentation), marque « Cacao Cameroun ». \\ \hline
Coton & $\sim 5\%$ & \textbf{20\%} & Relance CICAM (filature/tissage), filière textile-habillement, préférence nationale marchés publics (tenues scolaires, militaires). \\ \hline
Palme/Hévéa & $60-70\%$ (1\textsuperscript{re} transfo) & \textbf{85-90\%} (produits finis) & Oléochimie (biodiesel, tensioactifs), industrie du gant/pneu, valorisation effluents (biogaz). \\ \hline
Aluminium (ALUCAM) & $100\%$ (transfo locale) & \textbf{Intégration aval} & Projet extension (potline 4), énergie dédiée (Nachtigal+), produits laminés/extrudés (câbles, profilés). \\ \hline
\end{tabularx}
\end{center}
\textbf{Investissements industriels ciblés 2020-2030 :} \textbf{$\sim$ 3 000 milliards FCFA} (Public + Privé + PPP). Clé : \textbf{Pacte National pour l'Industrialisation (2023)} $\to$ Guichet unique investisseur, Zones Franches (Kribi, Douala, Garoua), Préférence nationale marchés publics (produits transformés locaux).
\end{exemplebox}

\begin{methode}[Calculer un taux de transformation à partir de données douanières/industrielles]
Deux approches complémentaires selon la data disponible :
\begin{enumerate}
    \item \textbf{Approche Douane (Commerce extérieur)} : 
    \[
    \text{Taux} = \frac{\text{Valeur Export Produits Transformés (SH 4/6 digits)}}{\text{Valeur Export Produits Bruts + Valeur Export Produits Transformés}} \times 100
    \]
    \textit{Exemple Coton :} Export Fibre brute (SH 5201) vs Export Fil (5205), Tissu (5208), Vêtements (61/62). Attention : ne compte pas la transformation pour le marché \emph{local} (non exporté).
    \item \textbf{Approche Production/Industrie (Bilan Matière)} : 
    \[
    \text{Taux} = \frac{\text{Consommation Locale Matière Première par les Usines}}{\text{Production Nationale Totale Matière Première}} \times 100
    \]
    \textit{Exemple Bois :} Volume grumes entrées scieries/usines placage / Volume total grumes abattues (officiel + informel estimé). Plus pertinent pour mesurer l'industrialisation réelle (marché local inclus).
\end{enumerate}
\textbf{Piège} : Distinguer « 1\textsuperscript{re} transfo » (sciage, égrenage, huilerie brute) et « transfo profonde » (meubles, vêtements, savon) dans les nomenclatures douanières (SH).
\end{methode}

\courssub{Environnement \& RSE : la contrainte incontournable}

L'industrialisation ne peut se faire au prix de la destruction du capital naturel. Les firmes modernes sont soumises à un cadre légal strict et à des pressions marchés (certifications).

\begin{definition}[Cadre légal environnemental camerounais (Loi n°96/12 du 5 août 1996 + Décrets d'application)]
\begin{itemize}
    \item \textbf{EIES (Étude d'Impact Environnemental et Social)} : Obligatoire pour \emph{tout} projet industriel, plantation $> 50$ ha, infrastructure lourde. Doit inclure Plan de Gestion Environnementale et Sociale (PGES). Validée par le MINEE/MINEPDED $\to$ Certificat de Conformité Environnementale (CCE) renouvelable (souvent 3 ans).
    \item \textbf{Audit environnemental} : Périodique pour installations existantes (classées ICPE - Installations Classées pour la Protection de l'Environnement).
    \item \textbf{Pollueur-Payeur} : Principe constitutionnel. Amendes, obligation de réhabilitation, fermeture administrative possible.
\end{itemize}
\end{definition}

\begin{exemplebox}[Principaux impacts \& Réponses sectorielles]
\begin{center}
\begin{tabularx}{\linewidth}{|l|X|X|}\hline
\textbf{Secteur / Site} & \textbf{Impacts majeurs} & \textbf{Mesures / Certifications} \\ \hline
\textbf{Plantations (Palme, Hévéa, Banane)} & Déforestation (conversion forêt $\to$ monoculture), érosion, pollution pesticides/engrais, effluents huileries (POME - Palm Oil Mill Effluent), perte biodiversité. & \textbf{RSPO} (Roundtable on Sustainable Palm Oil) : certification durabilité (zéro déforestation HCV/HCS, gestion tourbières, droits travailleurs/communautés). \textbf{ISO 14001} (SME). Traitement anaérobie POME $\to$ biogaz + compost. \\ \hline
\textbf{Industrie Lourde (ALUCAM, SONARA, CIMENCAM)} & Émissions atmosphériques (CO2, SOx, NOx, fluorures, poussières), rejets thermiques/eaux usées, déchets dangereux (bâches, réfractaires, catalyseurs). & \textbf{ISO 14001}, \textbf{ISO 50001} (énergie). ALUCAM : pots à précombustion, capture fluorures. SONARA : épuration eaux, projet CCR (craquage catalytique) pour carburants propres. CIMENCAM : substitution combustibles (pneus, balles riz, biomasse) $\to$ moins clinker. \\ \hline
\textbf{Agro-industrie légère (Brasseries, Savonneries, Textile)} & Rejets organiques forts (DBO/DCO), colorants, sels, emballages plastiques. & Stations d'épuration biologiques (boues activées, lagunage), économie circulaire (drêches $\to$ aliment bétail, CO2 récupéré), \textbf{consigne bouteilles} (verre), collecte plastique. \\ \hline
\textbf{Bois (Scieries, Placage)} & Déforestation (approvisionnement illégal), sciures/copeaux (brûlage à l'air libre), colles (formaldéhyde). & \textbf{FSC} (Forest Stewardship Council) / \textbf{PEFC} / \textbf{OLB} (Origine et Légalité Bois) : traçabilité, gestion durable concessions. Valorisation résidus (briquettes, panneaux particules, énergie). \\ \hline
\end{tabularx}
\end{center}
\end{exemplebox}

\begin{attention}[Greenwashing vs RSE réelle]
Méfiance vis-à-vis des rapports RSE « vitrine ». Vérifier : 1) Certification \textbf{tierce partie indépendante} (RSPO, FSC, ISO 14001 certifiés par SGS, Bureau Veritas, Ecocert -- pas auto-déclaration). 2) \textbf{Publication des indicateurs} : volumes effluents traités, t CO2 eq évités, \% énergie renouvelable, nombre plaintes communautés traitées. 3) \textbf{Intégration dans la gouvernance} : Comité RSE au Conseil d'Administration, rémunération dirigeants indexée sur objectifs ESG.
\end{attention}

\courssub{Financement, Compétitivité \& Pacte National pour l'Industrialisation}

Pourquoi l'industrialisation tarde-t-elle malgré les atouts ? Les « coûts de faire des affaires » (Doing Business) restent élevés.

\begin{propriete}[Les 5 piliers de la compétitivité industrielle camerounaise (Diagnostic PND30/Pacte 2023)]
\begin{enumerate}
    \item \textbf{Énergie} : Coupures fréquentes, tarif industriel élevé ($\sim$ 50-60 FCFA/kWh vs $<30$ au Maroc/Égypte), déficit de pointe. \textbf{Solution} : Barrages Nachtigal (420 MW, 2024/25), Grand Eweng (800-1200 MW), Lom Pangar (régulation), solaire/biomasse industriels. \textbf{Clé} : Énergie \emph{fiable, abondante, bon marché}.
    \item \textbf{Financement} : Taux d'intérêt courts 7-9\% (vs 3-4\% zone CFA benchmark), exigences de garanties foncières (titres fonciers rares), maturité courte (3-5 ans) inadaptée aux cycles industriels (10-15 ans). \textbf{Solution} : Banque de Développement (BCD) renforcée, fonds de garantie PME, marché obligataire corporate, Islamic Finance (Sukuk).
    \item \textbf{Logistique / Infrastructures} : Port de Douala (congestion, temps séjour $\sim$ 12-15 j vs 3-5 j performants), route Douala-Ndjamena (corridor CEMAC), chemin de fer (Camrail : capacité saturée, vétusté), aéroport fret (Nsimalen, Douala). \textbf{Solution} : Port en eau profonde Kribi (phase 2), autoroute Yaoundé-Douala, réhabilitation rail, corridors secs.
    \item \textbf{Capital Humain} : Pénurie techniciens supérieurs (BTS, DUT, Ingénieurs) opérationnels « terrain » (maintenance, process, qualité, HSE). Déficit formation professionnelle duale (entreprise/école). \textbf{Solution} : Réforme ENSET/IUT/Lycées techniques, apprentissage obligatoire, chaires entreprises, centres d'excellence sectoriels (Bois, Agro, Mines, Numérique).
    \item \textbf{Fiscalité / Administration} : Multiplicité taxes (Centre des Impôts, Commune, Douanes, Ministères sectoriels, Préfeture), instabilité fiscale (lois de finances annuelles changeantes), lourdeurs administratives (délais licences, visas). \textbf{Solution} : \textbf{Pacte National pour l'Industrialisation (Juin 2023)} $\to$ Guichet unique (API/CEPI), harmonisation fiscalité, stabilité 10 ans pour nouveaux investissements, préférence nationale marchés publics (seuil 15-20\% prix sup. accepté pour local).
\end{enumerate}
\end{propriete}

\begin{exoresolu}[Analyse d'un projet d'usine de transformation cacao (Chocolaterie) -- Faisabilité simplifiée]
\textbf{Contexte :} Un investisseur veut créer une usine de 10 000 t/an de fèves $\to$ Liqueur/Beurre/Poudre/Chocolat à Douala. Investissement : 15 Mrd FCFA.
\textbf{Données :} Prix fève bord champ : 1 300 FCFA/kg. Prix export fève brute (FOB Douala) : 2 200 FCFA/kg. Prix export beurre/poudre (moyen) : 4 500 FCFA/kg (équivalent fève). Coût transformation (énergie, main-d'œuvre, emballage, logistique, amortissement) : 1 000 FCFA/kg fève entrée. Taux récupération beurre+poudre $\approx$ 80\% poids fève.
\medskip
\textbf{Question :} Le projet crée-t-il de la valeur ajoutée locale par rapport à l'export brut ? Calculer la marge brute par kg de fève transformée vs export brut.
\medskip
\textbf{Solution :}
\begin{align*}
\text{Recette Export Brut} &= 2\,200 \text{ FCFA/kg fève} \\
\text{Recette Produits Transformés (équiv. fève)} &= 4\,500 \text{ FCFA/kg} \times 0,80 \text{ (rendement)} = 3\,600 \text{ FCFA/kg fève entrée} \\
\text{Coût Achat Fève + Transfo} &= 1\,300 + 1\,000 = 2\,300 \text{ FCFA/kg fève entrée} \\
\text{Marge Brute Transformation} &= 3\,600 - 2\,300 = \mathbf{1\,300 \text{ FCFA/kg fève entrée}} \\
\text{Marge Export Brut (négoce pur)} &= 2\,200 - 1\,300 = 900 \text{ FCFA/kg (marge commerciale seulement)} \\
\text{Gain Valeur Ajoutée Locale} &= 1\,300 - 900 = \mathbf{400 \text{ FCFA/kg supplémentaire captée localement}}
\end{align*}
\textbf{Interprétation :} L'usine génère 400 FCFA/kg de VA \emph{supplémentaire} vs négoce brut. Sur 10 000 t/an $\to$ \textbf{4 Mrd FCFA/an de VA créée localement} (salaires, impôts, profits réinvestis, sous-traitance). \textbf{Condition sine qua non :} Énergie stable (coût transfo 1 000 FCFA/kg tenu), logistique export produit fini (conteneurs reefers pour beurre), accès financement long terme.
\end{exoresolu}

\courssub{Intégration sous-régionale : le marché CEMAC/ZLECAf comme débouché naturel}

Le marché camerounais (28 M hab.) est trop petit pour absorber une production industrielle de masse. L'intégration régionale change la donne.

\begin{popfigure}
\begin{tikzpicture}[scale=0.9, every node/.style={font=\small}]
% Base map: simplified CEMAC countries positions (approx lat/long projected)
% Coords: (x,y) relative. Douala ~ (0,0). Yaounde ~ (1.5, 1.5). N'Djamena ~ (8, 6). Bangui ~ (5, 3). Libreville ~ (1, -2). Malabo ~ (-1, -1). Brazzaville ~ (2, -1.5).
% Countries boxes
\node[draw=popblue!50, fill=popblue!5, rounded corners, minimum width=2.5cm, minimum height=1.8cm, align=center] (CMR) at (0,0){\textbf{CAMEROUN}\\\small(Douala, Yaoundé)};
\node[draw=popgreen!50, fill=popgreen!5, rounded corners, minimum width=2cm, minimum height=1.5cm, align=center] (TCD) at (8,5.5){\textbf{TCHAD}\\\small(N'Djaména)};
\node[draw=poporange!50, fill=poporange!5, rounded corners, minimum width=2cm, minimum height=1.5cm, align=center] (CAF) at (5,3){\textbf{RCA}\\\small(Bangui)};
\node[draw=poppurple!50, fill=poppurple!5, rounded corners, minimum width=2cm, minimum height=1.5cm, align=center] (GAB) at (1,-2.5){\textbf{GABON}\\\small(Libreville)};
\node[draw=poppink!50, fill=poppink!5, rounded corners, minimum width=2cm, minimum height=1.5cm, align=center] (GNQ) at (-1,-1){\textbf{GUINÉE EQ.}\\\small(Malabo)};
\node[draw=popgold!50, fill=popgold!5, rounded corners, minimum width=2cm, minimum height=1.5cm, align=center] (COG) at (2,-1.5){\textbf{CONGO}\\\small(Brazzaville)};
% Corridors (Road/Rail)
\draw[popdark, thick, dashed, -{Stealth[length=3mm]}] (CMR.east) -- ++(1,0) node[midway, above, font=\tiny] {Route/Rail} -- (CAF.west);
\draw[popdark, thick, dashed, -{Stealth[length=3mm]}] (CAF.north east) -- ++(1.5,1) node[midway, above, font=\tiny] {Corridor Douala-Ndjaména} -- (TCD.south west);
\draw[popdark, thick, dashed, -{Stealth[length=3mm]}] (CMR.south) -- (GAB.north);
\draw[popdark, thick, dashed, -{Stealth[length=3mm]}] (CMR.south west) -- (GNQ.north east);
\draw[popdark, thick, dashed, -{Stealth[length=3mm]}] (GAB.south east) -- (COG.north west);
% Export Flows from Cameroun (Douala/Yaounde area) -> Neighbors
% Arrow thickness proportional to trade volume (qualitative)
% Ciment, Savon, Bière, Farine, Bois ouvré
\def\flow#1#2#3#4#5#6{ % from, to, label, color, width
    \draw[#5, line width=#6, -{Stealth[length=2.5mm]}] (#1) .. controls #2 .. (#3) node[midway, sloped, above, font=\tiny, text=popdark] {#4};
}
% Flows from CMR center approx (0.5, 0.5)
\flow{CMR}{+(3,2) and +(2,1)}{TCD}{Ciment, Clinker, Farine, Savon}{popblue}{1.5pt}
\flow{CMR}{+(2,1) and +(1,0.5)}{CAF}{Ciment, Savon, Bière, Médicaments}{popgreen}{1.2pt}
\flow{CMR}{+(0.5,-1) and +(0,-0.5)}{GAB}{Bois ouvré, Produits agro, BTP}{poppurple}{1pt}
\flow{CMR}{+(-0.5,-0.5) and +(0,-0.3)}{GNQ}{Bois, Agro, BTP}{poppink}{0.8pt}
\flow{CMR}{+(1,-0.8) and +(0.5,-0.4)}{COG}{Ciment, Savon, Produits chimiques}{popgold}{1pt}
% ZLECAf label
\node[draw=popdark, fill=popcream, rounded corners, font=\small\bfseries, align=center, anchor=north] at (4,-3.5) {ZLECAf : Marché unique 1,3 Md hab.\\Suppression droits de douane 90\% lignes tarifaires\\Règles d'origine, paiements PAPSS};
\end{tikzpicture}
\legende{Schéma des flux commerciaux intra-CEMAC depuis le Cameroun (principal fournisseur industriel de la zone) et corridors logistiques. L'entrée en vigueur de la ZLECAf (Zone de Libre-Échange Continentale Africaine) élargit le débouché à 1,3 milliard de consommateurs, à condition de lever les barrières non tarifaires (certification, transport, paiement).}
\end{popfigure}

\begin{aretenir}[Synthèse des enjeux transversaux pour l'émergence]
\begin{itemize}
    \item \textbf{Agro-industrie} : Passer de 1\textsuperscript{re} transfo (faible VA) à 2\textsuperscript{ème}/3\textsuperscript{ème} transfo (forte VA, emploi qualifié) via intégration verticale et interdiction export brut.
    \item \textbf{Emploi} : Secteur moderne = 10-15\% emploi total mais 80\% emploi formel. Défi : adéquation formation-emploi (maintenance, process, HSE) et lutte contre précarité sous-traitance.
    \item \textbf{Environnement} : Cadre légal strict (EIES, ICPE, Pollueur-Payeur) + certifications marchés (RSPO, FSC, ISO 14001). Économie circulaire (effluents $\to$ biogaz, résidus $\to$ énergie) = compétitivité.
    \item \textbf{Compétitivité} : 5 piliers (Énergie, Financement, Logistique, Capital Humain, Fiscalité) adressés par Pacte National Industrialisation 2023 (guichet unique, stabilité fiscale, préférence nationale).
    \item \textbf{Intégration régionale} : CEMAC (6 pays, 60 M hab.) + ZLECAf (1,3 Md hab.) = débouché naturel pour ciment, agro-alimentaire, bois ouvré, pharmaceutique. Nécessite corridors fluides (Kribi, rail, route) et harmonisation normes.
\end{itemize}
\end{aretenir}

\newpage

\coursec{Activité d'intégration}

\coursec{L'économie moderne au Cameroun : Agriculture de plantations et Industrie}

\courssub{Objectifs de l'activité}
\begin{itemize}
    \item Mobiliser les connaissances sur l'\cle{agriculture de grandes plantations} (cacao), le \cle{secteur industriel} (agro-industrie), les \cle{facteurs de localisation} (énergie, transport, main-d'œuvre, matière première, foncier) et les \cle{enjeux transversaux} (emploi, environnement, émergence).
    \item Appliquer une \cle{démarche multicritères} pour analyser des données géographiques chiffrées et cartographiques.
    \item Réaliser un \cle{croquis de synthèse} respectant la sémiologie graphique (figuration surfacique, linéaire, ponctuelle, légende organisée).
    \item Rédiger une \cle{note d'argumentation} structurée, justifiant un choix d'aménagement du territoire au regard de la compétitivité, de l'inclusion sociale et de la durabilité.
    \item Communiquer oralement et débattre de manière contradictoire pour confronter les stratégies d'implantation.
\end{itemize}

\courssub{Situation}
\textbf{Contexte :} Dans le cadre de la \cle{Stratégie Nationale de Développement 2020-2030 (SND30)} et de la politique d'\cle{import-substitution} / \cle{transformation locale}, l'État du Cameroun, via le \cle{Ministère de l'Industrie, des Mines et du Développement Technologique (MINIMIDT)} et l'\cle{Agence de Promotion des Investissements (API)}, lance un appel à manifestation d'intérêt pour l'implantation d'une \textbf{unité industrielle de transformation de fèves de cacao en poudre, beurre et liqueur de cacao} (capacité : 50 000 tonnes de fèves/an). L'objectif national est de porter le taux de transformation locale de \SI{\sim 30}{\%} (niveau 2023) à \SI{50}{\%} à l'horizon 2030.

\textbf{Cahier des charges technique simplifié de l'usine :}
\begin{itemize}
    \item \textbf{Approvisionnement :} 50 000 t de fèves sèches/an (rendement moyen plantation villageoise : 400--600 kg/ha ; rendement agro-industriel : 1,5--2 t/ha).
    \item \textbf{Énergie :} Puissance appelée \SI{50}{MW} en continu (processus thermiques + froid industriel), stabilité réseau critique (variation $< \pm 5\%$).
    \item \textbf{Main-d'œuvre :} 200 emplois directs qualifiés (techniciens, ingénieurs, logisticiens) + 500 emplois indirects (transit, maintenance, services) ; bassin de main-d'œuvre agricole saisonnière pour la collecte.
    \item \textbf{Logistique/Export :} 80\% de la production exportée (CEMAC, UE, USA) $\rightarrow$ conteneurisée (port en eau profonde requis) ; 20\% marché local/ sous-régional (route/rail).
    \item \textbf{Foncier :} Assiette de \SI{25}{ha} (usine, stockage, aires de manœuvre, traitement effluents, extension future), constructible, viabilisable, titre foncier sécurisé (Loi 2013 sur les incitations aux investissements privés).
    \item \textbf{Environnement :} Étude d'impact environnemental et social (EIES) obligatoire ; gestion des effluents (eaux de lavage, coques), émissions (chaudières), certification \cle{ISO 14001} visée.
\end{itemize}

\textbf{Dossier documentaire fourni aux groupes (extraits) :}

\begin{center}
\begin{tabularx}{\linewidth}{|l|X|X|X|X|}\hline
\textbf{Indicateur} & \textbf{Douala/Bonabéri} & \textbf{Kribi/Zone portuaire} & \textbf{Ebolowa/Sangmélima} & \textbf{Bertoua/Est} \\ \hline
\textbf{Production cacao régionale (t/an, moy. 2020-23)} & 15 000 (Littoral) & 35 000 (Sud : Océan, Vallée-du-Ntem) & 25 000 (Sud : Mvila, cœur de bassin) & 35 000 (Est : Haut-Nyong, Kadey) \\ \hline
\textbf{Disponibilité énergie (Réseau interconnecté Sud - RIS)} & \textbf{Excellent} (Centrale de Logbaba 85 MW, Songloulou 388 MW, Edéa 263 MW, proximité gaz) & \textbf{Très bon} (Raccordement RIS, proximité barrage Memve'ele 211 MW, ligne 225 kV) & \textbf{Moyen/Faible} (Réseau 90 kV saturé, coupures fréquentes, pas de grande centrale proche) & \textbf{Faible} (Réseau isolé 30 kV, groupes thermiques, Lom Pangar 30 MW non raccordé industriel) \\ \hline
\textbf{Coût kWh industriel (FCFA, est. 2024)} & 45--55 (Gaz/ Hydro) & 50--60 (Hydro/ Import) & 70--90 (Hydro + pertes ligne/ groupes secours) & 100--130 (Thermique diesel/ fuel) \\ \hline
\textbf{Infrastructures transport (État 2024)} & \textbf{Route :} Y dédoublé (N3, N5), saturé. \textbf{Rail :} Camrail (Douala-Yaoundé-Ngaoundéré), gare triage. \textbf{Port :} Douala (saturation, tirant d'eau 11,5 m) & \textbf{Route :} N17 (Kribi-Ebolowa) bonne, N7 (Kribi-Douala) en réhab. \textbf{Rail :} Projet Mbalam-Kribi (minerai), pas fret cacao. \textbf{Port :} \textbf{Kribi Phase 1} (eau profonde 16 m, capacité 1,4 M EVP, sous-exploité) & \textbf{Route :} N2 (Yaoundé-Ebolowa) bonne, N9 (Ebolowa-Sangmélima) dégradée. \textbf{Rail :} Aucun. \textbf{Port :} Douala (300 km) ou Kribi (120 km via N17) & \textbf{Route :} N10 (Yaoundé-Bertoua) bonne, axes secondaires mauvais. \textbf{Rail :} Aucun. \textbf{Port :} Douala (650 km) \\ \hline
\textbf{Bassin d'emploi / Formation} & \textbf{Très grand} (Douala 3,5 M hab, IUT, ENSET, Lycées tech) & \textbf{Moyen} (Kribi 80k hab, Ebolowa 100k hab, Lycées tech) & \textbf{Faible} (Zone rurale, lycées agricoles) & \textbf{Faible} (Bertoua 150k hab, Lycée tech) \\ \hline
\textbf{Foncier (Disponibilité / Conflits / Coût)} & \textbf{Faible / Élevé / Conflits} (Zones franches saturées, litiges coutumiers, prix $\uparrow$) & \textbf{Élevé / Faible / Sécurisé} (Zone Franche Kribi, titres État, 25 ha disponibles) & \textbf{Moyen / Moyen / Coutumier} (Terres villageoises, procédure longue, risques spoliation) & \textbf{Élevé / Faible / Coutumier} (Grands espaces, chefferies, procédure API facilitée) \\ \hline
\textbf{Environnement / Réserves} & Zone urbaine dense, pollution existante, mangroves sensibles & \textbf{Parc National de Campo Ma'an} (proximité, corridor éléphants), forêt dense, EIES stricte & Forêt dense, biodiversité forte, bassins versants (Nyong, Ntem), plantations villageoises & Forêt dense, biodiversité, parc Boumba-Bek/Nki, populations autochtones (Baka) \\ \hline
\textbf{Incitations Loi 2013 / Zones Franches} & Zone Franche Douala (saturée) & \textbf{Zone Franche Kribi (nouvelle, exonérations 10 ans, TVA, droits douane)} & Zone d'appui prioritaire (rural), incitations spécifiques agriculture & Zone d'appui prioritaire (rural), incitations spécifiques agriculture \\ \hline
\end{tabularx}
\end{center}

\legende{Tableau 1 : Grille comparative multicritères des sites potentiels (Données synthétisées MINIMIDT, SONARA, ENEO, CAMRAIL, PAK, API, MINADER 2024).}

\textbf{Cartes thématiques jointes (simplifiées) :}
\begin{enumerate}
    \item Carte du bassin cacaoyer (Régions Centre, Sud, Est, Littoral, Sud-Ouest) : isoproduction (t/an), axes routiers structurants (N1, N2, N3, N5, N7, N10, N17), réseau ferroviaire (Camrail), barrages hydroélectriques (Songloulou, Edéa, Memve'ele, Lom Pangar), ports (Douala, Kribi), villes principales.
    \item Carte des zones franches et pôles industriels (Douala-Bonabéri, Kribi, Yaoundé-Nkolbisson, Bafoussam, Ngaoundéré).
\end{enumerate}

\courssub{Consignes}
\begin{enumerate}
    \item \textbf{Analyse documentaire (Groupe) :} À partir du Tableau 1 et des cartes, identifier \textbf{trois sites potentiels} parmi la liste (Douala, Kribi, Ebolowa, Sangmélima, Bertoua, Ngaoundéré). Justifier brièvement l'élimination des autres.
    \item \textbf{Grille multicritères (Groupe) :} Construire une grille d'analyse pondérée (note sur 20 par critère, coefficient selon importance : Énergie 3, Matière première 3, Export/Logistique 3, Foncier 2, Main-d'œuvre 2, Environnement 2, Incitations 1). Noter chaque site retenu. Calculer la note pondérée totale.
    \item \textbf{Choix du site optimal (Groupe) :} Sélectionner le site optimal. En cas d'égalité, départager par l'argument « \cle{compétitivité-coût complet} » (énergie + transport + foncier).
    \item \textbf{Croquis de localisation (Individuel / Groupe) :} Réaliser à l'échelle \textbf{1/5 000 000} un croquis intitulé \og \textbf{Localisation de l'unité de transformation cacao : Site de [Nom du site]} \fg. Figurer obligatoirement :
        \begin{itemize}
            \item Zone(s) d'approvisionnement en fèves (surfacique, hachures/orange).
            \item Emplacement précis de l'usine (ponctuel, symbole usine).
            \item Axes d'évacuation routiers (linéaire plein) et ferroviaires (linéaire pointillé) vers le port.
            \item Infrastructures énergie (barrage, ligne HT, centrale) alimentant le site.
            \item Port d'exportation (symbole port) et flux export (flèche proportionnelle).
            \item Légende organisée, échelle graphique, flèche Nord, titre.
        \end{itemize}
    \item \textbf{Note d'argumentation (Individuel, 1 page max) :} Rédiger une note structurée (Introduction, Analyse comparative, Choix retenu, Justification triptyque : Compétitivité / Inclusion sociale / Durabilité, Conclusion) justifiant le choix au regard du \cle{PND30} et de la \cle{Loi 2013}.
    \item \textbf{Restitution orale (Plénière) :} Chaque groupe présente son site, son croquis et sa note (5 min). Débat contradictoire animé par le professeur : \og Pourquoi pas Douala ? Pourquoi pas l'intérieur (Ebolowa/Bertoua) ? \fg. Synthèse professeur.
\end{enumerate}

\courssub{Ressources à mobiliser}
\begin{aretenir}[Notions clés du cours - Rappel]
    \textbf{Dualisme économique} : Coexistence secteur moderne (plantations, industrie formelle) / secteur traditionnel (paysannerie, informel).
    \textbf{Agriculture de plantation} : Capitalistique, salariale, mono-culture, exportation, foncier étatique/concession (ex: CDC, SOCAPALM, HEVECAM) vs paysannerie villageoise (petits exploitants, polyculture, vivrier + rente cacao/café).
    \textbf{Facteurs de localisation industrielle (Weber / Losch adaptés)} :
    \begin{itemize}
        \item \cle{Matière première} : Indice de perte de poids (cacao : faible, mais volume/coût transport collectage fort $\rightarrow$ proximité bassins).
        \item \cle{Énergie} : \cle{Électro-intensif} (50 MW) $\rightarrow$ impératif proximité source/réseau fiable (hydrocarbures/gaz/hydro).
        \item \cle{Main-d'œuvre} : Qualifiée (usine) + non qualifiée abondante (collecte).
        \item \cle{Marché/Export} : Coût transport final (conteneur) $\rightarrow$ proximité port eau profonde.
        \item \cle{Foncier} : Disponibilité, viabilisation, sécurité juridique (Titre foncier vs Droit coutumier).
        \item \cle{Incitations} : Zones franches (Loi 2013), exonérations fiscales/douanières.
    \end{itemize}
    \textbf{Pôles industriels Cameroun} : Douala (moteur, 60\% valeur ajoutée), Kribi (émergent, port/eau profonde/énergie), Yaoundé (administratif/secondaire), Garoua/Ngaoundéré (coton/agro).
    \textbf{Sémiologie croquis} : Surfacique (zones), Linéaire (flux/axes), Ponctuel (usines/ports/villes). Hiérarchie visuelle (épaisseur, couleur, taille). Légende structurée (Thèmes / Signes).
\end{aretenir}

\begin{methode}[Démarche multicritères d'aide à la décision (AMD)]
    \begin{enumerate}
        \item Définir les \textbf{critères} pertinents (C1...Cn).
        \item Attribuer des \textbf{poids (coefficients)} $\sum p_i = 1$ (ou total 100).
        \item Noter chaque alternative sur chaque critère (échelle commune, ex: 0--20).
        \item Calculer le \textbf{score pondéré} : $S_j = \sum (p_i \times note_{ij})$.
        \item Classer et \textbf{analyser la sensibilité} (que se passe-t-il si le poids Énergie passe de 3 à 4 ?).
    \end{enumerate}
\end{methode}

\courssub{Démarche et corrigé commenté}
\begin{exoresolu}[Dossier d'évaluation : Stratégie d'implantation unité cacao]
\textbf{Étape 1 : Identification de 3 sites potentiels \& Élimination}
\begin{itemize}
    \item \textbf{Retenus :} \textbf{Kribi} (Port eau profonde + Énergie + Zone Franche + Foncier), \textbf{Douala} (Bassin emploi + Énergie + Rail + Marché local), \textbf{Ebolowa} (Cœur bassin cacao Mvila + Proximité Kribi).
    \item \textbf{Éliminés :} \textbf{Sangmélima} (enclavement routier N9, pas d'énergie, pas de port), \textbf{Bertoua} (éloignement port 650 km, énergie thermique chère, enclavement), \textbf{Ngaoundéré} (hors bassin cacao, 1000 km port, énergie limitée).
\end{itemize}

\textbf{Étape 2 : Grille multicritères pondérée (Coefficients : Énergie 3, MP 3, Export 3, Foncier 2, MO 2, Env 2, Incit 1 = Total 16)}
\begin{center}
\begin{tabularx}{\linewidth}{|l|X|X|X|}\hline
\textbf{Critère (Coeff)} & \textbf{Kribi} & \textbf{Douala} & \textbf{Ebolowa} \\ \hline
Énergie (3) & Excellent (Memve'ele, RIS, Gaz) - 18/20 & Excellent (Logbaba, Songloulou, Gaz) - 19/20 & Faible (Réseau 90kV saturé) - 8/20 \\ \hline
Matière Première (3) & Bon (Bassin Océan/VN 35k t, rayon 80km) - 15/20 & Moyen (Littoral 15k t, concurrence) - 10/20 & Bon (Mvila 25k t, rayon 50km) - 16/20 \\ \hline
Export/Logistique (3) & \textbf{Excellent} (Port Kribi 16m, 5km, libre) - 19/20 & Moyen (Port Douala saturé, tirant eau limité) - 11/20 & Moyen (Route 120km vers Kribi, pas rail) - 12/20 \\ \hline
Foncier (2) & \textbf{Excellent} (ZES Kribi, titres État, 25ha dispo) - 18/20 & Faible (Saturé, cher, conflits) - 7/20 & Moyen (Coutumier, procédure longue) - 11/20 \\ \hline
Main-d'œuvre (2) & Moyen (Bassin 180k hab, formation à créer) - 11/20 & \textbf{Excellent} (3,5M hab, vivier qualifié) - 18/20 & Faible (Rural, peu qualifié) - 8/20 \\ \hline
Environnement (2) & Sensible (Parc Campo Ma'an, EIES lourde) - 10/20 & Dégradé (Urbain, pollution, moins sensible biodiversité) - 13/20 & Sensible (Forêt dense, bassins versants) - 10/20 \\ \hline
Incitations (1) & \textbf{Maximal} (ZES Kribi neuve, 10 ans) - 20/20 & Moyen (ZF Douala saturée) - 12/20 & Bon (Zone rurale prioritaire) - 15/20 \\ \hline
\textbf{SCORE PONDÉRÉ (sur 320)} & \textbf{254} & \textbf{208} & \textbf{181} \\ \hline
\end{tabularx}
\end{center}
\legende{Calcul détaillé : Kribi = $(18\times3)+(15\times3)+(19\times3)+(18\times2)+(11\times2)+(10\times2)+(20\times1) = 54+45+57+36+22+20+20 = \mathbf{254}$. Douala = $(19\times3)+(10\times3)+(11\times3)+(7\times2)+(18\times2)+(13\times2)+(12\times1) = 57+30+33+14+36+26+12 = \mathbf{208}$. Ebolowa = $(8\times3)+(16\times3)+(12\times3)+(11\times2)+(8\times2)+(10\times2)+(15\times1) = 24+48+36+22+16+20+15 = \mathbf{181}$. Total max = $20\times16 = 320$. \textbf{Kribi l'emporte nettement (79,4\%)}.}

\textbf{Étape 3 : Choix optimal $\rightarrow$ \textbf{KRIBI}.}
\textit{Justification départage :} Kribi combine l'avantage comparatif absolu sur \textbf{Énergie} (coût kWh bas, stabilité), \textbf{Export} (port eau profonde non saturé, coûts manutention inférieurs), \textbf{Foncier} (sécurisé, viabilisé, Zone Franche) et \textbf{Incitations}. Le déficit Matière Première (vs Ebolowa) est gérable par camions (80-120 km route bonne N17) ou futur rail. Le déficit Main-d'œuvre (vs Douala) est résolu par formation locale (lycée technique Kribi, partenariat usine/écoles) et attractivité salariale Zone Franche.

\textbf{Étape 4 : Croquis de localisation (Description pour réalisation élève)}
\begin{popfigure}
\begin{tikzpicture}[scale=0.85, every node/.style={font=\footnotesize}]
    % Échelle 1/5 000 000 -> 1 cm = 50 km. Zone : 2°-5°N / 8°-13°E ~ 330km x 550km -> ~6.6cm x 11cm
    % Fond de carte simplifié Cameroun (contour polygone grossier)
    \draw[popline, thick, fill=popcream!20] (0,0) -- (11,0) -- (11,6.6) -- (9.5,6.6) -- (8,7.5) -- (6,7.5) -- (4,6.6) -- (0,6.6) -- cycle;
    % Mer
    \draw[popblue, fill=popblueL!30] (-0.5,-0.5) rectangle (11.5,0);
    \node[popblue, rotate=90] at (-0.8, -0.2) {GOLFE DE GUINÉE};
    % Villes repères
    \node[circle, fill=popdark, inner sep=1.5pt, label=above:\textbf{Douala}] (Douala) at (2.5, 1.5) {};
    \node[circle, fill=popdark, inner sep=1.5pt, label=above:\textbf{Yaoundé}] (Yaounde) at (5.5, 3.0) {};
    \node[circle, fill=popdark, inner sep=1.5pt, label=below:\textbf{Kribi}] (Kribi) at (7.5, 0.5) {};
    \node[circle, fill=popdark, inner sep=1.5pt, label=left:\textbf{Ebolowa}] (Ebolowa) at (6.5, 1.8) {};
    \node[circle, fill=popdark, inner sep=1.5pt, label=above:\textbf{Bertoua}] (Bertoua) at (9.5, 4.0) {};
    % Barrages
    \node[regular polygon, regular polygon sides=3, fill=popblue, inner sep=2pt, label=above:Songloulou] at (3.0, 1.0) {};
    \node[regular polygon, regular polygon sides=3, fill=popblue, inner sep=2pt, label=above:Edéa] at (3.5, 1.5) {};
    \node[regular polygon, regular polygon sides=3, fill=popblue, inner sep=2pt, label=right:Memve'ele] at (7.0, 2.2) {};
    \node[regular polygon, regular polygon sides=3, fill=popblue, inner sep=2pt, label=above:Lom Pangar] at (9.0, 4.5) {};
    % Réseau HT (lignes pointillées)
    \draw[popblue, dashed, thick] (3.0,1.0) -- (5.5,3.0) -- (7.0,2.2); % Songloulou -> Yaoundé -> Memve'ele
    \draw[popblue, dashed, thick] (7.0,2.2) -- (7.5,0.5); % Memve'ele -> Kribi
    % Routes (Traits pleins)
    \draw[popdark, thick] (Douala) -- (Yaounde); % N3
    \draw[popdark, thick] (Yaounde) -- (Ebolowa); % N2
    \draw[popdark, thick] (Ebolowa) -- (Kribi); % N17
    \draw[popdark, thick] (Yaounde) -- (Bertoua); % N10
    % Rail (Pointillés épais)
    \draw[popdark, densely dashed, ultra thick] (Douala) -- (Yaounde) -- ++(2,1.5); % Vers Ngaoundéré
    % ZONE CACAO (Surfacique hachuré orange) - Bassin Sud (Océan, Vallée-du-Ntem, Mvila)
    \pattern[pattern=north east lines, pattern color=poporange!70!popdark] (5.5,1.0) rectangle (8.5,3.5);
    \node[poporange!70!popdark, rotate=30, opacity=0.8] at (7.0, 2.2) {\textbf{Bassin Cacaoyer Sud (~60k t/an)}};
    % USINE KRIBI (Symbole usine)
    \node[draw=popdark, thick, fill=poporange!20, rectangle, minimum width=1cm, minimum height=0.6cm, label=below:\textbf{USINE CACAO}] (Usine) at (7.5, 1.2) {Usine};
    % FLUX EXPORT (Flèche épaisse vers port)
    \node[anchor=west] at (8.5, 0.5) {\textbf{Port Kribi (Eau profonde 16m)}};
    \draw[popgreen, very thick, ->] (Usine.east) -- ++(0.8,0) node[midway, above] {80\% Export};
    % FLUX APPRO (Flèches fines depuis bassin)
    \draw[poporange, ->, thick] (6.5, 2.5) -- (Usine.north west);
    \draw[poporange, ->, thick] (7.0, 3.0) -- (Usine.north);
    % ÉNERGIE (Liaison usine - barrage)
    \draw[popblue, thick, ->] (7.0, 2.2) -- (Usine.north east) node[midway, right, sloped] {Ligne 225 kV};
    % NORD
    \node[draw, circle, fill=popgold, inner sep=2pt] at (10.5, 6.0) {N};
    \draw[->, thick] (10.5, 5.8) -- (10.5, 5.3);
    % ÉCHELLE GRAPHIQUE
    \draw[popdark, thick] (0.5, 0.3) -- (1.5, 0.3) node[midway, below] {50 km};
    \draw[popdark, thick] (0.5, 0.25) -- (0.5, 0.35);
    \draw[popdark, thick] (1.5, 0.25) -- (1.5, 0.35);
\end{tikzpicture}
\legende{Croquis 1 : Localisation de l'unité de transformation cacao - Site de Kribi (Échelle 1/5 000 000). Figuration : Zone d'approvisionnement (hachures orange), Usine (carré orange), Axes routiers (traits pleins), Rail (pointillés), Ligne HT (pointillés bleus), Flux export (flèche verte épaisse), Port (symbole ancre).}
\end{popfigure}

\textbf{Étape 5 : Note d'argumentation (Modèle de rédaction attendue)}

\textbf{Objet :} Note d'argumentation pour l'implantation de l'unité de transformation cacao (50 000 t/an) sur le site de Kribi (Région du Sud).

\textbf{1. Introduction}
La transformation locale du cacao est un pilier de la \cle{SND30} (Objectif : 50\% transformation à l'horizon 2030). Le choix du site d'implantation conditionne la \cle{compétitivité} de la filière, l'\cle{inclusion} des planteurs villageois et la \cle{durabilité} environnementale. L'analyse multicritères de trois sites (Kribi, Douala, Ebolowa) conduit à retenir \textbf{Kribi} comme site optimal.

\textbf{2. Analyse comparative synthétique}
\begin{itemize}
    \item \textbf{Douala} : Avantages indéniables (énergie, main-d'œuvre qualifiée, marché local, rail). \textbf{Inconvénients majeurs} : Saturation portuaire (coûts/logistique), foncier rare/cher/conflit, zone franche saturée, pollution existante. Coût complet logistique export élevé.
    \item \textbf{Ebolowa} : Avantage \textbf{Matière Première} relatif (cœur de bassin Mvila, 25k t/an, circuit court collecte). \textbf{Inconvénients bloquants} : Déficit énergétique critique (réseau 90 kV instable, coût kWh $\times 2$ vs Kribi), enclavement portuaire (120 km route N17, pas de rail), foncier coutumier complexe, main-d'œuvre qualifiée absente. Risque d'usine « île » non compétitive à l'export.
    \item \textbf{Kribi} : \textbf{Meilleur compromis global} (Score 254/320). Énergie abondante/bon marché (Memve'ele, Gaz, RIS fiable), Port eau profonde (Kribi Phase 1, capacité libre, tirant d'eau 16m), Foncier sécurisé/viabilisé (Zone Franche Kribi, titres État), Incitations maximales (Loi 2013, exonérations 10 ans). Défauts (Distance bassin cacao 80-120 km, main-d'œuvre locale limitée) gérables et investissables.
\end{itemize}

\textbf{3. Justification triptyque du choix Kribi}

\textbf{A. Compétitivité-Coût Complet (Levier économique)}
\begin{itemize}
    \item \textbf{Énergie :} Coût kWh industriel \SI{\sim 55}{FCFA} (vs $> \SI{80}{FCFA}$ à Ebolowa, \SI{\sim 50}{FCFA} Douala mais gaz disponible). Stabilité process (50 MW continus) garantie par Memve'ele (211 MW) + RIS + potentiel gaz. \textbf{Gain estimé : 15-20\% sur OPEX énergie vs Ebolowa.}
    \item \textbf{Logistique Export :} Port Kribi à 5 km (vs 300 km Douala depuis Ebolowa, 30 km Douala depuis Douala mais congestion). Évitement surcoûts congestion Douala (délais, surestaries). Accès direct marchés UE/CEMAC/USA. \textbf{Gain logistique export : \SI{\sim 15-25}{USD/t} produit fini.}
    \item \textbf{Foncier/Incitations :} Zone Franche Kribi = Exonération IS (10 ans), TVA, droits douane intrants/équipements. Foncier viabilisé \SI{25}{ha} disponible immédiatement. Réduction CAPEX infrastructure (\SI{\sim 3-5}{Mds FCFA} économisés vs site vierge).
\end{itemize}

\textbf{B. Inclusion Sociale - Planteurs Villageois (Levier territorial)}
\begin{itemize}
    \item \textbf{Collecte organisée :} Rayon 80-120 km (Bassin Sud : Mvila, Océan, Vallée-du-Ntem). Création de \textbf{centres de collecte/fermentation primaire} (CCFP) villageois (coopératives) le long des axes N2/N17/N9. Valorisation qualité (fermentation contrôlée) $\rightarrow$ prime qualité planteurs.
    \item \textbf{Contrats cadre :} Usine $\leftrightarrow$ Coopératives (prix plancher indexé cours mondial, préfinancement intrants). Stabilisation revenus 5 000--10 000 planteurs.
    \item \textbf{Emplois locaux :} 200 directs (formation lycée technique Kribi / IUT Douala délocalisé), 500 indirects (transport, maintenance, services). Priorité recrutement bassins d'approvisionnement. Développement local (routes, électricité, santé) via RSE usine.
\end{itemize}

\textbf{C. Durabilité Environnementale (Levier conformité/acceptabilité)}
\begin{itemize}
    \item \textbf{EIES rigoureuse :} Impacts maîtrisés (effluents : méthanisation coques $\rightarrow$ biogaz/énergie + compost ; eaux : station épuration biologique ; air : filtres manches chaudières biomasse/coques).
    \item \textbf{Énergie verte :} Mix Hydro (Memve'ele) + Biomasse (coques cacao) + Gaz (transition) $\rightarrow$ \textbf{Usine bas-carbone} (objectif \cle{Net Zéro} scope 1\&2 à 10 ans). Certification \cle{ISO 14001} / \cle{Rainforest Alliance} / \cle{ISO 50001}.
    \item \textbf{Biodiversité :} Zone Franche Kribi hors Parc Campo Ma'an (zone tampon respectée). Programme reboisement / corridors écologiques financé par usine. Pas de déforestation directe (site Zone Franche déjà dégagé/viabilisé).
\end{itemize}

\textbf{4. Conclusion}
Le site de \textbf{Kribi} présente la \textbf{meilleure adéquation} aux exigences du cahier des charges (Énergie 50 MW, Export 80\%, Foncier 25 ha, Incitations) et aux objectifs stratégiques nationaux (SND30, Transformation locale, ZES Kribi). Il transforme la contrainte « distance bassin cacao » en opportunité d'\cle{organisation de la filière amont} (collecte/qualité/contrats) et résout les goulots d'étranglement énergie/logistique/foncier qui pénalisent Douala et Ebolowa. L'investissement initial dans la logistique amont (camions, CCFP) est marginal ($< \SI{5}{\%}$ CAPEX total) face aux gains structurels OPEX (énergie, port, fisc) sur 20 ans. \textbf{Kribi est le choix rationnel, durable et inclusif.}

\textbf{Étape 6 : Points clés débat oral (Synthèse professeur)}
\begin{itemize}
    \item \textbf{Argument « Douala » :} « Usine là où est le marché/la main-d'œuvre ». \textit{Réponse :} Usine exportatrice (80\%), le marché est le port. Douala = goulot portuaire + foncier. Délocaliser l'énergie (gaz) à Kribi est plus facile que délocaliser le port.
    \item \textbf{Argument « Ebolowa/Bertoua » :} « Usine au plus près de la matière première (Weber) ». \textit{Réponse :} Indice de perte de poids cacao faible (fèves sèches stables). Facteur énergie (électro-intensif) + Facteur export (conteneur) dominants. « Usine dans la brousse » = coûts énergie/logistique prohibitifs, échec compétitivité.
    \item \textbf{Leviers d'action publique :} Réhabilitation N17 (Ebolowa-Kribi), Ligne HT Yaoundé-Memve'ele-Kribi (renforcement), Formation professionnelle ciblée Kribi, Appui coopératives certification.
\end{itemize}
\end{exoresolu}

\courssub{Bilan de l'activité}
Cette activité d'intégration a permis de mettre en évidence :
\begin{itemize}
    \item La \textbf{primauté des facteurs « Énergie » et « Logistique d'export »} sur le seul facteur « Matière première » pour une industrie electro-intensive exportatrice (cacao), validant la théorie de la localisation adaptée au contexte camerounais (poids du port, fiabilité réseau).
    \item Le rôle structurant des \textbf{Zones Économiques Spéciales (ZES/Kribi)} et de la \textbf{Loi 2013} comme leviers d'attractivité (foncier, fiscalité, énergie) pour rééquilibrer l'espace productif hors Douala.
    \item La nécessité d'une \textbf{vision systémique} : l'usine n'est pas un point isolé mais le nœud d'un \textbf{système productif localisé} (bassin cacao $\rightarrow$ collecte/qualité $\rightarrow$ usine $\rightarrow$ port $\rightarrow$ marché mondial) nécessitant des infrastructures de liaison (route, rail, énergie) et une gouvernance filière (contrats, certification).
    \item La faisabilité de l'\textbf{inclusion sociale} (planteurs villageois) par l'organisation de l'amont (CCFP, contrats, prime qualité) plutôt que par la simple proximité géographique de l'usine.
    \item L'exigence de \textbf{durabilité environnementale} (EIES, énergie verte, économie circulaire coques) comme condition d'acceptabilité (populations, Parc Campo Ma'an, marchés UE \cle{CSRD}/\cle{CSDDD}) et de compétitivité long terme.
    \item La maîtrise par les élèves de la \textbf{sémiologie graphique} (croquis 1/5 000 000) et de l'\textbf{argumentation écrite/orale} structurée, compétences clés pour l'examen du \cle{Probatoire/Baccalauréat Technique}.
\end{itemize}

\newpage

\coursec{Méthodes et exercices résolus}

\courssub{Savoir-faire 1 : Appliquer les notions de l'unité à une situation concrète (étude de cas d'une firme agro-industrielle)}

\begin{methode}[Analyser une firme agro-industrielle camerounaise]
Pour traiter une étude de cas sur une grande plantation (CDC, SOCAPALM, PHP, HEVECAM, SOSUCAM), suivre une démarche en 5 étapes :
\begin{enumerate}
\item \textbf{Identifier} la firme : nom, localisation, date de création, capital (public, privé, mixte, étranger).
\item \textbf{Décrire la spéculation} : produit (banane, huile de palme, hévéa, sucre), surfaces, tonnages, destination (exportation ou marché intérieur).
\item \textbf{Analyser l'organisation spatiale} : plantations en blocs, usines de transformation sur place, voies d'évacuation (routes, chemin de fer, port), villages de travailleurs.
\item \textbf{Dégager les effets} : emplois créés, devises rapportées, infrastructures, mais aussi problèmes fonciers et environnementaux.
\item \textbf{Conclure} en reliant la firme au dualisme de l'économie camerounaise (secteur moderne / secteur traditionnel).
\end{enumerate}
\textbf{Réflexes} : toujours citer au moins un chiffre (surface, tonnage, effectif) ; toujours nommer la région et le port d'exportation concerné.
\end{methode}

\begin{definition}[Agriculture de plantation]
L'\cle{agriculture de plantation} (ou agro-industrie) est une forme d'exploitation agricole moderne caractérisée par : de \textbf{grandes surfaces} (plusieurs milliers d'hectares), une \textbf{monoculture} de rente (banane, palme, hévéa, sucre), d'importants \textbf{capitaux} (publics ou privés, souvent étrangers), une \textbf{main-d'œuvre salariée} nombreuse, une \textbf{usine de transformation} intégrée et une production destinée essentiellement à l'\textbf{exportation}. Elle s'oppose à l'agriculture vivrière familiale, pratiquée par les petits paysans pour l'autoconsommation.
\end{definition}

\begin{exoresolu}[La CDC, firme agro-industrielle du Sud-Ouest]
Énoncé : La Cameroon Development Corporation (CDC), créée en 1947, est l'une des plus grandes entreprises agro-industrielles du Cameroun. Elle exploite des plantations de bananes, d'hévéas et de palmiers à huile dans la région du Sud-Ouest, autour de Tiko, Buea et dans la plaine de Mungo. Elle emploie plusieurs milliers de travailleurs et exporte l'essentiel de sa production par le port de Douala.
\begin{enumerate}
\item Présentez la CDC (statut, localisation, spéculations).
\item Montrez comment la CDC illustre l'organisation spatiale d'une agriculture de plantation.
\item Dégagez deux effets positifs et une limite de ce type d'exploitation.
\end{enumerate}
\tcblower
\textbf{1. Présentation de la CDC.}
La CDC est une entreprise à capital majoritairement public (État camerounais), héritière des plantations allemandes du Kamerun, confisquées après la Première Guerre mondiale puis regroupées en 1947. Elle est localisée dans la région du Sud-Ouest (Tiko, Buea, plaine de Mungo), zone de forte pluviométrie et de sols volcaniques fertiles (pentes du mont Cameroun), favorables aux cultures pérennes. Ses trois spéculations sont : la \cle{banane} d'exportation, l'\cle{hévéa} (caoutchouc) et le \cle{palmier à huile}.

\textbf{2. Organisation spatiale.}
La CDC illustre le modèle classique de la plantation : de grands \textbf{blocs de cultures} monospécifiques de plusieurs milliers d'hectares ; des \textbf{usines de transformation} installées au cœur des plantations (huileries, stations de conditionnement des bananes, usines de caoutchouc) ; des \textbf{villages-camps} construits pour loger les travailleurs ; un réseau de \textbf{pistes internes} relié aux routes nationales et au \textbf{port de Douala}, porte de sortie des exportations. L'espace est donc entièrement organisé autour de la production et de l'évacuation du produit vers l'extérieur.

\textbf{3. Effets et limites.}
Effets positifs : (a) création de milliers d'emplois salariés dans une économie dominée par l'agriculture vivrière ; (b) rentrées de devises grâce aux exportations de bananes, de caoutchouc et d'huile de palme, et construction d'infrastructures (routes, dispensaires, écoles). Limite : la plantation concentre de vastes terres (problème foncier pour les communautés locales) et repose sur la monoculture, vulnérable aux fluctuations des cours mondiaux et source d'appauvrissement des sols.

\textbf{Conclusion} : la CDC est un cas typique d'agriculture moderne de plantation, intégrée au marché mondial, qui coexiste avec l'agriculture vivrière traditionnelle : elle illustre le \cle{dualisme} de l'économie camerounaise.
\end{exoresolu}

\courssub{Savoir-faire 2 : Appliquer les notions à la localisation industrielle}

\begin{methode}[Expliquer la localisation d'une industrie]
Pour expliquer pourquoi une industrie s'installe à un endroit donné, mobiliser les \cle{facteurs de localisation} dans l'ordre suivant :
\begin{enumerate}
\item \textbf{Matières premières} : proximité des ressources (bauxite, bois, pétrole, produits agricoles).
\item \textbf{Énergie} : disponibilité d'électricité (barrages hydroélectriques : Songloulou et Édéa sur la Sanaga, Lagdo sur la Bénoué).
\item \textbf{Main-d'œuvre et marché} : population nombreuse et pouvoir d'achat (Douala, Yaoundé).
\item \textbf{Transports} : port, chemin de fer, routes (axe Douala--Yaoundé, Transcamerounais).
\item \textbf{Politique de l'État} : zones industrielles, incitations, capitale administrative.
\end{enumerate}
\textbf{Formule à retenir} : une localisation industrielle = combinaison de plusieurs facteurs, jamais un seul. Toujours hiérarchiser (facteur dominant + facteurs secondaires).
\end{methode}

\begin{exoresolu}[Pourquoi Douala est-il le premier pôle industriel du Cameroun ?]
Énoncé : Douala concentre la majorité des unités industrielles du Cameroun (brasseries, cimenteries, agroalimentaire, première transformation du bois). À l'aide des facteurs de localisation, expliquez cette concentration industrielle, puis citez un autre pôle industriel camerounais et sa spécialité.
\tcblower
\textbf{Raisonnement} : il faut appliquer la grille des facteurs de localisation au cas de Douala.

\textbf{1. Le port} : Douala est le premier port du pays, par où transitent l'essentiel des importations (matières premières, biens d'équipement) et des exportations. Les industries d'import-substitution et d'exportation s'y installent pour réduire les coûts de transport.

\textbf{2. Le marché et la main-d'œuvre} : Douala est la plus grande ville du Cameroun (plus de 3 millions d'habitants) : elle offre une main-d'œuvre abondante et un marché de consommation solvable pour les brasseries, l'agroalimentaire, les cimenteries.

\textbf{3. L'énergie et les matières premières} : la ville est proche des barrages hydroélectriques du bassin de la Sanaga (Songloulou, Édéa) qui fournissent l'électricité, et elle reçoit les produits agricoles de l'arrière-pays (cacao, café, bananes, coton, bois) qu'elle transforme avant exportation.

\textbf{4. Les infrastructures} : nœud ferroviaire (ligne Douala--Yaoundé--Ngaoundéré, le Transcamerounais) et routier, zone industrielle de Bonabéri, présence des banques et des services aux entreprises.

\textbf{Autre pôle} : Édéa, spécialisé dans l'électrométallurgie (aluminium, usine ALUCAM) grâce à l'électricité abondante et bon marché du barrage d'Édéa ; Limbé avec le raffinage du pétrole (SONARA) ; ou Yaoundé (industries de biens de consommation, fonctions de capitale administrative).

\textbf{Conclusion} : Douala doit sa primauté industrielle à la combinaison d'un port, d'un grand marché, d'énergie et d'infrastructures ; c'est le facteur \textbf{portuaire} qui joue le rôle dominant.
\end{exoresolu}

\courssub{Savoir-faire 3 : Traiter et interpréter des données chiffrées (calculs de parts et d'évolutions)}

\begin{methode}[Calculer une part en pourcentage et un taux de variation]
Deux calculs sont exigibles dans les sujets portant sur l'économie moderne :
\begin{enumerate}
\item \textbf{Part en pourcentage} : $\text{part} = \dfrac{\text{valeur de la partie}}{\text{valeur du total}} \times 100$. Toujours vérifier que la somme des parts fait environ $100\,\%$.
\item \textbf{Taux de variation} entre deux dates : $t = \dfrac{V_f - V_i}{V_i} \times 100$, où $V_i$ est la valeur initiale et $V_f$ la valeur finale. Un résultat positif = hausse ; négatif = baisse.
\end{enumerate}
\textbf{Réflexes} : arrondir à $0,1\,\%$ près ; écrire les unités ; conclure par une phrase d'interprétation géographique, pas seulement par le chiffre.
\end{methode}

\begin{exoresolu}[Le poids du secteur secondaire dans l'économie camerounaise]
Énoncé : En valeur ajoutée, on estime que le PIB du Cameroun se répartit approximativement ainsi : primaire $15\,\%$, secondaire $26\,\%$, tertiaire $59\,\%$. La production industrielle manufacturière est passée, dans une région donnée, de 400 milliards de FCFA à 500 milliards de FCFA entre deux exercices.
\begin{enumerate}
\item Calculez la part cumulée des secteurs primaire et secondaire. Que représente le tertiaire ?
\item Calculez le taux de variation de la production manufacturière. Interprétez.
\end{enumerate}
\tcblower
\textbf{1. Part cumulée primaire + secondaire :}
\[ 15 + 26 = 41\,\% \]
Le tertiaire représente $59\,\%$ du PIB, soit plus de la moitié. Vérification : $15 + 26 + 59 = 100\,\%$ : la répartition est cohérente.

\textbf{2. Taux de variation de la production manufacturière :}
\begin{align*}
t &= \dfrac{V_f - V_i}{V_i} \times 100 \\
t &= \dfrac{500 - 400}{400} \times 100 \\
t &= \dfrac{100}{400} \times 100 = 25\,\%
\end{align*}
\textbf{Interprétation} : la production manufacturière a augmenté de $25\,\%$ entre les deux exercices, ce qui traduit un dynamisme du secteur secondaire. Toutefois, avec seulement $26\,\%$ du PIB, le secondaire reste faible comparé au tertiaire : l'industrie camerounaise progresse mais ne joue pas encore le rôle moteur attendu dans la trajectoire d'émergence.

\textbf{Conclusion} : les calculs confirment une économie camerounaise dominée par le tertiaire, dont l'industrie, bien qu'en croissance ($+25\,\%$), demeure un secteur secondaire limité.
\end{exoresolu}

\courssub{Savoir-faire 4 : Rédiger et justifier une démarche (croquis de synthèse)}

\begin{methode}[Construire un croquis « Les pôles de l'économie moderne au Cameroun »]
\begin{enumerate}
\item \textbf{Tracer le fond de carte} : contour simplifié du Cameroun en polygone, flèche du nord, titre en haut.
\item \textbf{Placer les villes-pôles} par des points nommés : Douala, Yaoundé, Édéa, Ngaoundéré, Limbé, Garoua.
\item \textbf{Figurer les plantations} par des zones colorées : Sud-Ouest (CDC : banane, palme), Littoral (SOCAPALM, PHP), Centre-Sud (HEVECAM, SOSUCAM).
\item \textbf{Figurer l'industrie} : symboles pour Douala (port + industries), Édéa (aluminium), Limbé (raffinage), et les barrages (Songloulou, Édéa, Lagdo).
\item \textbf{Tracer les axes} : chemin de fer Transcamerounais (Douala--Yaoundé--Ngaoundéré), routes, et des flèches d'exportation vers le port de Douala.
\item \textbf{Rédiger la légende numérotée} : chaque symbole doit y figurer, dans l'ordre.
\end{enumerate}
\textbf{Réflexes} : un croquis sans titre, sans nord ou sans légende est pénalisé ; les flèches de flux doivent être orientées vers le port.
\end{methode}

\begin{exoresolu}[Croquis : les pôles de l'économie moderne au Cameroun]
Énoncé : Réalisez un croquis montrant les pôles de l'économie moderne au Cameroun (grandes plantations, pôles industriels, axes d'évacuation), accompagné d'une légende organisée.
\tcblower
\textbf{Démarche} : fond de carte simplifié, puis report des éléments du plus important au plus secondaire.

\begin{popfigure}
\begin{tikzpicture}[scale=1.0]
% Contour simplifié du Cameroun
\draw[thick,popdark,fill=popcream] (0,0) -- (2.2,-0.2) -- (3.4,0.6) -- (3.6,1.8) -- (4.6,2.6) -- (4.2,3.6) -- (3.2,4.0) -- (2.6,5.4) -- (1.8,6.2) -- (1.2,5.2) -- (0.4,4.6) -- (0.8,3.4) -- (0.2,2.4) -- (-0.4,1.2) -- cycle;
% Nord
\draw[-{Stealth},thick,popdark] (5.2,5.4) -- (5.2,6.1) node[above]{\textbf{N}};
% Zones de plantations
\fill[popgreen!45] (0.5,1.6) circle (0.45);
\fill[popgreen!45] (1.6,1.0) circle (0.4);
\fill[popgreen!45] (2.6,0.6) circle (0.4);
% Villes
\fill[popink] (0.6,1.3) circle (0.09) node[left=1pt,font=\small]{Douala};
\fill[popink] (2.2,1.9) circle (0.09) node[right=1pt,font=\small]{Yaoundé};
\fill[popink] (1.1,1.7) circle (0.07) node[above=1pt,font=\small]{Édéa};
\fill[popink] (2.6,4.6) circle (0.09) node[right=1pt,font=\small]{Ngaoundéré};
\fill[popink] (0.1,2.2) circle (0.07) node[left=1pt,font=\small]{Limbé};
\fill[popink] (2.4,5.6) circle (0.07) node[right=1pt,font=\small]{Garoua};
% Chemin de fer
\draw[very thick,poporange,dashed] (0.6,1.3) -- (2.2,1.9) -- (2.6,4.6);
% Flèche exportation vers Douala
\draw[-{Stealth},very thick,popblue] (2.9,0.5) .. controls (1.8,0.2) .. (0.75,1.15);
% Barrages
\node[font=\small,popblue] at (1.7,2.3) {$\blacksquare$ Songloulou};
\node[font=\small,popblue] at (2.9,5.2) {$\blacksquare$ Lagdo};
\end{tikzpicture}
\legende{Croquis schématique (échelle indicative). 1 : zones de grandes plantations (CDC au Sud-Ouest ; SOCAPALM/PHP au Littoral ; HEVECAM/SOSUCAM au Centre-Sud, en vert). 2 : pôles industriels et urbains (points noirs). 3 : chemin de fer Transcamerounais (pointillés orange). 4 : flux d'exportation vers le port de Douala (flèche bleue). 5 : barrages hydroélectriques (carrés bleus).}
\end{popfigure}

\textbf{Justification de la démarche} : le choix d'un fond simplifié permet de hiérarchiser l'information ; les plantations sont localisées dans le Sud forestier humide (conditions agro-climatiques), les industries dans les villes (facteurs de localisation), et le flux converge vers Douala, ce qui traduit la dépendance de l'économie moderne au port.

\textbf{Conclusion} : le croquis montre une économie moderne concentrée dans le Sud et sur l'axe Douala--Yaoundé--Ngaoundéré, signe d'un déséquilibre spatial Nord/Sud.
\end{exoresolu}

\courssub{Savoir-faire 5 : Rédiger et justifier une conclusion (sujet de synthèse)}

\begin{methode}[Rédiger un paragraphe argumenté et une conclusion au Bac]
\begin{enumerate}
\item \textbf{Annoncer} l'idée principale dans une phrase d'introduction claire (reformulation de la consigne).
\item \textbf{Développer} : une idée = un argument + un exemple précis (firme, chiffre, lieu) + une explication.
\item \textbf{Articuler} avec des connecteurs : d'abord, ensuite, en effet, cependant, enfin.
\item \textbf{Nuancer} : présenter les limites ou les contre-arguments (le correcteur valorise l'esprit critique).
\item \textbf{Conclure} en répondant exactement à la question posée, puis ouvrir brièvement (ex. : émergence à l'horizon 2035).
\end{enumerate}
\textbf{À retenir} : jamais d'exemple sans explication, jamais d'affirmation sans preuve chiffrée ou localisée.
\end{methode}

\begin{exoresolu}[L'économie moderne camerounaise : moteur ou façade du développement ?]
Énoncé : En un paragraphe argumenté d'une quinzaine de lignes, montrez que l'économie moderne (grandes plantations, industrie) est à la fois un atout et un facteur de déséquilibres pour le Cameroun.
\tcblower
\textbf{Proposition de rédaction :}

L'économie moderne camerounaise, fondée sur les grandes plantations et l'industrie, constitue un atout réel mais demeure source de déséquilibres. D'abord, les agro-industries comme la CDC (Sud-Ouest), la SOCAPALM (Littoral) ou la SOSUCAM (Centre) fournissent des milliers d'emplois salariés et rapportent des devises grâce aux exportations de bananes, d'huile de palme, de caoutchouc et de sucre. Ensuite, l'industrie, concentrée à Douala (brasseries, cimenteries, agroalimentaire), à Édéa (aluminium d'ALUCAM grâce au barrage hydroélectrique) et à Limbé (raffinage de la SONARA), participe à environ $26\,\%$ du PIB et réduit la dépendance aux importations. Cependant, ce secteur moderne reste une façade fragile : il est contrôlé en partie par des capitaux étrangers, il repose sur la monoculture vulnérable aux cours mondiaux, et il se concentre dans le Sud et le long de l'axe Douala--Yaoundé--Ngaoundéré, marginalisant le Grand Nord. Enfin, les plantations posent des problèmes fonciers (expropriation des communautés) et environnementaux (déforestation, appauvrissement des sols). En conclusion, l'économie moderne est un moteur indispensable mais insuffisant : pour devenir un véritable levier d'émergence à l'horizon 2035, le Cameroun doit transformer davantage ses matières premières sur place et rééquilibrer l'organisation spatiale de ses activités.

\textbf{Justification de la démarche} : le paragraphe suit le plan annoncé (atouts $\rightarrow$ limites $\rightarrow$ conclusion nuancée), chaque argument est étayé par un exemple localisé ou un chiffre, et la conclusion répond précisément à la question posée.
\end{exoresolu}

\courssub{Savoir-faire 6 : Distinguer et comparer les statuts fonciers des plantations}

\begin{methode}[Construire un tableau comparatif rigoureux]
Pour comparer les types d'exploitations agricoles (plantations publiques, privées, villageoises) :
\begin{enumerate}
\item \textbf{Choisir les critères} de comparaison : statut du capital, taille, produit, main-d'œuvre, destination de la production.
\item \textbf{Remplir chaque case} avec une information courte et exacte (pas de phrase longue).
\item \textbf{Dégager une conclusion} : ce qui rapproche et ce qui oppose les types.
\end{enumerate}
\textbf{Réflexe} : un tableau comparatif doit toujours être suivi d'une phrase de synthèse rédigée.
\end{methode}

\begin{exoresolu}[Plantations publiques et privées au Cameroun]
Énoncé : Construisez un tableau comparant les plantations publiques (ex. CDC) et privées (ex. SOCAPALM, PHP) selon quatre critères, puis tirez une conclusion.
\tcblower
\begin{center}
\begin{tabularx}{\linewidth}{|l|X|X|}
\hline
\rowcolor{popblueL} \textbf{Critère} & \textbf{Plantation publique (CDC)} & \textbf{Plantation privée (SOCAPALM, PHP)} \\
\hline
Capital & Majoritairement État camerounais & Capitaux privés, souvent mixtes ou étrangers \\
\hline
Localisation & Sud-Ouest (Tiko, Mungo) & Littoral (Apouh, Dibombari, Mbongo, Penja) \\
\hline
Produits & Banane, hévéa, palme à huile & Huile de palme, hévéa, banane (PHP) \\
\hline
Destination & Exportation + marché intérieur & Exportation + marché intérieur (huile) \\
\hline
\end{tabularx}
\end{center}
\textbf{Conclusion} : publiques ou privées, les plantations partagent la même logique (monoculture de rente orientée vers l'exportation, emploi salarié) ; elles se distinguent surtout par la nature du capital, ce qui illustre la place croissante des investisseurs privés dans l'agro-industrie camerounaise.
\end{exoresolu}

\begin{attention}[Les 5 erreurs qui coûtent des points au Bac]
\begin{enumerate}
\item \textbf{Confondre agriculture de plantation et agriculture vivrière} : la plantation est une exploitation de grande taille, monospécifique, capitalistique, orientée vers l'exportation ; l'agriculture vivrière est familiale et destinée à l'autoconsommation. Ne jamais mélanger les deux dans une définition.
\item \textbf{Citer une firme sans la localiser} : écrire « la SOCAPALM » sans préciser « Littoral » ou « CDC » sans « Sud-Ouest » fait perdre le point de l'exemple. Toute firme citée doit être localisée.
\item \textbf{Oublier les éléments obligatoires du croquis} : titre, flèche du nord, légende numérotée. Un croquis sans légende est sanctionné, même s'il est exact.
\item \textbf{Se tromper dans le taux de variation} : diviser par la valeur \emph{finale} au lieu de la valeur \emph{initiale}. Rappel : $t = \dfrac{V_f - V_i}{V_i}\times 100$. Vérifier le signe (hausse $>0$, baisse $<0$).
\item \textbf{Affirmer sans chiffrer ni nuancer} : écrire « l'industrie domine l'économie » est faux (le secondaire ne représente qu'environ $26\,\%$ du PIB, derrière le tertiaire). Toute affirmation doit être étayée par un chiffre et nuancée dans la conclusion.
\end{enumerate}
\end{attention}

\newpage

\coursec{Exercices}

\courssub{Je vérifie mes connaissances}

\begin{exercice}[QCM : Dualisme et poids de l'économie moderne]
Parmi les affirmations suivantes, cochez celle qui caractérise le mieux le dualisme de l'économie camerounaise.
\begin{enumerate}
    \item Coexistence harmonieuse entre une agriculture vivrière de subsistance et une agriculture de plantation d'exportation.
    \item Opposition structurelle entre un secteur moderne, capitaliste, tourné vers l'exportation (plantations, industrie) et un secteur traditionnel, paysan, vivrier.
    \item Part égale du secteur primaire et du secteur secondaire dans le PIB national.
    \item Absence de liens financiers et commerciaux entre les firmes agro-industrielles et les planteurs villageois.
\end{enumerate}
\end{exercice}

\begin{exercice}[Vrai / Faux justifié : Statuts fonciers et acteurs]
Indiquez si les affirmations sont vraies (V) ou fausses (F) et justifiez brièvement votre réponse en vous appuyant sur les ordonnances foncières de 1974 (domaine national).
\begin{enumerate}
    \item Les terres occupées par les grandes plantations (CDC, SOCAPALM, HEVECAM) relèvent du domaine national, sur lequel l'État consent des baux de longue durée.
    \item Les « outgrowers » (planteurs sous contrat) sont des salariés de la firme agro-industrielle.
    \item Le bail consenti à une firme agro-industrielle peut porter sur une durée allant jusqu'à 99 ans.
    \item Lors de l'extension d'une plantation sur le domaine national, l'indemnisation porte sur les cultures et les constructions (biens meubles et immeubles par incorporation), et non sur le sol nu, propriété de l'État.
\end{enumerate}
\end{exercice}

\begin{exercice}[Définitions et calcul direct : Indicateurs industriels]
\begin{enumerate}
    \item Définissez le \cle{taux de transformation locale} (TTL) pour une filière agricole donnée.
    \item Une usine de transformation de cacao a réceptionné 85\,000 tonnes de fèves brutes sur l'année. Elle a produit 62\,000 tonnes de produits dérivés (pâte, beurre, poudre). Calculez le taux de transformation locale de cette usine. Exprimez le résultat en pourcentage arrondi au dixième.
    \item Citez la principale certification internationale exigée pour l'huile de palme durable sur le marché européen.
\end{enumerate}
\end{exercice}

\begin{exercice}[Questions courtes : Localisation et corridors industriels]
\begin{enumerate}
    \item Citez \textbf{trois} facteurs géographiques qui ont déterminé l'implantation de l'usine ALUCAM à Edéa.
    \item Quel axe urbain concentre l'essentiel (environ 70\,\%) de la valeur ajoutée industrielle du Cameroun ? Citez les villes pôles qui le structurent.
    \item Différenciez un \cle{planteur villageois indépendant} d'un \cle{outgrower} (planteur sous contrat) sur le plan de la commercialisation de la production.
\end{enumerate}
\end{exercice}

\courssub{Je m'entraîne}

\begin{exercice}[Analyse de document : Évolution de l'agro-industrie dans les exportations]
\textbf{Document :} Extrait simulé d'un rapport INS/BCEC « Structure des exportations non pétrolières 2010--2023 ».
\begin{center}
\begin{tabular}{|l|c|c|c|}\hline
\textbf{Produit} & \textbf{Part 2010 (\%)} & \textbf{Part 2023 (\%)} & \textbf{Variation (points)} \\ \hline
Bois et articles en bois & 18,2 & 12,5 & $-5{,}7$ \\ \hline
Cacao (fèves + dérivés) & 22,5 & 28,3 & $+5{,}8$ \\ \hline
Coton (fibre + graine) & 4,1 & 3,2 & $-0{,}9$ \\ \hline
Caoutchouc naturel & 3,8 & 5,6 & $+1{,}8$ \\ \hline
Huile de palme & 1,2 & 4,1 & $+2{,}9$ \\ \hline
Banane / Plantain & 6,5 & 7,8 & $+1{,}3$ \\ \hline
Autres (café, poivre, etc.) & 43,7 & 38,5 & $-5{,}2$ \\ \hline
\end{tabular}
\end{center}
\legende{Source : INS / BCEC, données simulées à des fins pédagogiques. Parts en \% des exportations non pétrolières.}
\begin{enumerate}
    \item Identifiez les deux filières agro-industrielles dont la part a le plus progressé entre 2010 et 2023. Calculez leur part cumulée en 2023.
    \item La filière bois régresse en part relative. Proposez \textbf{deux} causes structurelles liées à la politique de transformation locale (loi de 1994, restriction des exportations de grumes brutes).
    \item En vous appuyant sur le tableau et vos connaissances, proposez \textbf{deux} mesures concrètes de politique économique pour augmenter la part de l'agro-industrie dans les exportations non pétrolières à l'horizon 2030 (PND30).
\end{enumerate}
\end{exercice}

\begin{exercice}[Étude de cas : Conflit foncier CDC / Communautés riveraines]
\textbf{Contexte :} La Cameroon Development Corporation (CDC), créée en 1947, exploite de grandes plantations de palmiers à huile, d'hévéas et de bananiers dans la région du Sud-Ouest (départements du Fako et du Kupe-Muanenguba, autour de Tiko, Muyuka et Mbonge), sur des terres du domaine national concédées par bail de longue durée. Des communautés riveraines contestent l'assiette de la plantation, dénonçant l'absence de consultation préalable, la non-reconnaissance des droits coutumiers sur les terres de jachère et l'insuffisance des retombées au titre de la Responsabilité Sociétale des Entreprises (RSE). Des tensions ont éclaté lors d'une tentative d'extension sur des terres revendiquées par les villages.
\begin{enumerate}
    \item Identifiez les \textbf{trois} catégories d'acteurs principaux en présence et précisez le fondement juridique de leur position (droit écrit, droit coutumier, bail).
    \item Analysez les arguments juridiques de la CDC (bail, ordonnances foncières de 1974) et les arguments coutumiers des communautés (occupation ancestrale, jachère assimilée à propriété).
    \item Proposez une issue négociée structurée en \textbf{trois} volets : (1) Médiation / instance de dialogue, (2) Ajustement de l'assiette (réduction, zonage), (3) Mécanisme de partage de la valeur (fonds de développement local, redevances surfaciques majorées, embauche locale).
\end{enumerate}
\end{exercice}

\begin{exercice}[Croquis de synthèse guidé : Pôles de l'économie moderne]
À l'aide d'une carte muette du Cameroun fournie en annexe (ou tracée au tableau), réalisez un croquis légendé intitulé \emph{« Les pôles de l'économie moderne au Cameroun »}. Respectez impérativement la légende imposée ci-dessous.
\textbf{Légende imposée :}
\begin{itemize}
    \item \textbf{Plantations (point + sigle + spéculation) :} CDC (palmier à huile / hévéa / banane) -- Tiko, Muyuka (Sud-Ouest) ; SOCAPALM (palmier à huile) -- Dibombari, Mbongo, Eséka, Kribi ; HEVECAM (hévéa) -- Nieté (près de Kribi) ; SAFACAM (palmier / hévéa) -- Dizangué ; PHP (banane) -- Penja.
    \item \textbf{Unités industrielles (carré + nom + type) :} ALUCAM (aluminium) -- Edéa ; SONARA (raffinage) -- Limbé ; CIMENCAM (ciment) -- Douala/Figuil ; SOSUCAM (sucre) -- Nkoteng/Mbandjock.
    \item \textbf{Énergie (triangle + nom) :} barrages d'Edéa et Song Loulou (sur la Sanaga), Lom Pangar (sur le Lom) ; pipeline Tchad--Cameroun (aboutissant à Kribi) -- trait discontinu.
    \item \textbf{Ports (ancre) :} Douala (port autonome), Kribi (port en eau profonde).
    \item \textbf{Zone franche (hachures) :} zone franche industrielle de Kribi.
    \item \textbf{Corridor industriel (flèche épaisse) :} Douala -- Edéa -- Yaoundé -- Kribi.
\end{itemize}
\textbf{Consignes graphiques :} orientation (flèche du Nord), échelle approximative, titre, légende organisée par thèmes, propreté.
\end{exercice}

\begin{exercice}[Dossier documentaire : Filère bois et transformation locale]
\textbf{Document 1 :} Extrait de la Loi n°94/01 du 20 janvier 1994 portant régime des forêts : l'exploitation forestière doit s'accompagner d'une transformation locale ; l'exportation des grumes brutes est progressivement restreinte afin d'atteindre un taux de transformation locale élevé (objectif de 70\,\%).
\textbf{Document 2 :} Taux de transformation locale effectif déclaré (MINFOF, données simulées) : 2015 : 42\,\% ; 2018 : 55\,\% ; 2021 : 63\,\% ; 2023 : 68\,\%.
\textbf{Document 3 :} Principaux obstacles identifiés par les industriels du bois : (a) coût de l'énergie (électricité instable, prix du kWh industriel élevé) ; (b) vétusté des machines et manque d'accès au crédit long terme ; (c) concurrence du bois scié artisanal non certifié ; (d) enclavement des sites de production (pistes impraticables en saison des pluies).
\begin{enumerate}
    \item Calculez le taux de croissance moyen annuel du TTL entre 2015 et 2023. Le Cameroun est-il en voie d'atteindre l'objectif de 70\,\% ?
    \item Classez les quatre obstacles du Document 3 en \textbf{facteurs internes} à l'entreprise et \textbf{facteurs externes} (environnement des affaires).
    \item Rédigez un paragraphe argumenté (10--12 lignes) expliquant pourquoi la transformation locale du bois est un levier majeur d'industrialisation (emplois, valeur ajoutée, diversification des exportations) mais bute sur la gouvernance forestière et le coût des facteurs de production.
\end{enumerate}
\end{exercice}

\courssub{Je me prépare au Bac}

\begin{exercice}[Sujet type Bac -- Dissertation de géographie : Aménagement du territoire]
\textbf{Sujet :} \emph{« À l'aide d'exemples précis tirés de l'agriculture de grandes plantations et de l'industrie, montrez comment l'économie moderne participe à l'aménagement du territoire camerounais. »}
\textbf{Consignes :} Vous traiterez ce sujet sous la forme d'une dissertation structurée (introduction, développement en deux ou trois parties équilibrées, conclusion). Vous ferez apparaître obligatoirement les notions de \cle{dualisme}, \cle{polarisation} (Douala/Edéa/Yaoundé/Kribi), \cle{enclavement} (Nord, Est), \cle{infrastructures} (routes, rail, ports, énergie), \cle{conflits fonciers} et la trajectoire d'\cle{émergence} (PND30). Un croquis localisé dans le développement est vivement recommandé.
\end{exercice}

\begin{exercice}[Sujet type Bac -- Exercice cartographique commenté]
\textbf{Sujet :} \emph{« À l'aide de la carte fournie (carte muette du Cameroun avec fond hydrographique et villes principales), réalisez le croquis suivant : "Les infrastructures et flux de l'économie moderne au Cameroun". Vous y ferez figurer et nommerez : 5 plantations (spéculation + firme), 4 unités industrielles (type + nom), 3 barrages hydroélectriques, 2 ports, 1 pipeline, 1 zone franche. Vous tracerez les principaux flux d'exportation (flèches proportionnelles) et le corridor industriel majeur. Vous rédigerez une légende organisée et un commentaire de 15 lignes analysant la disparité spatiale des équipements. »}
\textbf{Données pour le croquis (à placer) :}
\begin{itemize}
    \item Plantations : CDC (banane / palmier / hévéa) -- Sud-Ouest ; SOCAPALM (palmier) -- Littoral/Sud ; HEVECAM (hévéa) -- Sud ; SAFACAM (palmier / hévéa) -- Dizangué ; PHP (banane) -- Penja.
    \item Industries : ALUCAM (aluminium) -- Edéa ; SONARA (raffinage) -- Limbé ; CIMENCAM (ciment) -- Douala/Figuil ; SOSUCAM (sucre) -- Centre.
    \item Énergie : barrages d'Edéa, Song Loulou, Lom Pangar ; pipeline Tchad--Kribi.
    \item Ports : Douala, Kribi.
    \item Zone franche : Kribi.
    \item Flux : évacuation des productions (bois, cacao, coton, aluminium, pétrole) vers Douala et Kribi.
\end{itemize}
\end{exercice}

\begin{exercice}[Sujet type Bac -- Dissertation technique : Industrialisation par la transformation]
\textbf{Sujet :} \emph{« L'industrialisation du Cameroun par la transformation locale des matières premières agricoles et forestières : atouts, contraintes et leviers d'action (PND30). »}
\textbf{Consignes :} Votre démonstration s'appuiera sur des exemples concrets (cacao, bois, coton, huile de palme, caoutchouc). Vous structurerez votre réponse en trois parties : I. Les atouts structurels (ressources, marché régional CEMAC, main-d'œuvre, PND30) ; II. Les contraintes persistantes (énergie, financement, logistique, gouvernance, bas TTL, secteur informel) ; III. Les leviers d'action prioritaires du PND30 (zones économiques spéciales, incitations fiscales et douanières, guichet unique, formation professionnelle technique, infrastructures énergétiques -- Nachtigal -- et corridors logistiques). La conclusion ouvrira sur le défi de la \cle{compétitivité} face aux importations asiatiques.
\end{exercice}



\newpage

\coursec{Corrigés des exercices}

\coursec{L'économie moderne au Cameroun}
\courssub{Corrigés des exercices et sujets type Bac}

\begin{corrige}[Exercice 1 --- QCM : Dualisme et poids de l'économie moderne]
Bonne réponse : \textbf{2}. Le \cle{dualisme} désigne la coexistence, au sein d'une même économie, de deux secteurs aux logiques opposées : un secteur moderne (capital intensif, salariat, marché mondial) et un secteur traditionnel (faible capital, autoconsommation, marché local). La proposition 1 est fausse car la coexistence n'est pas « harmonieuse » (concurrence pour la terre et la main-d'œuvre) ; la 3 est fausse (le secteur primaire traditionnel domine largement l'emploi, mais le secondaire moderne dépasse le primaire en PIB) ; la 4 est fausse car il existe des liens (contract farming, achat de récoltes, avances).
\end{corrige}

\begin{corrige}[Exercice 2 --- Vrai / Faux justifié : Statuts fonciers et acteurs]
\begin{enumerate}
    \item \textbf{Vrai.} Les ordonnances n\textsuperscript{o} 74-1, 74-2 et 74-3 de 1974 placent les terres non immatriculées dans le domaine national de l'État ; l'État y consent des baux (concessions provisoires puis définitives) aux agro-industries.
    \item \textbf{Faux.} L'\cle{outgrower} (planteur contractuel) est un exploitant indépendant lié à la firme par un contrat d'achat (fourniture d'intrants, encadrement technique, prix garanti) ; il n'est pas salarié et reste propriétaire-exploitant de sa parcelle.
    \item \textbf{Vrai.} Les baux de longue durée (bail emphytéotique ou concession définitive) peuvent atteindre 99 ans, ce qui sécurise les investissements lourds des plantations pérennes (palmier, hévéa, bananier).
    \item \textbf{Vrai.} Le sol du domaine national appartenant à l'État, seules les mises en valeur effectives (cultures, cases, arbres plantés) ouvrent droit à indemnisation en cas de reprise, jamais la terre nue.
\end{enumerate}
\end{corrige}

\begin{corrige}[Exercice 3 --- Définitions et calcul direct : Indicateurs industriels]
\begin{enumerate}
    \item Le \cle{taux de transformation locale} (TTL) est le rapport, exprimé en pourcentage, entre la quantité de matière première effectivement transformée sur le territoire national et la quantité totale produite (ou réceptionnée par les usines) :
    \[ \text{TTL} = \frac{\text{quantité transformée localement}}{\text{quantité totale produite}} \times 100 \]
    \item Application numérique :
    \[ \text{TTL} = \frac{62\,000}{85\,000} \times 100 \approx 72{,}9\,\% \]
    L'usine transforme localement environ 72,9\,\% des fèves réceptionnées.
    \item La certification \cle{RSPO} (\emph{Roundtable on Sustainable Palm Oil}, Table ronde sur l'huile de palme durable) garantit une production respectueuse de critères environnementaux et sociaux.
\end{enumerate}
\textbf{Points de vigilance :} ne pas oublier le $\times 100$ ; arrondir au dixième comme demandé ; vérifier que $62\,000/85\,000 \approx 0{,}7294$.
\end{corrige}

\begin{corrige}[Exercice 4 --- Questions courtes : Localisation et corridors industriels]
\begin{enumerate}
    \item ALUCAM à Edéa : (a) l'énergie hydroélectrique abondante et bon marché du barrage d'Edéa sur la Sanaga (l'électrolyse de l'alumine est très électro-intensive) ; (b) la proximité du port de Douala (40 km) pour l'importation de l'alumine (Guinée, Australie) et l'exportation de l'aluminium ; (c) la desserte par la voie ferrée Transcamerounais I (Douala--Yaoundé) et par la route nationale N3.
    \item L'axe (corridor) \textbf{Douala -- Edéa -- Yaoundé}, prolongé vers Kribi, concentre l'essentiel de l'industrie camerounaise ; ses pôles sont Douala (premier pôle industriel et portuaire), Edéa (énergie, aluminium, chimie) et Yaoundé (agroalimentaire, matériaux de construction, pharmacie).
    \item Le planteur villageois indépendant vend librement sa production (aux coopératives, aux négociants ou sur le marché local), sans contrat préalable ; l'outgrower est lié par contrat à une firme qui lui fournit intrants et encadrement et à laquelle il doit livrer sa récolte à un prix convenu.
\end{enumerate}
\end{corrige}

\begin{corrige}[Exercice 5 --- Analyse de document : Évolution de l'agro-industrie dans les exportations]
\begin{enumerate}
    \item Les deux plus fortes progressions en points de pourcentage sont le \textbf{cacao} ($+5{,}8$ points) et l'\textbf{huile de palme} ($+2{,}9$ points). Part cumulée en 2023 : $28{,}3 + 4{,}1 = 32{,}4\,\%$ des exportations non pétrolières.
    \item Causes possibles de la baisse des grumes : (a) la restriction de l'exportation des grumes brutes (loi de 1994, renforcée depuis) a réduit les volumes exportés par les exploitants incapables de transformer sur place ; (b) la transformation locale reste freinée par le coût de l'énergie et la vétusté des scieries, si bien que la production de bois ouvré ne compense pas la baisse des grumes.
    \item Mesures pour relancer la transformation : (a) développer les unités de transformation de proximité (usines de trituration du palmier à huile, chocolateries, usines de contreplaqué/panneaux) avec incitations fiscales et douanières ; (b) améliorer l'énergie et la logistique (corridors vers Douala et Kribi) et sécuriser le foncier des planteurs pour accroître les volumes transformables.
\end{enumerate}
\textbf{Points de vigilance :} bien distinguer variation en \emph{points de pourcentage} et taux de variation en \% ; les mesures proposées doivent être concrètes et reliées au PND30.
\end{corrige}

\begin{corrige}[Exercice 6 --- Étude de cas : Conflit foncier CDC / Communautés riveraines]
\begin{enumerate}
    \item (a) \textbf{L'État} (MINDCAF, administrations régionale et départementale) : propriétaire-gestionnaire du domaine national selon les ordonnances de 1974 ; (b) \textbf{la CDC} (Cameroon Development Corporation) : titulaire d'un bail de longue durée (concession), droit écrit d'occupation et d'exploitation ; (c) \textbf{les communautés riveraines} : détentrices de droits coutumiers antérieurs (terres de culture, jachères, forêts communautaires, sites sacrés).
    \item Pour la CDC, le bail consenti par l'État sur le domaine national est un titre juridique opposable : les terres non immatriculées appartiennent à l'État, qui peut les concéder. Pour les communautés, l'occupation ancestrale, les tombes, les cultures pérennes et les jachères (terres en repos mais appropriées) fondent un droit coutumier antérieur que la loi de 1974 a ignoré en classant ces terres dans le domaine national sans enquête préalable systématique.
    \item Issue négociée structurée en trois volets : (1) création d'une \textbf{commission de médiation} associant préfet, chefs traditionnels, élus locaux, CDC et représentants des jeunes et des femmes, avec cartographie participative des terres ; (2) \textbf{ajustement de l'assiette foncière} : restitution des jachères et sites sacrés, création de corridors villageois et de zones tampons, abandon de l'extension contestée ; (3) \textbf{partage de la valeur} : fonds de développement local abondé par hectare exploité, quota d'embauche locale prioritaire, contrats d'outgrowers pour les jeunes, redevances surfaciques reversées aux communes.
\end{enumerate}
\textbf{Points de vigilance :} citer explicitement les ordonnances foncières de 1974 ; présenter les deux logiques juridiques sans les opposer de manière manichéenne ; structurer la médiation en trois volets comme demandé.
\end{corrige}

\begin{corrige}[Exercice 7 --- Croquis de synthèse guidé : Pôles de l'économie moderne]
Barème indicatif de réussite : (1) titre explicite et orientation (flèche Nord) présents ; (2) au moins 4 plantations correctement localisées avec leur spéculation : CDC et PHP (bananier, palmier, hévéa) dans le Sud-Ouest/Littoral humide, SOCAPALM et SAFACAM (palmier, hévéa) dans la cuvette littorale (Moungo, Sanaga-Maritime), HEVECAM (hévéa) dans la région du Sud (Mvomeka'a, Bipindi) ; (3) les 4 unités industrielles placées : ALUCAM à Edéa (sur la Sanaga), SONARA à Limbé (côte), CIMENCAM à Douala (Bonabéri) et Figuil (Nord), SOSUCAM à Nkoteng (Centre) ; (4) les 3 barrages sur les bons fleuves : Edéa et Song Loulou (Sanaga), Lom Pangar (Lom) ; (5) les 2 ports (Douala, Kribi) et la zone franche de Kribi ; (6) le corridor Douala--Edéa--Yaoundé--Kribi en flèche épaisse ; (7) légende hiérarchisée par thèmes (plantations, industries, énergie, transports, ports). Le croquis doit faire apparaître la concentration littorale des pôles modernes et l'enclavement relatif du Nord et de l'Est.
\textbf{Points de vigilance :} un croquis sans titre, sans nord ou sans légende organisée est pénalisé ; chaque figuré doit être unique et constant entre carte et légende.
\end{corrige}

\begin{corrige}[Exercice 8 --- Dossier documentaire : Filière bois et transformation locale]
\begin{enumerate}
    \item Taux de croissance moyen annuel (taux actuariel) entre 2015 et 2023 (8 ans) :
    \[ t = \left(\frac{68}{42}\right)^{1/8} - 1 \approx 0{,}062 \quad \text{soit environ } 6{,}2\,\% \text{ par an.} \]
    En points de pourcentage, la progression est de $68 - 42 = 26$ points en 8 ans, soit environ $3{,}25$ points par an. À ce rythme, le seuil de 70\,\% est quasiment atteint (68\,\% en 2023) : le Cameroun est en bonne voie, mais le TTL déclaré reste à consolider et à vérifier sur le terrain (traçabilité, sciage artisanal).
    \item Facteurs \textbf{internes} : (b) vétusté des machines et difficultés de financement des équipements modernes. Facteurs \textbf{externes} : (a) coût et instabilité de l'énergie électrique, (c) concurrence du sciage artisanal informel non taxé, (d) enclavement et mauvais état des pistes forestières.
    \item Paragraphe attendu : la transformation locale du bois (sciage, déroulage, contreplaqué, meubles) crée des emplois industriels stables dans les zones de production (Est, Sud), augmente la valeur ajoutée par mètre cube exporté (un mètre cube de contreplaqué vaut plusieurs fois une grume brute) et diversifie les exportations non pétrolières. Elle constitue donc un levier majeur de l'industrialisation voulue par le PND30. Cependant, elle se heurte à des contraintes : gouvernance forestière perfectible (permis, traçabilité, concurrence du sciage artisanal informel non taxé), coût élevé de l'électricité et du crédit, enclavement des massifs forestiers qui renchérit le transport vers Douala et Kribi. Sans énergie compétitive, financement long et application effective de la loi de 1994, l'objectif de transformation reste fragile.
\end{enumerate}
\textbf{Points de vigilance :} bien distinguer taux actuariel (en \%/an) et progression en points ; le paragraphe doit mobiliser les trois documents et respecter la longueur demandée (10--12 lignes).
\end{corrige}

\begin{corrige}[Exercice 9 --- Sujet type Bac : Dissertation --- Aménagement du territoire]
\textbf{Plan attendu.} Introduction : définition de l'économie moderne (secteur capitaliste de plantations et d'industries hérité de la colonisation) et du \cle{dualisme} ; problématique : dans quelle mesure ce secteur structure-t-il l'espace camerounais ?
\textbf{I. L'économie moderne polarise le territoire} : concentration des plantations (CDC, PHP, SOCAPALM, SAFACAM, HEVECAM) dans le Littoral et le Sud-Ouest humides ; industries sur l'axe Douala--Edéa--Yaoundé (ALUCAM, CIMENCAM, SONARA à Limbé) ; rôle structurant des ports de Douala et Kribi et des barrages (Edéa, Song Loulou, Lom Pangar).
\textbf{II. Un aménagement inégalitaire et source de tensions} : \cle{enclavement} du Nord (seul pôle industriel : CIMENCAM à Figuil) et de l'Est (peu industrialisé, évacuation par le rail/route) ; \cle{conflits fonciers} entre baux des firmes et droits coutumiers ; dépendance aux exportations de matières premières brutes.
\textbf{III. Vers un aménagement rééquilibré (PND30)} : zones économiques spéciales, port en eau profonde et zone franche de Kribi, pipeline Tchad--Cameroun, projets hydroélectriques (Nachtigal, Grand Eweng), transformation locale des matières premières (bois, cacao, coton).
Conclusion : l'économie moderne aménage mais \cle{polarise} l'espace ; l'\cle{émergence} suppose de diffuser ses retombées par les \cle{infrastructures} et la transformation locale.
\textbf{Barème indicatif :} introduction problématisée (3 pts), développement équilibré avec exemples précis et localisés (10 pts), croquis légendé inséré (3 pts), conclusion avec ouverture (2 pts), langue et organisation (2 pts).
\textbf{Points de vigilance :} employer obligatoirement les six notions imposées (\cle{dualisme}, \cle{polarisation}, \cle{enclavement}, \cle{infrastructures}, \cle{conflits fonciers}, \cle{émergence}) ; chaque affirmation doit être étayée par un exemple localisé.
\end{corrige}

\begin{corrige}[Exercice 10 --- Sujet type Bac : Exercice cartographique commenté]
\textbf{Commentaire attendu (idées directrices).} Le croquis révèle une forte disparité spatiale : l'économie moderne se concentre dans le quart sud-ouest du pays, autour du corridor Douala--Edéa--Yaoundé--Kribi. Douala, port historique, draine l'essentiel des flux d'exportation (cacao, bois, coton, banane, aluminium) ; Kribi, avec son port en eau profonde, sa zone franche et le terminal du pipeline Tchad--Cameroun, émerge comme second pôle. Les plantations s'alignent sur le littoral humide (Sud-Ouest, Littoral, Sud), là où pluies abondantes et proximité portuaire se combinent. L'énergie hydroélectrique de la Sanaga (Edéa, Song Loulou) explique la localisation d'ALUCAM. En revanche, le Nord (hors Figuil et la plaine cotonnière) et l'Est forestier restent enclavés, évacuant leurs produits bruts par de longs corridors routiers et ferroviaires. Cette \cle{polarisation} traduit l'héritage colonial (économie de traite tournée vers la mer) et pose le problème du rééquilibrage visé par le PND30 : désenclaver l'intérieur, transformer sur place et diffuser la valeur ajoutée.
\textbf{Barème indicatif :} exactitude des localisations (8 pts), qualité graphique et légende organisée (4 pts), flux proportionnels et corridor (3 pts), commentaire structuré de 15 lignes (5 pts).
\textbf{Points de vigilance :} les flèches de flux doivent être proportionnelles à l'importance des exportations ; le commentaire doit analyser (causes et conséquences), pas seulement décrire la carte.
\end{corrige}

\begin{corrige}[Exercice 11 --- Sujet type Bac : Dissertation technique --- Industrialisation par la transformation]
\textbf{Plan attendu.} Introduction : le Cameroun exporte l'essentiel de ses matières premières à l'état brut ; la transformation locale est au cœur de la stratégie d'\cle{émergence} du PND30.
\textbf{I. Atouts} : ressources abondantes et diversifiées (cacao -- le Cameroun est 4\textsuperscript{e} producteur mondial --, bois du bassin du Congo, coton du Nord, palmier à huile et hévéa du littoral) ; main-d'œuvre jeune et abondante ; marché régional CEMAC et débouchés de la ZLECAf ; cadre stratégique du PND30.
\textbf{II. Contraintes} : déficit et coût de l'énergie électrique ; accès difficile au crédit long terme ; logistique coûteuse (enclavement, saturation du port de Douala) ; TTL encore insuffisant (bois : 68\,\% en 2023 ; cacao : majoritairement exporté en fèves) ; poids du secteur informel et gouvernance perfectible.
\textbf{III. Leviers} : zones économiques spéciales (Kribi) et zone franche ; incitations fiscales et douanières ; guichet unique de l'investissement ; formation technique et professionnelle ; barrage de Nachtigal et corridors routiers, ferroviaires et portuaires.
Conclusion : la transformation locale est la voie de l'émergence, à condition de gagner la bataille de la \cle{compétitivité} face aux produits manufacturés asiatiques.
\textbf{Barème indicatif :} introduction avec définitions et problématique (3 pts), trois parties équilibrées et illustrées (11 pts), conclusion ouverte sur la compétitivité (3 pts), maîtrise de la langue et rigueur du plan (3 pts).
\textbf{Points de vigilance :} chaque partie doit citer au moins deux filières concrètes (cacao, bois, coton, huile de palme, caoutchouc) ; les chiffres (TTL du bois, objectif de 70\,\%) doivent être exacts ; la conclusion doit obligatoirement ouvrir sur la \cle{compétitivité}.
\end{corrige}

\newpage

\coursec{Fiche bilan}

\begin{popfigure}
\begin{center}\tcbox[colback=popblue,colframe=popblue,coltext=white,arc=9pt,boxsep=4pt,left=6pt,right=6pt]{\bfseries\small L'économie moderne au Cameroun}\end{center}
\vspace{-2pt}
\begin{tcbraster}[raster columns=3,raster equal height=rows,raster column skip=6pt,raster row skip=6pt,enhanced,blankest]
\begin{tcolorbox}[enhanced,colback=poporange,colframe=poporange,coltext=white,boxrule=0.9pt,arc=3pt,left=4pt,right=4pt,top=3pt,bottom=3pt,boxsep=1pt,halign=left]\bfseries\scriptsize Dualisme : moderne / traditionnel\end{tcolorbox}
\begin{tcolorbox}[enhanced,colback=popgreen,colframe=popgreen,coltext=white,boxrule=0.9pt,arc=3pt,left=4pt,right=4pt,top=3pt,bottom=3pt,boxsep=1pt,halign=left]\bfseries\scriptsize Grandes plantations : CDC, SOCAPALM, HEVECAM\end{tcolorbox}
\begin{tcolorbox}[enhanced,colback=poppurple,colframe=poppurple,coltext=white,boxrule=0.9pt,arc=3pt,left=4pt,right=4pt,top=3pt,bottom=3pt,boxsep=1pt,halign=left]\bfseries\scriptsize Foncier : concessions, cahier des charges\end{tcolorbox}
\begin{tcolorbox}[enhanced,colback=popgold,colframe=popgold,coltext=white,boxrule=0.9pt,arc=3pt,left=4pt,right=4pt,top=3pt,bottom=3pt,boxsep=1pt,halign=left]\bfseries\scriptsize Industrie : typologie et localisation\end{tcolorbox}
\begin{tcolorbox}[enhanced,colback=poppink,colframe=poppink,coltext=white,boxrule=0.9pt,arc=3pt,left=4pt,right=4pt,top=3pt,bottom=3pt,boxsep=1pt,halign=left]\bfseries\scriptsize Pôles : Douala, Édéa, Kribi\end{tcolorbox}
\begin{tcolorbox}[enhanced,colback=popblueL,colframe=popblueL,coltext=white,boxrule=0.9pt,arc=3pt,left=4pt,right=4pt,top=3pt,bottom=3pt,boxsep=1pt,halign=left]\bfseries\scriptsize Enjeux : emploi, environnement\end{tcolorbox}
\begin{tcolorbox}[enhanced,colback=popgreenL,colframe=popgreenL,coltext=white,boxrule=0.9pt,arc=3pt,left=4pt,right=4pt,top=3pt,bottom=3pt,boxsep=1pt,halign=left]\bfseries\scriptsize Trajectoire d'émergence 2035\end{tcolorbox}
\end{tcbraster}

\legende{Carte mentale de synthèse : les sept idées-forces de la leçon « L'économie moderne au Cameroun ». À partir du centre, chaque branche correspond à une partie du plan : le cadre général (dualisme), les grandes plantations, le foncier, l'industrie, les pôles, les enjeux et la perspective d'émergence.}
\end{popfigure}

\begin{aretenir}[L'essentiel de la leçon]
\textbf{Définitions clés}
\begin{itemize}
\item \cle{Économie moderne} : secteur monétisé, organisé en entreprises, intégré aux échanges nationaux et internationaux (grandes plantations, industries, banques, transports modernes).
\item \cle{Dualisme économique} : coexistence au Cameroun d'un secteur moderne minoritaire en emplois mais dominant en valeur, et d'un secteur traditionnel (agriculture vivrière, artisanat, informel) majoritaire en population.
\item \cle{Agriculture de plantation} : grande exploitation capitaliste, spécialisée dans une culture d'exportation (banane, palmier à huile, hévéa, cacao), associant champs, usine de transformation et main-d'œuvre salariée.
\item \cle{Agro-industrie} : intégration de la production agricole et de sa transformation industrielle au sein d'une même firme.
\item \cle{Concession foncière} : mise à disposition de terres du domaine national à une entreprise par l'État, contre un cahier des charges (emplois, équipements sociaux, redevances).
\item \cle{Facteur de localisation industrielle} : élément qui explique l'implantation d'une usine (matières premières, énergie, main-d'œuvre, marché, transports, capitaux, décisions de l'État).
\end{itemize}

\textbf{Repères chiffrés et acteurs à connaître par cœur}
\begin{center}
\begin{tabularx}{\linewidth}{|l|X|X|}\hline
\rowcolor{popblueL} \textbf{Firme} & \textbf{Production} & \textbf{Localisation principale} \\ \hline
CDC & banane, palmier à huile, hévéa & région du Sud-Ouest (Tiko, Mbonge) \\ \hline
SOCAPALM & huile de palme & Littoral (Dibombari, Édéa), Sud \\ \hline
HEVECAM & caoutchouc (hévéa) & Sud (Nieté, près de Kribi) \\ \hline
ALUCAM & aluminium (alumine importée) & Édéa (énergie hydroélectrique) \\ \hline
CIMENCAM & ciment & Douala (Bonabéri), Figuil \\ \hline
SOSUCAM & sucre & Nkoteng et Bandjock (Centre) \\ \hline
\end{tabularx}
\end{center}

\begin{itemize}
\item Douala : premier pôle industriel et portuaire (environ 60 \% de la production industrielle) ; Édéa : métallurgie et énergie ; Kribi : port en eau profonde et zone industrielle en croissance.
\item Le secteur secondaire représente environ un quart du PIB ; l'industrie manufacturière reste dominée par l'agroalimentaire.
\end{itemize}

\textbf{Méthodes express}
\begin{itemize}
\item \emph{Croquis} : contour simplifié du pays, titre, légende numérotée, flèche du nord ; figer les pôles par des points ou zones colorées, les flux par des flèches proportionnelles.
\item \emph{Étude de cas} : présenter la firme (statut, production), décrire l'organisation spatiale (plantation, usine, cités ouvrières, voies d'évacuation), puis discuter les enjeux (emploi, foncier, environnement).
\item \emph{Justifier une localisation} : citer au moins trois facteurs et les relier au site concret.
\end{itemize}
\end{aretenir}

\begin{aretenir}[Checklist avant le Bac]
\begin{itemize}
\item[$\square$] Je sais définir l'économie moderne et le dualisme économique camerounais.
\item[$\square$] Je connais les trois grandes firmes agro-industrielles (CDC, SOCAPALM, HEVECAM), leurs productions et leurs localisations.
\item[$\square$] Je sais expliquer la genèse coloniale des grandes plantations.
\item[$\square$] Je maîtrise les statuts fonciers (domaine national, concession, cahier des charges).
\item[$\square$] Je sais décrire l'organisation spatiale d'une plantation agro-industrielle.
\item[$\square$] Je connais la typologie des industries camerounaises (extractive, agroalimentaire, chimique, métallurgique, matériaux).
\item[$\square$] Je sais citer et appliquer les facteurs de localisation industrielle à un exemple précis.
\item[$\square$] Je situe les pôles industriels majeurs : Douala, Édéa, Kribi, Yaoundé.
\item[$\square$] Je sais réaliser un croquis « Les pôles de l'économie moderne au Cameroun » avec titre, légende et nord.
\item[$\square$] Je connais les enjeux : emploi, conflits fonciers, déforestation, pollution.
\item[$\square$] Je sais relier l'économie moderne à la vision d'émergence du Cameroun à l'horizon 2035.
\item[$\square$] Je rédige une conclusion justifiée en trois temps (rappel, bilan nuancé, ouverture).
\end{itemize}
\end{aretenir}

\findecours

\end{document}
